\documentclass[11pt,reqno]{amsart}

\usepackage{amsfonts}
\usepackage[utf8]{inputenc}
\usepackage{amsmath,amssymb,amsthm}
\usepackage{changepage}
\usepackage{ragged2e}
\usepackage{float}
\usepackage{array}
\usepackage{graphicx}
\usepackage{setspace}
\usepackage{xcolor}
\usepackage{enumitem}
\usepackage[left=0.8in,right=0.8in,bottom=1in,top=1in,headheight=15pt,headsep=40pt]{geometry}
\usepackage{tikz}
\usepackage{fancyvrb}
\usepackage{fancyhdr}
\usepackage{titlesec,titletoc}
\usepackage[colorlinks=true,linkcolor=blue,citecolor=green,urlcolor=magenta]{hyperref}
\usepackage{aliascnt}
\usepackage[nameinlink]{cleveref}
\usepackage{listings}
\usepackage{booktabs}
\providecommand{\TEXTsymbol}[1]{#1}
\setcounter{MaxMatrixCols}{10}
\usepackage{longtable}
\newcommand{\Prime}{\mathsf{Prime}}
\newcommand{\floor}{\mathsf{floor}}
\newcommand{\Ingham}{\mathsf{Ingham}}
\newcommand{\MillsRec}{\mathsf{MillsRec}}
\newcommand{\IsMills}{\mathsf{IsMills}}
\newcommand{\N}{\mathbb{N}}
\newcommand{\Z}{\mathbb{Z}}
\newcommand{\R}{\mathbb{R}}
\restylefloat{figure}
\hypersetup{pageanchor=false}
\makeatletter
\@ifundefined{theorem}{}{\let\theorem\relax}
\@ifundefined{definition}{}{\let\definition\relax}
\@ifundefined{lemma}{}{\let\lemma\relax}
\@ifundefined{corollary}{}{\let\corollary\relax}
\@ifundefined{proposition}{}{\let\proposition\relax}
\@ifundefined{remark}{}{\let\remark\relax}
\@ifundefined{statement}{}{\let\statement\relax}
\makeatother
\DeclareUnicodeCharacter{03B9}{\ensuremath{\iota}}
\DeclareUnicodeCharacter{2200}{\ensuremath{\forall}}
\DeclareUnicodeCharacter{2203}{\ensuremath{\exists }}
\DeclareUnicodeCharacter{2208}{\ensuremath{\in}}
\DeclareUnicodeCharacter{2218}{\ensuremath{\circ}}
\DeclareUnicodeCharacter{2264}{\ensuremath{\leq}}
\DeclareUnicodeCharacter{2265}{\ensuremath{\geq}}
\DeclareUnicodeCharacter{2192}{\ensuremath{\to}}
\DeclareUnicodeCharacter{21A6}{\ensuremath{\mapsto}}
\DeclareUnicodeCharacter{2260}{\ensuremath{\neq}}
\DeclareUnicodeCharacter{00D7}{\ensuremath{\times}}
\DeclareUnicodeCharacter{2115}{\ensuremath{\mathbb{N}}}
\DeclareUnicodeCharacter{211D}{\ensuremath{\mathbb{R}}}
\DeclareUnicodeCharacter{27E8}{\ensuremath{\langle}}
\DeclareUnicodeCharacter{27E9}{\ensuremath{\rangle}}
\providecommand{\R}{\mathbb{R}}
\providecommand{\N}{\mathbb{N}}
\providecommand{\Q}{\mathbb{Q}}
\providecommand{\Z}{\mathbb{Z}}
\providecommand{\period}{\,\allowbreak}
\newcommand{\pagenumside}{C}
\newlength{\pageNumberOffset}
\setlength{\pageNumberOffset}{0pt}
\newcommand{\PagenumSideC}{C}
\newcommand{\applypagehead}[1]{\expandafter\fancyhead\expandafter[\pagenumside]{#1}\ifx\pagenumside\PagenumSideC\else\ifdim\pageNumberOffset=0pt\else\expandafter\fancyheadoffset\expandafter[\pagenumside]{\pageNumberOffset}\fi\fi}
\newcommand{\applypagefoot}[1]{\expandafter\fancyfoot\expandafter[\pagenumside]{#1}\ifx\pagenumside\PagenumSideC\else\ifdim\pageNumberOffset=0pt\else\expandafter\fancyfootoffset\expandafter[\pagenumside]{\pageNumberOffset}\fi\fi}
\pagestyle{fancy}
\fancyhf{}
\applypagehead{\thepage}
\renewcommand{\headrulewidth}{0pt}
\renewcommand{\footrulewidth}{0pt}
\fancypagestyle{plain}{\fancyhf{}\applypagehead{\thepage}\renewcommand{\headrulewidth}{0pt}}
\setlength{\arraycolsep}{1.5pt}
\setlength{\tabcolsep}{1.5pt}
\renewcommand{\arraystretch}{1.0}
\setlength{\arrayrulewidth}{0.4pt}
\setlength{\doublerulesep}{1pt}
\setlength{\extrarowheight}{0pt}
\newcolumntype{c}{>{\centering\arraybackslash\scriptsize}m{40pt}}
\newcolumntype{r}{>{\centering\arraybackslash\scriptsize}m{2.5cm}}
\newcolumntype{l}{>{\centering\arraybackslash\scriptsize}m{10.5cm}}
\fvset{frame=lines,framerule=0.7pt,rulecolor=\color{blue!60!black},framesep=2pt,fontsize=\scriptsize,xleftmargin=4pt,xrightmargin=4pt}
\RecustomVerbatimEnvironment{verbatim}{Verbatim}{}
\DefineVerbatimEnvironment{verbatimbox}{Verbatim}{frame=single,framerule=0.7pt,rulecolor=\color{blue!60!black},framesep=2pt,fontsize=\scriptsize,xleftmargin=4pt,xrightmargin=4pt}
\makeatletter
\def\@setsubjclass{\ifx\@subjclass\@empty\else\@subjclass\fi}
\makeatother
\subjclass{\textbf{2020 Mathematics Subject Classification.}\,Primary 11A41, 11N05; Secondary 03B35, 03B38, 03F35, 03F60, 68V15, 68V20, 11-03.}
\keywords{prime-representing function, Mills' theorem, prime gaps, Ingham's estimate, Bertrand's postulate, finitary predicates, typed predicates, reverse mathematics, Lean formalisation, constructive mathematics}
\renewcommand{\abstractname}{\fontsize{9}{11}\selectfont Abstract}
\newcommand{\abstractfont}{\normalfont\fontsize{8}{10}\selectfont}
\newcommand{\abstractheadingsep}{\vspace{0.5em}}
\makeatletter
\renewenvironment{abstract}{\ifx\maketitle\relax\ClassWarning{\@classname}{Abstract should precede\protect\maketitle\space in AMS document classes reported}\fi\global\setbox\abstractbox=\vtop \bgroup\abstractfont{\centering\abstractname\par}\abstractheadingsep\list{}{\labelwidth\z@\leftmargin3pc \rightmargin\leftmargin\listparindent\normalparindent \itemindent\z@\parsep\z@ \@plus\p@\let\fullwidthdisplay\relax}\item\relax}{\endlist\egroup\ifx\@setabstract\relax \@setabstracta \fi}
\makeatother
\setcounter{secnumdepth}{3}
\newcommand{\SecAbove}{1.0ex plus 0.3ex minus 0.1ex}
\newcommand{\SecBelow}{0pt}
\newcommand{\SubsecAbove}{0.8ex plus 0.2ex minus 0.1ex}
\newcommand{\SubsecBelow}{0pt}
\newcommand{\SubsubsecAbove}{0.6ex plus 0.2ex minus 0.1ex}
\newcommand{\SubsubsecBelow}{0pt}
\makeatletter
\titleformat{\section}{\normalfont\fontsize{16}{18}\selectfont\bfseries}{\thesection.}{0.5em}{}
\titleformat{\subsection}{\normalfont\fontsize{14}{16}\selectfont\bfseries}{\thesubsection.}{0.5em}{}
\titleformat{\subsubsection}{\normalfont\fontsize{12}{14}\selectfont\bfseries}{\thesubsubsection.}{0.5em}{}
\titlespacing*{\section}{\z@}{\SecAbove}{\SecBelow}
\titlespacing*{\subsection}{\z@}{\SubsecAbove}{\SubsecBelow}
\titlespacing*{\subsubsection}{\z@}{\SubsubsecAbove}{\SubsubsecBelow}
\makeatother
\newcounter{axctr}\newcounter{posctr}\newcounter{defctr}
\newcounter{lemctr}\newcounter{thmctr}\newcounter{asserctr}\newcounter{corctr}
\renewcommand{\theaxctr}{\arabic{axctr}}
\renewcommand{\theposctr}{\arabic{posctr}}
\renewcommand{\thedefctr}{\arabic{defctr}}
\renewcommand{\thelemctr}{\arabic{lemctr}}
\renewcommand{\thethmctr}{\arabic{thmctr}}
\renewcommand{\theasserctr}{\arabic{asserctr}}
\renewcommand{\thecorctr}{\arabic{corctr}}
\renewcommand{\theHaxctr}{\thesection.\arabic{axctr}}
\renewcommand{\theHposctr}{\thesection.\arabic{posctr}}
\renewcommand{\theHdefctr}{\thesection.\arabic{defctr}}
\renewcommand{\theHlemctr}{\thesection.\arabic{lemctr}}
\renewcommand{\theHthmctr}{\thesection.\arabic{thmctr}}
\renewcommand{\theHasserctr}{\thesection.\arabic{asserctr}}
\providecommand{\theHasser}{\theequation}
\renewcommand{\theHcorctr}{\thesection.\arabic{corctr}}
\newcommand{\swpax}[2]{\refstepcounter{axctr}\hypertarget{#1}{\mbox{\textbf{Ax\,\theaxctr}}}\label{#2}}
\newcommand{\swppos}[2]{\refstepcounter{posctr}\hypertarget{#1}{\mbox{\textbf{Pos\,\theposctr}}}\label{#2}}
\newcommand{\swpdef}[2]{\refstepcounter{defctr}\hypertarget{#1}{\mbox{\textbf{Def\,\thedefctr}}}\label{#2}}
\newcommand{\swplem}[2]{\refstepcounter{lemctr}\hyperlink{#1}{\mbox{\textbf{Lem\,\thelemctr}}}\label{#2}}
\newcommand{\swpthm}[2]{\refstepcounter{thmctr}\hyperlink{#1}{\mbox{\textbf{Thm\,\thethmctr}}}\label{#2}}
\newcommand{\swpcor}[2]{\refstepcounter{corctr}\hypertarget{#1}{\mbox{\textbf{Cor\,\thecorctr}}}\label{#2}}
\newcommand{\tagged}[2]{\hypertarget{#1}{\mbox{\textbf{#2}}}}
\newcommand{\autorefto}[2]{\hyperlink{#1}{#2}}
\newcommand{\mcite}[1]{\text{\cite{#1}}}
\newtheoremstyle{blueplain}{2pt plus 1pt minus 0.5pt}{1pt plus 0.5pt minus 0.3pt}{\itshape}{}{\bfseries\color{blue}}{.}{\newline}{\thmname{#1}\thmnumber{ #2}\thmnote{ {\bfseries\color{blue}(#3)}}}
\theoremstyle{blueplain}
\newtheorem{theorem}{Theorem}[section]
\newtheorem{lemma}{Lemma}[section]
\newtheorem{corollary}{Corollary}[section]
\newtheorem{proposition}{Proposition}[section]
\newtheorem{remark}{Remark}[section]
\newtheorem{definition}{Definition}[section]
\newtheorem{statement}{Statement}[section]
\definecolor{proofhead}{RGB}{11,83,148}
\definecolor{proofqed}{RGB}{204,0,0}
\renewcommand{\qedsymbol}{\textcolor{proofqed}{\ensuremath{\square}}}
\renewenvironment{proof}[1][Proof]{\par\noindent\textbf{\textcolor{proofhead}{#1.}}\par\nobreak}{\par\nobreak\hfill\qedsymbol}
\newenvironment{swpproof}[1][Proof]{\begin{proof}[#1]}{\end{proof}}
\renewcommand{\theHtheorem}{\thesection.\arabic{theorem}}
\renewcommand{\theHstatement}{\thesection.\arabic{statement}}
\providecommand{\axctrautorefname}{Axiom}
\providecommand{\posctrautorefname}{Postulate}
\providecommand{\defctrautorefname}{Definition}
\providecommand{\lemctrautorefname}{Lemma}
\providecommand{\thmctrautorefname}{Theorem}
\providecommand{\asserctrautorefname}{Proposition}
\providecommand{\corctrautorefname}{Corollary}
\providecommand{\statementautorefname}{Statement}
\crefname{axctr}{axiom}{axioms}
\crefname{posctr}{postulate}{postulates}
\crefname{defctr}{definition}{definitions}
\crefname{lemctr}{lemma}{lemmas}
\crefname{thmctr}{theorem}{theorems}
\crefname{asserctr}{proposition}{propositions}
\crefname{corctr}{corollary}{corollaries}
\Crefname{axctr}{Axiom}{Axioms}
\Crefname{posctr}{Postulate}{Postulates}
\Crefname{defctr}{Definition}{Definitions}
\Crefname{lemctr}{Lemma}{Lemmas}
\Crefname{thmctr}{Theorem}{Theorems}
\Crefname{asserctr}{Proposition}{Propositions}
\Crefname{corctr}{Corollary}{Corollaries}
\crefname{statement}{statement}{statements}
\Crefname{statement}{Statement}{Statements}
\crefname{section}{section}{sections}
\Crefname{section}{Section}{Sections}
\crefname{equation}{eq}{eqs}
\Crefname{equation}{Eq}{Eqs}
\crefformat{equation}{#2\textcolor{blue}{eq~#1}#3}
\Crefformat{equation}{#2\textcolor{blue}{Eq~#1}#3}
\newcounter{inq}
\renewcommand{\theinq}{\theequation}
\makeatletter
\renewcommand{\p@inq}{\p@equation}
\makeatother
\crefname{inq}{ineq}{inequalities}
\Crefname{inq}{Ineq}{Inequalities}
\crefformat{inq}{#2\textcolor{blue}{inequality~#1}#3}
\Crefformat{inq}{#2\textcolor{blue}{Inequality~#1}#3}
\crefmultiformat{inq}{#2\textcolor{blue}{inequalities~#1}#3}{\textcolor{blue}{ ,\,~#2#1#3}}{\textcolor{blue}{, #2#1#3}}{\textcolor{blue}{ ,\,~#2#1#3}}
\Crefmultiformat{inq}{#2\textcolor{blue}{Inequalities~#1}#3}{\textcolor{blue}{ ,\,~#2#1#3}}{\textcolor{blue}{, #2#1#3}}{\textcolor{blue}{ ,\,~#2#1#3}}
\makeatletter
\newenvironment{ineq}
{\begin{equation}\refstepcounter{inq}}
{\end{equation}}
\makeatother
\newcounter{type}
\renewcommand{\thetype}{\theequation}
\makeatletter
\renewcommand{\p@type}{\p@equation}
\makeatother
\crefname{type}{type}{types}
\Crefname{type}{Type}{Types}
\crefformat{type}{#2\textcolor{blue}{type~#1}#3}
\Crefformat{type}{#2\textcolor{blue}{Type~#1}#3}
\crefmultiformat{type}{#2\textcolor{blue}{types~#1}#3}{\textcolor{blue}{ ,\,~#2#1#3}}{\textcolor{blue}{, #2#1#3}}{\textcolor{blue}{ ,\,~#2#1#3}}
\Crefmultiformat{type}{#2\textcolor{blue}{Types~#1}#3}{\textcolor{blue}{ ,\,~#2#1#3}}{\textcolor{blue}{, #2#1#3}}{\textcolor{blue}{ ,\,~#2#1#3}}
\makeatletter
\newenvironment{type}
{\begin{equation}\refstepcounter{type}}
{\end{equation}}
\makeatother
\newcounter{asser}
\renewcommand{\theasser}{\theequation}
\makeatletter
\renewcommand{\p@asser}{\p@equation}
\makeatother
\crefname{asser}{Assertion}{assertions}
\Crefname{asser}{Assertion}{Assertions}
\crefformat{asser}{#2\textcolor{blue}{assertion~#1}#3}
\Crefformat{asser}{#2\textcolor{blue}{Assertion~#1}#3}
\crefmultiformat{asser}{#2\textcolor{blue}{assertions~#1}#3}{\textcolor{blue}{ ,\,~#2#1#3}}{\textcolor{blue}{, #2#1#3}}{\textcolor{blue}{ ,\,~#2#1#3}}
\Crefmultiformat{asser}{#2\textcolor{blue}{Assertions~#1}#3}{\textcolor{blue}{ ,\,~#2#1#3}}{\textcolor{blue}{, #2#1#3}}{\textcolor{blue}{ ,\,~#2#1#3}}
\makeatletter
\newenvironment{asser}
{\begin{equation}\refstepcounter{asser}}
{\end{equation}}
\makeatother
\newcommand{\thmref}[1]{\hyperref[#1]{Theorem~\ref*{#1}}}
\newcommand{\lemref}[1]{\hyperref[#1]{Lemma~\ref*{#1}}}
\newcommand{\corref}[1]{\hyperref[#1]{Corollary~\ref*{#1}}}
\newcommand{\propref}[1]{\hyperref[#1]{Proposition~\ref*{#1}}}
\newcommand{\defref}[1]{\hyperref[#1]{Definition~\ref*{#1}}}
\makeatletter
\let\orig@maketitle\maketitle
\renewcommand{\maketitle}{\begingroup\let\MakeUppercase\@firstofone\let\uppercase\@firstofone\orig@maketitle\endgroup}
\newlength{\titleToAuthorSkip}
\newlength{\authorToAbstractSkip}
\setlength{\titleToAuthorSkip}{6pt plus 2pt minus 1pt}
\setlength{\authorToAbstractSkip}{10pt plus 3pt minus 2pt}
\renewcommand{\@maketitle}{\normalfont\normalsize\let\@makefnmark\relax\let\@thefnmark\relax\ifx\@empty\@subjclass\else\@footnotetext{\@setsubjclass}\fi\ifx\@empty\@keywords\else\@footnotetext{\@setkeywords}\fi\ifx\@empty\@thanks\else\@footnotetext{\@setthanks}\fi\@settitle\ifx\@empty\@authors\else\vspace{\titleToAuthorSkip}\@setauthors\fi\ifx\@empty\@date\else\vspace{3pt}\@setdate\fi\ifx\@empty\@abstract\else\vspace{\authorToAbstractSkip}\fi\@setabstract}
\makeatother
\makeatletter
\renewcommand{\encodingdefault}{OT1}
\makeatother
\renewcommand{\rmdefault}{cmr}
\renewcommand{\sfdefault}{cmss}
\renewcommand{\ttdefault}{cmtt}
\fontencoding{OT1}\fontfamily{cmr}\selectfont
\renewcommand{\contentsname}{\texorpdfstring{{\color{blue!60!black}\sffamily\Large\selectfont\bfseries Contents}}{Contents}}
\renewcommand{\seriesdefault}{\bfdefault}
\boldmath
\providecommand{\comma}{,\,}
\IfFileExists{tcilatex.tex}{\input{tcilatex}}{}
\makeatletter
\def\dudley@listtighten{  \setlength{\topsep}{0pt}  \setlength{\partopsep}{0pt}  \setlength{\itemsep}{1pt plus 0.5pt minus 0.3pt}  \setlength{\parsep}{0pt}  \setlength{\labelsep}{0.4em}}
\def\dudley@itemlabel{  \ifnum\@listdepth=\@ne \labelitemi
	\else\ifnum\@listdepth=\tw@ \labelitemii
	\else\ifnum\@listdepth=\thr@@ \labelitemiii
	\else \labelitemiv\fi\fi\fi}
\def\dudley@enumlabel{  \ifnum\@listdepth=\@ne \arabic{enumi}.
	\else\ifnum\@listdepth=\tw@ \alph{enumii}.
	\else\ifnum\@listdepth=\thr@@ \roman{enumiii}.
	\else \Alph{enumiv}.\fi\fi\fi}
\def\dudley@enumcounter{  \ifnum\@listdepth=\@ne \usecounter{enumi}  \else\ifnum\@listdepth=\tw@ \usecounter{enumii}  \else\ifnum\@listdepth=\thr@@ \usecounter{enumiii}  \else \usecounter{enumiv}\fi\fi\fi}

\makeatother
\setlength{\parindent}{0pt}
\setlength{\parskip}{0.25em plus 0.05em minus 0.05em}
\makeatletter
\renewenvironment{quotation}
{\list{}{\listparindent\z@
		\itemindent\z@
		\rightmargin\leftmargin
		\parsep\z@}	\item\relax}
{\endlist}
\makeatother
\numberwithin{equation}{section}
\newcolumntype{p}[1]{>{\centering\arraybackslash\scriptsize}m{#1}|}
\lstdefinestyle{python}{
	language=Python,
	basicstyle=\ttfamily\scriptsize,
	keywordstyle=\color{blue},
	commentstyle=\color{gray},
	stringstyle=\color{orange},
	numbers=left,
	numberstyle=\tiny,
	breaklines=true,
	frame=single,
	showstringspaces=false,
	xleftmargin=4pt,
	xrightmargin=4pt
}
\lstdefinestyle{lean}{
	language=ML,
	basicstyle=\ttfamily\scriptsize,
	keywordstyle=\color{blue},
	commentstyle=\color{gray},
	stringstyle=\color{orange},
	numbers=left,
	numberstyle=\tiny,
	breaklines=true,
	frame=single,
	showstringspaces=false,
	xleftmargin=4pt,
	xrightmargin=4pt
}
\setcounter{tocdepth}{1}
\contentsmargin{1em}
\titlecontents{part}[0pt]{\addvspace{14pt}\large\bfseries}{\makebox[2.0em][l]{\thecontentslabel.}}{}{\titlerule*[0.6pc]{.}\contentspage}[\addvspace{4pt}]
\titlecontents{section}[2.0em]{\addvspace{10pt}\small\bfseries}{\makebox[2.0em][l]{\thecontentslabel.}}{\hspace*{-2.0em}}{\titlerule*[0.6pc]{.}\contentspage}[\addvspace{2pt}]
\titlecontents{subsection}[4.5em]{\addvspace{4pt}\footnotesize}{\makebox[2.5em][l]{\thecontentslabel.}}{\hspace*{-2.5em}}{\titlerule*[0.6pc]{.}\contentspage}[\addvspace{1pt}]
\setlist[itemize]{topsep=0pt,partopsep=0pt,itemsep=1pt plus 0.5pt minus 0.3pt,parsep=0pt,leftmargin=1.5em,labelsep=0.4em,label=\textbullet}
\setlist[enumerate]{topsep=0pt,partopsep=0pt,itemsep=1pt plus 0.5pt minus 0.3pt,parsep=0pt,leftmargin=1.8em,labelsep=0.4em}

\begin{document}
\title{\bfseries{\huge \emph{Mills' Theorem in Finitary Form.}}}
\author{ 
	\bfseries{\large \hspace{1.2in} Brahimi Mahdi Tahar\newline\text{ \
		} \newline\text{ \  } \hspace{0.5in}\text{ \  } University of Mohamed
		Boudiaf of M'sila, Algeria\newline\text{ \  }
		\newline\href{mailto:mahditahar.brahimi@univ-msila.dz}{mahditahar.brahimi@univ-msila.dz}}
}
\maketitle
\begin{abstract}
	\noindent \setstretch{1.2}
	
	\begin{itemize}
		\item We formalise, as typed predicates over $\mathbb{N}$, $\mathbb{Z}$,
		$\mathbb{R}$, and $\mathbb{C}$, the ten prime-representing assertions of
		Mills \cite{Mills1947}, Kuipers \cite{Kuipers1950}, Ansari
		\cite{Ansari1951}, Wright \cite{Wright1951,Wright1954}, Ore \cite{Ore1952},
		and Niven \cite{Niven1951}, together with their Lean counterparts
		\cite{FormalConjectures,MathlibGraph}.
		
	
		
		\item The central statement is \cref{asser:mills}: from the finitary
		Ingham estimate \cref{ineq:ingham} one derives
		\begin{equation*}
		\exists \theta :\mathbb{R}_{>1},\,\mathrm{IsMills}(\theta ),
		\qquad
		\mathrm{IsMills}(\theta ) :\iff
		\theta >1\wedge\forall n\in\mathbb{N}_{\geq 1},\
		\mathsf{Prime}\bigl(\mathsf{floor}(\theta ^{3^{n}})\bigr).
		\end{equation*}
		
		\item All definitions and theorems carry explicit type annotations
		$\mathsf{Object}:\mathsf{Type}$, and each block is paired with a Lean
		declaration of the same signature; the reverse-mathematical classification
		(\cref{sec:thm-summary}) is stated over the finitary systems
		$\mathrm{RCA}_0$, $\mathrm{WKL}_0$, $\mathrm{ACA}_0$ via the finite-level
		comprehension schemes $\Delta^0_1\text{-}\mathsf{CA}$, $\mathsf{WKL}$,
		$\mathsf{ACA}$.
		
		\item The subtypes $\mathbb{R}_{>1}$, $\mathbb{R}_{>0}$, and
		$\mathbb{N}_{\geq 1}$ replace set-theoretic comprehension throughout, and
		$\mathsf{LEM}$ is recorded as a hypothesis in \cref{sec:axioms} rather than
		built into the definition of $\mathsf{Prop}$, so that the constructive
		content of each route remains explicit.
	\end{itemize}
\end{abstract}

\newpage \thispagestyle{empty}

\bigskip

\tableofcontents

\bigskip

\newpage \setstretch{1.2} \pagestyle{plain} \pagenumbering{arabic} %
\hypersetup{pageanchor=true}

\section{Introduction}

The prime-representing function of W.~H.~Mills~\cite{Mills1947} states,
formally, \cref{asser:mills}.

The formula produces prime values exclusively, yet is computationally inert:
\begin{center}
Any construction of $\theta $ presupposes the ability to recognise
arbitrarily large primes.
\end{center}

We record ten logically independent routes to this conclusion, organised
through this document as follows:

\begin{itemize}
\item The definitions in \cref{sec:defs}
\item The axioms in \cref{sec:axioms}
\item The postulates in \cref{sec:postulates}
\item The theorem statements in \cref{sec:theorems}: prime-gap constructions
in \cref{sec:thm-prime-gap}, namely
\cref{asser:mills,asser:kuipers,asser:ansari}; elementary and abstract
constructions in \cref{sec:thm-elementary}, namely
\cref{asser:wright,asser:ore,asser:direct}; variants and metrical theorems in
\cref{sec:thm-variants}, namely \cref{asser:niven,asser:wright-metrical}; and
the Lean formulation in \cref{sec:thm-lean}, namely
\cref{asser:lean-exists,asser:lean-least,asser:lean-irrational,asser:lean-rh,asser:lean-detailed}.
\item The proofs in \cref{sec:proofs}
\item The scripts in \cref{sec:scripts}: Lean scripts in
\cref{sec:scripts-lean}, Python scripts in \cref{sec:scripts-python}
\item The appendices in \cref{sec:appendix}: tables in
\cref{sec:appendix-tables}, figures in \cref{sec:appendix-figures}.
\end{itemize}

The geometric content of the recurrence is shown in
\Cref{fig:mills-recurrence} (in \cref{sec:appendix-figures}). The assertions
\cref{asser:millsrec,asser:uv} formalise the recurrence on primes and the
two monotone sequences whose common limit defines $\theta $.

The Mills recurrence is
\begin{ineq}
\label{eq:mills-seq-intro}  q_n^3<q_{n+1}<(q_n+1)^3-1\qquad(n=1,2,\dots),
\end{ineq}
and the limit $\theta $ is obtained as the common limit of
\begin{equation}
u_{n}=q_{n}^{3^{-n}}\quad \text{and}\quad v_{n}=(q_{n}+1)^{3^{-n}}.
\end{equation}

\section{Definitions}

\label{sec:defs}

Every mathematical object in this document is either a \emph{type}, an \emph{%
operator}, or a \emph{predicate}, introduced in a numbered \texttt{definition%
} block. Each block has two parts:

\begin{itemize}
\item A \texttt{\textsf{Define:}} clause, followed by a numbered \texttt{type%
} block giving the signature
\item An \texttt{\textsf{Assume:}} clause, followed by a numbered \texttt{%
asser} block giving the axiom or defining equivalence.
\end{itemize}

\subsection{Primitive types}

\label{sec:defs-prim-types}

\begin{definition}
\label{def:prop} \textsf{Define:} the \cref{type:PropType} of propositions.
\begin{type}
\label{type:PropType} \textsf{Prop} : \textsf{Type}
\end{type}
\textsf{Assume:} \cref{asser:prop-axioms} (propositions are closed under the
standard logical connectives).
\begin{asser}
\label{asser:prop-axioms} \textsf{ClosedUnder}(\textsf{Prop}, \{\wedge,
\vee, \neg, \Rightarrow, \forall, \exists , \iff, =_{X}\})
\end{asser}
\end{definition}

\begin{definition}
\label{def:nat} \textsf{Define:} the \cref{type:NatType} of natural numbers.
\begin{type}
\label{type:NatType} \mathbb{N} : \textsf{Type}
\end{type}
\textsf{Assume:} \cref{asser:nat-axioms} (the Peano axioms for $\mathbb{N}$).
\begin{asser}
\label{asser:nat-axioms} \textsf{Peano}(\mathbb{N}, 0, \textsf{succ})
\end{asser}
\end{definition}

\begin{definition}
\label{def:int} \textsf{Define:} the \cref{type:IntType} of integers.
\begin{type}
\label{type:IntType} \mathbb{Z} : \textsf{Type}
\end{type}
\textsf{Assume:} $\mathbb{Z}$ is the initial integral domain of
characteristic zero (\cref{asser:int-axioms}).
\begin{asser}
\label{asser:int-axioms} \textsf{IntegralDomain}(\mathbb{Z}, +, \cdot, 0,
1) \wedge \textsf{Initial}(\mathbb{Z}) \wedge \textsf{char}(\mathbb{Z})
= 0
\end{asser}
\end{definition}

\begin{definition}
\label{def:real} \textsf{Define:} the \cref{type:RealType} of real numbers.
\begin{type}
\label{type:RealType} \mathbb{R} : \textsf{Type}
\end{type}
\textsf{Assume:} $\mathbb{R}$ is the Dedekind-complete ordered field (%
\cref{asser:real-axioms}).
\begin{asser}
\label{asser:real-axioms}
\begin{aligned}
&\textsf{Field}(\mathbb{R}, +, \cdot, 0, 1) \wedge \textsf{LinearOrder}(\mathbb{R}, <) \\
&\quad \wedge \textsf{OrderCompatible}(+, \cdot, <) \wedge \textsf{DedekindComplete}(\mathbb{R}, <)
\end{aligned}
\end{asser}
\end{definition}

\begin{definition}
\label{def:complex} \textsf{Define:} the \cref{type:ComplexType} of complex
numbers.
\begin{type}
\label{type:ComplexType} \mathbb{C} : \textsf{Type}
\end{type}
\textsf{Assume:} $\mathbb{C}$ is the algebraic closure of $\mathbb{R}$ (%
\cref{asser:complex-axioms}).
\begin{asser}
\label{asser:complex-axioms} \textsf{AlgebraicClosure}(\mathbb{C},
\mathbb{R}) \wedge \mathbb{R} \subseteq \mathbb{C}
\end{asser}
\end{definition}

\subsection{Primitive operators}

\label{sec:defs-prim-ops}

\begin{definition}
\label{def:exp-prim} \textsf{Define:} the \cref{type:ExpPrimType} of the
exponential function.
\begin{type}
\label{type:ExpPrimType} \exp : \mathbb{R} \to \mathbb{R}
\end{type}
\textsf{Assume:} \cref{asser:exp-prim-axioms} characterises $\exp$.
\begin{asser}
\label{asser:exp-prim-axioms}
\begin{aligned}
&\exp(0) = 1 \wedge \bigl(\forall x : \mathbb{R},\ \exp'(x) = \exp(x)\bigr) \\
&\quad \wedge \bigl(\forall x, y : \mathbb{R},\ \exp(x+y) = \exp(x)\cdot\exp(y)\bigr)
\end{aligned}
\end{asser}
\end{definition}

\begin{definition}
\label{def:log-prim} \textsf{Define:} the \cref{type:LogPrimType} of the
natural logarithm.
\begin{type}
\label{type:LogPrimType} \log : \mathbb{R}_{>0} \to \mathbb{R}
\end{type}
\textsf{Assume:} $\log$ is the two-sided inverse of $\exp$ (%
\cref{asser:log-prim-axioms}).
\begin{asser}
\label{asser:log-prim-axioms} \bigl(\forall x : \mathbb{R},\ \log(\exp(x))
= x\bigr) \wedge \bigl(\forall y : \mathbb{R}_{>0},\ \exp(\log(y)) = y
\bigr)
\end{asser}
\end{definition}

\begin{definition}
\label{def:exp2-prim} \textsf{Define:} the \cref{type:Exp2PrimType} of the
base-2 exponential.
\begin{type}
\label{type:Exp2PrimType} \exp_2 : \mathbb{R} \to \mathbb{R}
\end{type}
\textsf{Assume:} $\exp_2(x) := \exp(x \log 2)$ (\cref{asser:exp2-prim-axioms}).
\begin{asser}
\label{asser:exp2-prim-axioms} \forall x : \mathbb{R},\ \exp_2(x) = \exp(x \log 2)
\end{asser}
\end{definition}

\begin{definition}
\label{def:log2-prim} \textsf{Define:} the \cref{type:Log2PrimType} of the
base-2 logarithm.
\begin{type}
\label{type:Log2PrimType} \log_2 : \mathbb{R}_{>0} \to \mathbb{R}
\end{type}
\textsf{Assume:} $\log_2$ is the two-sided inverse of $\exp_2$ (%
\cref{asser:log2-prim-axioms}).
\begin{asser}
\label{asser:log2-prim-axioms} \bigl(\forall x : \mathbb{R},\ \log_2(\exp_2(x))
= x\bigr) \wedge \bigl(\forall y : \mathbb{R}_{>0},\ \exp_2(\log_2(y)) = y
\bigr)
\end{asser}
\end{definition}

\begin{definition}
\label{def:logc-prim} \textsf{Define:} the \cref{type:LogCPrimType} of the
base-$c$ logarithm, for $c : \mathbb{R}_{>0, \neq 1}$.
\begin{type}
\label{type:LogCPrimType} \log_c : \mathbb{R}_{>0} \to \mathbb{R}
\end{type}
\textsf{Assume:} $\log_c(x) := \log(x)/\log(c)$ (%
\cref{asser:logc-prim-axioms}).
\begin{asser}
\label{asser:logc-prim-axioms} \forall c : \mathbb{R}_{>0, \neq 1},\ \forall x : \mathbb{R}_{>0},\ \log_c(x) = \frac{\log x}{\log c}
\end{asser}
\end{definition}

\begin{definition}
\label{def:zeta} \textsf{Define:} the \cref{type:ZetaType} of the Riemann
zeta function.
\begin{type}
\label{type:ZetaType} \zeta : \mathbb{C} \to \mathbb{C}
\end{type}
\textsf{Assume:} $\zeta$ is the analytic continuation of the Dirichlet
series (\cref{asser:zeta-axioms}).
\begin{asser}
	\label{asser:zeta-axioms}
	\forall s : \mathbb{C},\ \Re(s) > 1 \Rightarrow
	\zeta(s) = \sum_{n\in\mathbb{N}_{>0}} n^{-s}
\end{asser}
\end{definition}

\begin{definition}
\label{def:re} \textsf{Define:} the \cref{type:ReType} of the real-part
operator.
\begin{type}
\label{type:ReType} \Re : \mathbb{C} \to \mathbb{R}
\end{type}
\textsf{Assume:} \cref{asser:re-axioms} characterises $\Re$.
\begin{asser}
\label{asser:re-axioms} \forall a, b : \mathbb{R},\ \Re(a + bi) = a
\end{asser}
\end{definition}

\begin{definition}
\label{def:continuous} \textsf{Define:} the \cref{type:ContinuousType} of
the continuity predicate.
\begin{type}
\label{type:ContinuousType} \textsf{Continuous} : (\mathbb{R} \to
\mathbb{R}) \to \textsf{Prop}
\end{type}
\textsf{Assume:} \cref{asser:continuous-axioms} (the $\varepsilon$-$\delta$
characterisation).
\begin{asser}
\label{asser:continuous-axioms}
\begin{aligned}
&\forall f : \mathbb{R} \to \mathbb{R},\ \textsf{Continuous}(f) :\iff \\
&\qquad \forall x : \mathbb{R},\ \forall \varepsilon : \mathbb{R}_{>0},\ \exists  \delta : \mathbb{R}_{>0},\ \forall y : \mathbb{R}, \\
&\qquad\qquad |y - x| < \delta \Rightarrow |f(y) - f(x)| < \varepsilon
\end{aligned}
\end{asser}
\end{definition}

\begin{definition}
\label{def:eventually-increasing} \textsf{Define:} the %
\cref{type:EventuallyIncreasingType} of the eventual-increase predicate.
\begin{type}
\label{type:EventuallyIncreasingType} \textsf{EventuallyIncreasing} : (
\mathbb{R} \to \mathbb{R}) \to \textsf{Prop}
\end{type}
\textsf{Assume:} \cref{asser:eventually-increasing-axioms} (eventual strict
monotonicity).
\begin{asser}
\label{asser:eventually-increasing-axioms}
\begin{aligned}
&\forall f : \mathbb{R} \to \mathbb{R},\ \textsf{EventuallyIncreasing}(f) :\iff \\
&\qquad \exists  N : \mathbb{R},\ \forall x, y : \mathbb{R},\ N < x < y \Rightarrow f(x) < f(y)
\end{aligned}
\end{asser}
\end{definition}

\begin{definition}
\label{def:bijective} \textsf{Define:} the \cref{type:BijectiveType} of the
bijectivity predicate.
\begin{type}
\label{type:BijectiveType} \textsf{Bijective} : \forall \alpha, \beta :
\textsf{Type},\ (\alpha \to \beta) \to \textsf{Prop}
\end{type}
\textsf{Assume:} \cref{asser:bijective-axioms} (injectivity together with
surjectivity).
\begin{asser}
\label{asser:bijective-axioms}
\begin{aligned}
&\forall \alpha, \beta : \textsf{Type},\ \forall f : \alpha \to \beta,\ \textsf{Bijective}(f) :\iff \\
&\qquad \bigl(\forall x, y : \alpha,\ f(x) = f(y) \Rightarrow x = y\bigr) \\
&\qquad \wedge \bigl(\forall y : \beta,\ \exists  x : \alpha,\ f(x) = y\bigr)
\end{aligned}
\end{asser}
\end{definition}

\begin{definition}
\label{def:measure} \textsf{Define:} the \cref{type:MeasureType} of Lebesgue
measure on $\mathbb{R}$.
\begin{type}
	\label{type:MeasureType} \mu : (\mathbb{R} \to \textsf{Prop}) \to
	\mathbb{R}_{\geq 0}
\end{type}
\textsf{Assume:} $\mu$ is the standard Lebesgue measure (%
\cref{asser:measure-axioms}).
\begin{asser}
\label{asser:measure-axioms} \textsf{LebesgueMeasure}(\mu, \mathbb{R})
\end{asser}
\end{definition}

\begin{definition}
\label{def:interior} \textsf{Define:} the \cref{type:InteriorType} of the
interior operator.
\begin{type}
\label{type:InteriorType} \textsf{Int} : (\mathbb{R} \to \textsf{Prop})
\to (\mathbb{R} \to \textsf{Prop})
\end{type}
\textsf{Assume:} \cref{asser:interior-axioms} (interior is the maximal open
subset).
\begin{asser}
\label{asser:interior-axioms}
\begin{aligned}
&\forall S : \mathbb{R} \to \textsf{Prop},\ \textsf{Int}(S) \subseteq S \wedge \textsf{IsOpen}(\textsf{Int}(S)) \\
&\quad \wedge \forall U : \mathbb{R} \to \textsf{Prop},\ \bigl(\textsf{IsOpen}(U) \wedge U \subseteq S\bigr) \Rightarrow U \subseteq \textsf{Int}(S)
\end{aligned}
\end{asser}
\end{definition}

\begin{definition}
\label{def:closure} \textsf{Define:} the \cref{type:ClosureType} of the
closure operator.
\begin{type}
\label{type:ClosureType} \textsf{cl} : (\mathbb{R} \to \textsf{Prop}) \to
( \mathbb{R} \to \textsf{Prop})
\end{type}
\textsf{Assume:} \cref{asser:closure-axioms} (closure is the minimal closed
superset).
\begin{asser}
\label{asser:closure-axioms}
\begin{aligned}
&\forall S : \mathbb{R} \to \textsf{Prop},\ S \subseteq \textsf{cl}(S) \wedge \textsf{IsClosed}(\textsf{cl}(S)) \\
&\quad \wedge \forall F : \mathbb{R} \to \textsf{Prop},\ \bigl(\textsf{IsClosed}(F) \wedge S \subseteq F\bigr) \Rightarrow \textsf{cl}(S) \subseteq F
\end{aligned}
\end{asser}
\end{definition}

\subsection{Prime and floor}

\label{sec:defs-prim-pred}

\begin{definition}
\label{def:prime} \textsf{Define:} the \cref{type:PrimeType} of the
computable prime numbers.
\begin{type}
\label{type:PrimeType} \textsf{Prime} \colon \mathbb{N} \to \textsf{Prop}
\end{type}
\textsf{Assume:} the predicate \cref{asser:prime} holds.
\begin{asser}
\label{asser:prime} \forall p : \mathbb{N},\ \textsf{Prime}(p) :\iff (p
\geq 2) \wedge \bigl(\forall d : \mathbb{N},\ (d \geq 2 \wedge d \mid p)
\Rightarrow d = p\bigr)
\end{asser}
\end{definition}

\begin{definition}
\label{def:floor} \textsf{Define:} the \cref{type:FloorType} of the
computable integer part of a number.
\begin{type}
\label{type:FloorType} \textsf{floor} \colon \mathbb{R} \to \mathbb{Z}
\end{type}
\textsf{Assume:} the predicate \cref{asser:floor} holds.
\begin{asser}
\label{asser:floor} \forall x : \mathbb{R},\ \forall n : \mathbb{Z},\
\textsf{floor}(x) = n \iff (n \leq x < n + 1)
\end{asser}
\end{definition}

\subsection{Sequences and Mills predicates}

\label{sec:defs-mills}

\begin{definition}
\label{def:pseq} \textsf{Define:} the \cref{type:PSeqType} of strictly
increasing prime-valued sequences.
\begin{type}
\label{type:PSeqType} \textsf{PSeq} \colon (\mathbb{N} \to \mathbb{N})
\to \textsf{Prop}
\end{type}
\textsf{Assume:} the predicate \cref{asser:pseq} holds.
\begin{asser}
\label{asser:pseq}
\begin{aligned}
&\forall p : \mathbb{N} \to \mathbb{N},\ \textsf{PSeq}(p) :\iff \bigl(p_1 = 2\bigr) \\
&\quad \wedge \bigl(\forall n : \mathbb{N}_{\geq 1},\ \textsf{Prime}(p_n)\bigr) \\
&\quad \wedge \bigl(\forall n : \mathbb{N}_{\geq 1},\ p_n < p_{n+1}\bigr)
\end{aligned}
\end{asser}
\end{definition}

\begin{definition}
\label{def:ismills} \textsf{Define:} the \cref{type:IsMillsType} of the
Mills predicate on reals.
\begin{type}
\label{type:IsMillsType} \textsf{IsMills} \colon \mathbb{R} \to \textsf{
Prop}
\end{type}
\textsf{Assume:} the predicate \cref{asser:ismills} holds.
\begin{asser}
\label{asser:ismills} \forall A : \mathbb{R},\ \textsf{IsMills}(A) :\iff
(A > 1) \wedge \bigl(\forall n : \mathbb{N}_{\geq 1},\ \textsf{Prime}(
\textsf{floor}(A^{3^{n}}))\bigr)
\end{asser}
\end{definition}

\begin{definition}
\label{def:isminmills} \textsf{Define:} the \cref{type:IsMinMillsType} of
the least-Mills predicate.
\begin{type}
\label{type:IsMinMillsType} \textsf{IsMinMills} \colon \mathbb{R} \to
\textsf{Prop}
\end{type}
\textsf{Assume:} the predicate \cref{asser:isminmills} holds.
\begin{asser}
\label{asser:isminmills} \forall A : \mathbb{R},\ \textsf{IsMinMills}(A)
:\iff \textsf{IsMills}(A) \wedge \bigl(\forall B : \mathbb{R},\ \textsf{
IsMills}(B) \Rightarrow A \leq B\bigr)
\end{asser}
\end{definition}

\begin{definition}
\label{def:millsrec} \textsf{Define:} the \cref{type:MillsRecType} of the
Mills recurrence on prime sequences.
\begin{type}
\label{type:MillsRecType} \textsf{MillsRec} \colon (\mathbb{N} \to
\mathbb{N}) \to \textsf{Prop}
\end{type}
\textsf{Assume:} the predicate \cref{asser:millsrec} holds.
\begin{asser}
\label{asser:millsrec} \forall q : \mathbb{N} \to \mathbb{N},\ \textsf{
MillsRec}(q) :\iff \forall n : \mathbb{N}_{\geq 1},\ q_n^{3} < q_{n+1} <
(q_n + 1)^{3} - 1
\end{asser}
\end{definition}

\begin{definition}
\label{def:uv} \textsf{Define:} the \cref{type:UVType} of the bracketing
pair $(u,v)$ for a sequence $q$.
\begin{type}
\label{type:UVType} \textsf{UV} \colon (\mathbb{N} \to \mathbb{N})
\times (\mathbb{N} \to \mathbb{R}) \times (\mathbb{N} \to \mathbb{R}%
) \to \textsf{Prop}
\end{type}
\textsf{Assume:} the predicate \cref{asser:uv} holds.
\begin{asser}
\label{asser:uv}
\begin{aligned}
&\forall q : \mathbb{N} \to \mathbb{N},\ \forall u, v : \mathbb{N} \to \mathbb{R},\ \textsf{UV}(q, u, v) :\iff \\
&\qquad \forall n : \mathbb{N}_{\geq 1},\ \bigl(u_n = q_n^{3^{-n}}\bigr) \wedge \bigl(v_n = (q_n + 1)^{3^{-n}}\bigr)
\end{aligned}
\end{asser}
\end{definition}

\subsection{Kuipers, Ansari, Wright, Niven predicates}

\label{sec:defs-kanw}

\begin{definition}
\label{def:iskuipers} \textsf{Define:} the \cref{type:IsKuipersType} of the
Kuipers predicate.
\begin{type}
\label{type:IsKuipersType} \textsf{IsKuipers} \colon \mathbb{N}_{\geq 3}
\times \mathbb{R} \to \textsf{Prop}
\end{type}
\textsf{Assume:} the predicate \cref{asser:iskuipers} holds.
\begin{asser}
\label{asser:iskuipers}
\begin{aligned}
&\forall c : \mathbb{N}_{\geq 3},\ \forall A : \mathbb{R},\ \textsf{IsKuipers}(c, A) :\iff \\
&\qquad (A > 1) \wedge \bigl(\forall n : \mathbb{N}_{\geq 1},\ \textsf{Prime}(\textsf{floor}(A^{c^{n}}))\bigr)
\end{aligned}
\end{asser}
\end{definition}

\begin{definition}
\label{def:isansari} \textsf{Define:} the \cref{type:IsAnsariType} of the
Ansari predicate.
\begin{type}
\label{type:IsAnsariType} \textsf{IsAnsari} \colon \mathbb{R}_{>8/3}
\times \mathbb{R} \to \textsf{Prop}
\end{type}
\textsf{Assume:} the predicate \cref{asser:isansari} holds.
\begin{asser}
\label{asser:isansari}
\begin{aligned}
&\forall c : \mathbb{R}_{>8/3},\ \forall A : \mathbb{R},\ \textsf{IsAnsari}(c, A) :\iff \\
&\qquad (A > 1) \wedge \bigl(\forall n : \mathbb{N}_{\geq 1},\ \textsf{Prime}(\textsf{floor}(A^{c^{n}}))\bigr)
\end{aligned}
\end{asser}
\end{definition}

\begin{definition}
\label{def:iswright} \textsf{Define:} the \cref{type:IsWrightType} of the
Wright predicate (iterated base-2 exponential).
\begin{type}
\label{type:IsWrightType} \textsf{IsWright} \colon \mathbb{R} \to \textsf{Prop}
\end{type}
\textsf{Assume:} the predicate \cref{asser:iswright} holds.
\begin{asser}
\label{asser:iswright} \forall A : \mathbb{R},\ \textsf{IsWright}(A) :\iff
(A > 1) \wedge \bigl(\forall n : \mathbb{N}_{\geq 1},\ \textsf{Prime}(
\textsf{floor}(\textsf{exp}_2^{[n]}(A)))\bigr)
\end{asser}
\end{definition}

\begin{definition}
\label{def:isniven} \textsf{Define:} the \cref{type:IsNivenType} of the
Niven predicate.
\begin{type}
\label{type:IsNivenType} \textsf{IsNiven} \colon \mathbb{R}_{>1} \times
\mathbb{R} \to \textsf{Prop}
\end{type}
\textsf{Assume:} the predicate \cref{asser:isniven} holds.
\begin{asser}
\label{asser:isniven} \forall c : \mathbb{R}_{>1},\ \forall \theta :
\mathbb{R},\ \textsf{IsNiven}(c, \theta) :\iff (\theta > 0) \wedge \bigl(
\forall n : \mathbb{N}_{\geq 1},\ \textsf{Prime}(\textsf{floor}
(c^{\theta^{n}}))\bigr)
\end{asser}
\end{definition}

\subsection{Ore's selection predicates}

\label{sec:defs-ore}

\begin{definition}
\label{def:isoresel} \textsf{Define:} the \cref{type:IsOreSelType} of the
Ore selection predicate.
\begin{type}
\label{type:IsOreSelType} \textsf{IsOreSel} \colon (\mathbb{R} \to
\mathbb{R}) \to \textsf{Prop}
\end{type}
\textsf{Assume:} the predicate \cref{asser:isoresel} holds.
\begin{asser}
\label{asser:isoresel} \forall f : \mathbb{R} \to \mathbb{R},\ \textsf{
IsOreSel}(f) :\iff \textsf{Continuous}(f) \wedge \textsf{EventuallyIncreasing%
}(f)
\end{asser}
\end{definition}

\begin{definition}
\label{def:orebound} \textsf{Define:} the \cref{type:OreBoundType} of the
Ore bracketing predicate.
\begin{type}
\label{type:OreBoundType} \textsf{OreBound} \colon (\mathbb{R} \to
\mathbb{R}) \times (\mathbb{N} \to \mathbb{N}) \times (\mathbb{N}
\to \mathbb{R}_{\geq 0}) \times (\mathbb{N} \to \mathbb{R}_{\geq 0})
\to \textsf{Prop}
\end{type}
\textsf{Assume:} the predicate \cref{asser:orebound} holds.
\begin{asser}
\label{asser:orebound}
\begin{aligned}
&\forall f : \mathbb{R} \to \mathbb{R},\ \forall q : \mathbb{N} \to \mathbb{N},\ \forall \delta, \epsilon : \mathbb{N} \to \mathbb{R}_{\geq 0},\ \\
&\qquad \textsf{OreBound}(f, q, \delta, \epsilon) :\iff \\
&\qquad \forall n : \mathbb{N}_{\geq 1},\ f(q_n - \delta_n) < q_{n+1} - \delta_{n+1} \\
&\qquad\qquad \wedge q_{n+1} + \epsilon_{n+1} < f(q_n + \epsilon_n)
\end{aligned}
\end{asser}
\end{definition}

\subsection{Polynomial-gap predicate}

\label{sec:defs-polygap}

\begin{definition}
\label{def:polygap} \textsf{Define:} the \cref{type:PolyGapType} of the
polynomial prime-gap predicate.
\begin{type}
\label{type:PolyGapType} \textsf{PolyGap} \colon (0,1) \times \mathbb{R}
_{>0} \to \textsf{Prop}
\end{type}
\textsf{Assume:} the predicate \cref{asser:polygap} holds.
\begin{asser}
\label{asser:polygap} \forall \alpha : (0,1),\ \forall k : \mathbb{R}
_{>0},\ \textsf{PolyGap}(\alpha, k) :\iff \forall n : \mathbb{N}_{\geq
1},\ p_{n+1} - p_n < k\, p_n^{\alpha}
\end{asser}
\end{definition}

\subsection{Null, nowhere-dense, continuum, irrational, RH}

\label{sec:defs-rm}

\begin{definition}
\label{def:isnull} \textsf{Define:} the \cref{type:IsNullType} of the
null-set predicate.
\begin{type}
\label{type:IsNullType} \textsf{IsNull} \colon (\mathbb{R} \to \textsf{Prop}) \to \textsf{Prop}
\end{type}
\textsf{Assume:} \cref{asser:isnull} (null sets are exactly those of
Lebesgue measure zero).
\begin{asser}
\label{asser:isnull} \forall S : \mathbb{R} \to \textsf{Prop},\ \textsf{
IsNull}(S) :\iff \mu(S) = 0
\end{asser}
\end{definition}

\begin{definition}
\label{def:nowheredense} \textsf{Define:} the \cref{type:NowhereDenseType}
of the nowhere-dense predicate.
\begin{type}
\label{type:NowhereDenseType} \textsf{NowhereDense} \colon (\mathbb{R} \to
\textsf{Prop}) \to \textsf{Prop}
\end{type}
\textsf{Assume:} \cref{asser:nowheredense} (nowhere-dense sets have empty
interior of closure).
\begin{asser}
\label{asser:nowheredense} \forall S : \mathbb{R} \to \textsf{Prop},\
\textsf{NowhereDense}(S) :\iff \textsf{Int}(\textsf{cl}(S)) = \varnothing
\end{asser}
\end{definition}

\begin{definition}
\label{def:continuumcard} \textsf{Define:} the \cref{type:ContinuumCardType}
of the continuum-cardinality predicate.
\begin{type}
\label{type:ContinuumCardType} \textsf{ContinuumCard} \colon (\mathbb{R}
\to \textsf{Prop}) \to \textsf{Prop}
\end{type}
\textsf{Assume:} the predicate \cref{asser:continuumcard} holds.
\begin{asser}
\label{asser:continuumcard}
\begin{aligned}
&\forall P : \mathbb{R} \to \textsf{Prop},\ \textsf{ContinuumCard}(P) :\iff \\
&\qquad \exists  f : \mathbb{R} \to (\Sigma\, x : \mathbb{R},\ P(x)),\ \textsf{Bijective}(f)
\end{aligned}
\end{asser}
\end{definition}

\begin{definition}
\label{def:irrational} \textsf{Define:} the \cref{type:IrrationalType} of
the irrationality predicate.
\begin{type}
\label{type:IrrationalType} \textsf{Irrational} \colon \mathbb{R} \to
\textsf{Prop}
\end{type}
\textsf{Assume:} the predicate \cref{asser:irrational} holds.
\begin{asser}
\label{asser:irrational} \forall A : \mathbb{R},\ \textsf{Irrational}(A)
:\iff \neg\bigl(\exists p : \mathbb{Z},\ \exists q : \mathbb{N}_{\geq
1},\ A = p/q\bigr)
\end{asser}
\end{definition}

\begin{definition}
\label{def:rh} \textsf{Define:} the \cref{type:RHType} of the
Riemann-hypothesis proposition.
\begin{type}
\label{type:RHType} \textsf{RH} : \textsf{Prop}
\end{type}
\textsf{Assume:} \cref{asser:rh} states the Riemann hypothesis.
\begin{asser}
\label{asser:rh} \textsf{RH} :\iff \forall s : \mathbb{C},\ \bigl(\zeta(s)
= 0 \wedge \Re(s) > 1/2\bigr) \Rightarrow \Re(s) = 1/2
\end{asser}
\end{definition}

\subsection{Theta predicates for the metrical theorem}

\label{sec:defs-theta}

\begin{definition}
\label{def:theta-ansari} \textsf{Define:} the \cref{type:ThetaAnsariType} of
the $\theta$-predicate for Ansari.
\begin{type}
\label{type:ThetaAnsariType} \textsf{ThetaAnsari} \colon \mathbb{R} \to
\textsf{Prop}
\end{type}
\textsf{Assume:} the predicate \cref{asser:theta-ansari} holds.
\begin{asser}
\label{asser:theta-ansari} \forall \theta : \mathbb{R},\ \textsf{
ThetaAnsari}(\theta) :\iff \exists c : \mathbb{R}_{>8/3},\ \textsf{IsAnsari} (c, \theta)
\end{asser}
\end{definition}

\begin{definition}
\label{def:theta-wright} \textsf{Define:} the \cref{type:ThetaWrightType} of
the $\theta$-predicate for Wright.
\begin{type}
\label{type:ThetaWrightType} \textsf{ThetaWright} \colon \mathbb{R} \to
\textsf{Prop}
\end{type}
\textsf{Assume:} the predicate \cref{asser:theta-wright} holds.
\begin{asser}
\label{asser:theta-wright} \forall \theta : \mathbb{R},\ \textsf{
ThetaWright}(\theta) :\iff \textsf{IsWright}(\theta)
\end{asser}
\end{definition}

\begin{definition}
\label{def:theta-niven} \textsf{Define:} the \cref{type:ThetaNivenType} of
the $\theta $-predicate for Niven.
\begin{type}
\label{type:ThetaNivenType} \textsf{ThetaNiven} \colon \mathbb{R} \to
\textsf{Prop}
\end{type}
\textsf{Assume:} the predicate \cref{asser:theta-niven} holds.
\begin{asser}
\label{asser:theta-niven} \forall \theta : \mathbb{R},\ \textsf{ThetaNiven}
(\theta) :\iff \exists c : \mathbb{R}_{>1},\ \textsf{IsNiven}(c, \theta)
\end{asser}
\end{definition}

\subsection{Iterated operators}

\label{sec:defs-iter}

\begin{definition}
\label{def:expiter} \textsf{Define:} for $n : \mathbb{N}$, the %
\cref{type:ExpIterType} of the $n$-fold iterated base-2 exponential.
\begin{type}
\label{type:ExpIterType} \textsf{exp}_2^{[n]} \colon \mathbb{R} \to
\mathbb{R}
\end{type}
\textsf{Assume:} the predicate \cref{asser:expiter} holds.
\begin{asser}
\label{asser:expiter} \bigl(\forall x : \mathbb{R},\ \textsf{exp}_2^{[0]}(x)
= x\bigr) \wedge \bigl(\forall n : \mathbb{N},\ \textsf{exp}_2^{[n+1]} =
\textsf{exp}_2 \circ \textsf{exp}_2^{[n]}\bigr)
\end{asser}
\end{definition}

\begin{definition}
\label{def:logiter} \textsf{Define:} for $n : \mathbb{N}$, the %
\cref{type:LogIterType} of the $n$-fold iterated base-2 logarithm.
\begin{type}
\label{type:LogIterType} \textsf{log}_2^{[n]} \colon \mathbb{R}_{>0} \to
\mathbb{R}
\end{type}
\textsf{Assume:} the predicate \cref{asser:logiter} holds.
\begin{asser}
\label{asser:logiter} \bigl(\forall x : \mathbb{R}_{>0},\ \textsf{log}_2
^{[0]}(x) = x\bigr) \wedge \bigl(\forall n : \mathbb{N},\ \textsf{log}_2
^{[n+1]} = \textsf{log}_2^{[n]} \circ \textsf{log}_2\bigr)
\end{asser}
\end{definition}

\section{Axioms}

\label{sec:axioms}

\subsection{Reverse-mathematical systems}

The following table fixes the systems used throughout the classification (%
\cref{sec:thm-summary}).

{\small
\noindent \newline
\noindent
\begin{tabular}{|c|m{6.0cm}|m{3.0cm}|m{4.0cm}|}
	\hline
	\textbf{Tag} & \textbf{Full name} & \textbf{System} & \textbf{Definition} \\ \hline
	$\mathrm{RCA}_0$ & Recursive Comprehension Axiom & Recursive CA & $\mathrm{FOL}+\Delta^0_1\text{-CA}+\Sigma^0_1\text{-Ind}$ \\
\end{tabular}%
\noindent \newline
\noindent
\begin{tabular}{|c|m{6.0cm}|m{3.0cm}|m{4.0cm}|}
	\hline
	$\mathrm{WKL}_0$ & Weak K\"onig's Lemma & Weak K\"onig & $\mathrm{RCA}_0+\mathrm{WKL}$ \\ \hline
	$\mathrm{ACA}_0$ & Arithmetical Comprehension Axiom & Arithmetical CA & $\mathrm{RCA}_0+\text{Arithmetical CA}$ \\
\end{tabular}%
\noindent \newline
\noindent
\begin{tabular}{|c|m{6.0cm}|m{3.0cm}|m{4.0cm}|}
	\hline
	Ingham & Ingham's prime-gap estimate & Ingham's estimate & \Cref{ineq:ingham} \\
\end{tabular}%
}

\subsection{Analytic input}

The analytic input common to Routes 1, 2, 3, 7 is Ingham's estimate.

\begin{ineq}
\label{ineq:ingham} \textsf{Ingham} :\iff \exists k : \mathbb{R}_{>0},\
\forall n : \mathbb{N}_{\geq 1},\ p_{n+1} - p_n < k\, p_n^{5/8}
\end{ineq}

\subsection{Elementary input}

The elementary input used by Route 4 (and by Route 6 with a general
polynomial gap) is Bertrand's postulate.

\begin{ineq}
\label{ineq:bertrand} \textsf{Bertrand} :\iff \forall N : \mathbb{N}_{\geq
2},\ \exists p : \mathbb{N},\ (N < p < 2N) \wedge \textsf{Prime}(p)
\end{ineq}
\subsection{Law of excluded middle}

The reverse-mathematical classification records \textsf{LEM} explicitly as a
hypothesis rather than building it into the definition of \textsf{Prop}.

\begin{asser}
	\label{asser:lem}
	\textsf{LEM} :\iff \forall P : \textsf{Prop},\ P \vee \neg P
\end{asser}
\section{Postulates}

\label{sec:postulates}

The hypotheses underpinning the constructions are collected in the following
table; each row refers to the formal statement in the place indicated.

{\small
\noindent \newline
\noindent
\begin{tabular}{|c|m{11.0cm}|m{2.0cm}|m{2.5cm}|}
	\hline
	\textbf{Tag} & \textbf{Statement} & \textbf{Name} & \textbf{Strength} \\ \hline
	Pos 1 & $\exists k : \mathbb{R}_{>0},\ \forall n : \mathbb{N}_{\geq 1},\ p_{n+1}-p_n < k\, p_n^{5/8}$ & Ingham & \cref{ineq:ingham} \\
\end{tabular}%
\noindent \newline
\noindent
\begin{tabular}{|c|m{11.0cm}|m{2.0cm}|m{2.5cm}|}
	\hline
	Pos 2 & $\forall n : \mathbb{N}_{\geq 1},\ p_{n+1}-p_n < p_n$ & Bertrand & $\mathrm{RCA}_0$ \\ \hline
	Pos 3 & $\exists k : \mathbb{R}_{>0},\ \exists \alpha : (0,1),\ \forall n,\ p_{n+1}-p_n < k\, p_n^{\alpha}$ & PolyGap & \cref{asser:polygap} \\
\end{tabular}%
\noindent \newline
\noindent
\begin{tabular}{|c|m{11.0cm}|m{2.0cm}|m{2.5cm}|}
	\hline
	Pos 4 & $\forall N \gg 1,\ \exists p : \mathbb{N},\ N^3 < p < (N+1)^3-1 \wedge \mathrm{Prime}(p)$ & MillsRec & \cref{ineq:ingham} \\ \hline
	Pos 5 & $\forall x : \mathbb{R},\ \forall n : \mathbb{Z},\ \mathrm{floor}(x)=n \iff n \leq x < n+1$ & Floor & \cref{asser:floor} \\
\end{tabular}%
\noindent \newline
\noindent
\begin{tabular}{|c|m{11.0cm}|m{2.0cm}|m{2.5cm}|}
	\hline
	Pos 6 & $\forall q,u,v,\ \mathrm{UV}(q,u,v) \wedge (\forall n,\ u_{n}<u_{n+1}<v_{n+1}<v_{n}) \wedge (v_{n}-u_{n}\rightarrow 0)$ & Cauchy & $\mathrm{RCA}_0$ \\ \hline
\end{tabular}%
}

\section{Theorems}

\label{sec:theorems}

Each theorem below states a hypothesis under \textbf{\textsf{Assume:}} and a
conclusion under \textbf{\textsf{Conclude:}}. The hypothesis is either
referenced (by \cref{ineq:ingham}, \cref{ineq:bertrand}, %
\cref{asser:lean-least}, etc.) or spelled out inline; the conclusion is
placed in a numbered \texttt{asser} (resp.\ \texttt{ineq}) block whose label
is the one used throughout the paper. Proofs are given in \cref{sec:proofs}.

\subsection{Prime-gap constructions}

\label{sec:thm-prime-gap}

\begin{theorem}[Mills \protect\cite{Mills1947}]
\label{thm:mills}
\textbf{\textsf{Assume:}} \cref{ineq:ingham} (\textsf{Ingham}'s estimate).
\textbf{\textsf{Conclude:}}
\begin{asser}
	\label{asser:mills} \exists \theta : \mathbb{R}_{>1},\ \textsf{IsMills}%
	(\theta)
\end{asser}

\end{theorem}

\begin{remark}[Transcription typo]
\label{rem:typo} The historical transcription of Dudley \cite{Dudley1969}
writes
\begin{equation}
u_{n}=q_{n}^{5^{-n}},\quad v_{n}=(q_{n}+1)^{5^{-n}},
\end{equation}
and concludes
\begin{equation}
q_{n}=[\theta ^{5^{n}}].
\end{equation}
That is a typo. With exponent $5^{-n}$ the sequence $\{u_{n}\}$ fails
monotonicity at $n=1$: since
\begin{ineq}
q_{n+1}>q_{n}^{3}
\end{ineq}
and
\begin{ineq}
3\cdot 5^{-(n+1)}=3/5\cdot 5^{-n}<5^{-n},
\end{ineq}
we get
\begin{equation}
u_{n+1}<u_{n}.
\end{equation}
Only the exponent $3^{-n}$ (matching the recurrence \cref{eq:mills-seq-intro}%
) restores monotonicity.\newline
See \Cref{sec:appendix-tables} for a numerical confirmation.
\end{remark}

\begin{theorem}[Monotonicity of UV]
	\label{thm:millsseq}
	\textbf{\textsf{Assume:}} 
	\begin{center}
$	q : \mathbb{N} \to \mathbb{N}$, $u, v : \mathbb{N} \to
	\mathbb{R}$ with \cref{asser:uv,asser:millsrec}.
	\end{center}
	\textbf{\textsf{Conclude:}}
\begin{ineq}
\label{ineq:millsseq} \forall n : \mathbb{N}_{\geq 1},\ u_{n}<
u_{n+1}<v_{n+1}<v_{n}
\end{ineq}
\end{theorem}

\begin{theorem}[Kuipers \protect\cite{Kuipers1950}]
\label{thm:kuipers}
\textbf{\textsf{Assume:}} \cref{ineq:ingham}.\\
\textbf{\textsf{Conclude:}}
\begin{asser}
\label{asser:kuipers} \forall c : \mathbb{N}_{\geq 3},\ \exists \theta:
\mathbb{R}_{>1},\ \textsf{IsKuipers}(c, \theta)
\end{asser}
\end{theorem}

\begin{theorem}[Ansari \protect\cite{Ansari1951}]
\label{thm:ansari}
\textbf{\textsf{Assume:}} \cref{ineq:ingham}.\\
\textbf{\textsf{Conclude:}}
\begin{asser}
\label{asser:ansari} \forall c : \mathbb{R}_{>8/3},\ \exists \theta :
\mathbb{R}_{>1},\ \textsf{IsAnsari}(c, \theta)
\end{asser}
\end{theorem}

\subsection{Elementary and abstract constructions}

\label{sec:thm-elementary}

\begin{theorem}[Wright \protect\cite{Wright1951}]
\label{thm:wright}
\textbf{\textsf{Assume:}} \cref{ineq:bertrand}.\\
\textbf{\textsf{Conclude:}}
\begin{asser}
\label{asser:wright} \exists \theta : \mathbb{R}_{>1},\ \textsf{IsWright}%
(\theta)
\end{asser}
\end{theorem}

\begin{theorem}[Ore \protect\cite{Ore1952}]
\label{thm:ore}
\textbf{\textsf{Assume:}} $f : \mathbb{R} \to \mathbb{R}$, $q : \mathbb{N}
\to \mathbb{N}$, $\delta, \epsilon : \mathbb{N} \to \mathbb{R}_{\geq 0}$
with \cref{asser:isoresel,asser:orebound}.\\
\textbf{\textsf{Conclude:}}
\begin{asser}
\label{asser:ore} \exists \theta : \mathbb{R}_{>0},\ \forall n : \mathbb{%
N}_{\geq 1},\ q_n - \delta_n \leq f^{[n]}(\theta) \leq q_n + \epsilon_n
\end{asser}
\end{theorem}

\begin{theorem}[Polynomial-gap construction]
\label{thm:direct}
\textbf{\textsf{Assume:}} $\alpha : (0,1)$, $k : \mathbb{R}_{>0}$ with %
\cref{asser:polygap}, and $c : \mathbb{R}_{>1/(1-\alpha)}$.\\
\textbf{\textsf{Conclude:}}
\begin{asser}
\label{asser:direct} \exists \theta : \mathbb{R}_{>0},\ \forall n :
\mathbb{N}_{\geq 1},\ \textsf{Prime}(\textsf{floor}(\theta^{c^{n}}))
\end{asser}
\end{theorem}

\begin{remark}
\label{rem:direct-valid}
The condition $c>1/(1-\alpha)$ is exactly what is needed so that the
polynomial gap $p_{n+1}-p_n < k\,p_n^\alpha$ lifts to a gap below
$N^{c-1}$ inside the interval $(N^c,(N+1)^c-1)$. Taking $c=1/\alpha$,
as in some informal accounts, is insufficient unless $\alpha=1/2$.
\end{remark}

\subsection{Variants and metrical theorems}

\label{sec:thm-variants}

\begin{theorem}[Niven \protect\cite{Niven1951}]
\label{thm:niven}
\textbf{\textsf{Assume:}} \cref{ineq:ingham}.\\
\textbf{\textsf{Conclude:}}
\begin{asser}
\label{asser:niven} \forall c : \mathbb{R}_{>1},\ \exists \theta :
\mathbb{R}_{>0},\ \textsf{IsNiven}(c, \theta)
\end{asser}
\end{theorem}

\begin{theorem}[Wright \protect\cite{Wright1954}]
\label{thm:wright-metrical}
\textbf{\textsf{Assume:}} each of the three constructions in %
\cref{asser:ansari,asser:wright,asser:niven} yields, at every stage $n$, a
non-empty admissible interval for $q_{n+1}$.\\
\textbf{\textsf{Conclude:}}
\begin{asser}
	\label{asser:wright-metrical}
	\begin{array}{l}
		\mathsf{IsNull}(\mathsf{ThetaAnsari})\wedge
		\mathsf{NowhereDense}(\mathsf{ThetaAnsari}) \\
		\wedge \mathsf{ContinuumCard}(\mathsf{ThetaAnsari})\wedge
		\mathsf{IsNull}(\mathsf{ThetaWright}) \\
		\wedge \mathsf{NowhereDense}(\mathsf{ThetaWright})\wedge
		\mathsf{ContinuumCard}(\mathsf{ThetaWright}) \\
		\wedge \mathsf{IsNull}(\mathsf{ThetaNiven})\wedge
		\mathsf{NowhereDense}(\mathsf{ThetaNiven}) \\
		\wedge \mathsf{ContinuumCard}(\mathsf{ThetaNiven})
	\end{array}
\end{asser}
\end{theorem}

\subsection{The Lean formulation}

\label{sec:thm-lean}

\begin{theorem}[Lean-form Mills theorem]
\label{thm:lean-exists}
\textbf{\textsf{Assume:}} \cref{ineq:ingham}.\\
\textbf{\textsf{Conclude:}}
\begin{asser}
\label{asser:lean-exists} \exists A : \mathbb{R}_{>1},\ \textsf{IsMills}(A)
\end{asser}
\end{theorem}

\begin{theorem}[Existence of the least Mills constant]
\label{thm:lean-least}
\textbf{\textsf{Assume:}} \cref{asser:lean-exists}.\\
\textbf{\textsf{Conclude:}}
\begin{asser}
\label{asser:lean-least} \exists A : \mathbb{R},\ \textsf{IsMinMills}(A)
\end{asser}
\end{theorem}

\begin{theorem}[Irrationality of the least Mills constant, \protect\cite{Saito2025}]
\label{thm:lean-irrational}
\textbf{\textsf{Assume:}} \cref{asser:lean-least}.\\
\textbf{\textsf{Conclude:}}
\begin{asser}
\label{asser:lean-irrational} \forall A : \mathbb{R},\ \textsf{IsMinMills}%
(A) \Rightarrow \textsf{Irrational}(A)
\end{asser}
\end{theorem}

\begin{theorem}[Lower bound under RH,\protect\cite{CaldwellCheng2005}]
\label{thm:lean-rh}
\textbf{\textsf{Assume:}} \cref{asser:rh} and \cref{asser:lean-least}.\\
\textbf{\textsf{Conclude:}}
\begin{asser}
\label{asser:lean-rh} \forall A : \mathbb{R},\ \textsf{IsMinMills}(A)
\Rightarrow 1.3063778838 < A < 1.3063778839
\end{asser}
\end{theorem}

\begin{theorem}[Lean-adapted Mills limit]
\label{thm:lean-detailed}
\textbf{\textsf{Assume:}} $q : \mathbb{N} \to \mathbb{N}$ with %
\cref{asser:millsrec} and $q_1>k^{8}$ (for $k$ the constant of %
\cref{ineq:ingham}).\\
\textbf{\textsf{Conclude:}}
\begin{asser}
\label{asser:lean-detailed} \exists \theta : \mathbb{R}_{>0},\ \bigl(%
\forall n : \mathbb{N}_{\geq 1},\ q_n = \textsf{floor}(\theta^{3^{n}}%
)\bigr) \wedge \textsf{IsMills}(\theta)
\end{asser}
\end{theorem}

\subsection{Reverse-mathematical summary}

\label{sec:thm-summary}

{\small
\noindent \newline
\noindent
\begin{tabular}{|m{0.6cm}|m{6.4cm}|m{2.4cm}|m{3cm}|}
	\hline
	\textbf{\#} & \textbf{Mechanism} & \textbf{Reference} & \textbf{Strength} \\ \hline
	1 & Prime-gap estimate + $3^n$ iteration & \Cref{asser:mills} & $\mathrm{RCA}_0$/Ingham \\
\end{tabular}%
\noindent \newline
\noindent
\begin{tabular}{|m{0.6cm}|m{6.4cm}|m{2.4cm}|m{3cm}|}
	\hline
	2 & Prime-gap estimate + $c^n$ iteration & \Cref{asser:kuipers} & $\mathrm{RCA}_0$/Ingham \\ \hline
	3 & Prime-gap estimate + $c>8/3$ & \Cref{asser:ansari} & $\mathrm{RCA}_0$/Ingham \\
\end{tabular}%
\noindent \newline
\noindent
\begin{tabular}{|m{0.6cm}|m{6.4cm}|m{2.4cm}|m{3cm}|}
	\hline
	4 & Bertrand + $2$-tower & \Cref{asser:wright} & $\mathrm{RCA}_0$ \\ \hline
	5 & Abstract selection framework & \Cref{asser:ore} & $\mathrm{RCA}_0$ \\
\end{tabular}%
\noindent \newline
\noindent
\begin{tabular}{|m{0.6cm}|m{6.4cm}|m{2.4cm}|m{3cm}|}
	\hline
	6 & General polynomial gap ($c>1/(1-\alpha)$) & \Cref{asser:direct} & gap-dependent \\ \hline
	7 & Niven fixed-base variant & \Cref{asser:niven} & $\mathrm{RCA}_0$/Ingham \\
\end{tabular}%
\noindent \newline
\noindent
\begin{tabular}{|m{0.6cm}|m{6.4cm}|m{2.4cm}|m{3cm}|}
	\hline
	8 & Wright metrical theorem & \Cref{asser:wright-metrical} & $\mathrm{WKL}_0$ \\ \hline
	9 & Lean pointer & \Cref{sec:thm-lean} & formal \\
\end{tabular}%
\noindent \newline
\noindent
\begin{tabular}{|m{0.6cm}|m{6.4cm}|m{2.4cm}|m{3cm}|}
	\hline
	10 & Lean detailed adaptation & \Cref{asser:lean-detailed} & formal \\
\end{tabular}%
}

% =========================================================================
% =========================================================================
\section{Proofs}
\label{sec:proofs}
% =========================================================================

One starred subsection per route, in route order.\newline
Each proof is a column dependency table: first row \textsf{Hyp}, last row
$\rightarrow \mathsf{I}(1,\dots ,N-1)$.\newline
Justifications use only \textsf{Hyp}, \textsf{def}, \textsf{Ax}, \textsf{Pos},
\textsf{Def}, \textsf{Lem}, \textsf{Thm}, $\mathsf{Eq}(\dots )$, \textsf{Predicates},
or a logical inference rule
\begin{equation}
(\rightarrow \mathsf{I},\wedge \mathsf{E},\vee \mathsf{E},\forall \mathsf{I},
\forall \mathsf{E},\exists \mathsf{I},\exists \mathsf{E},\mathsf{RAA}),
\end{equation}
discharging the hypothesis and yielding the implication stated in
\cref{sec:theorems}.

% -------------------------------------------------------------------------
\subsection{Preliminary lemmas}
\label{sec:prelim}
% -------------------------------------------------------------------------

\begin{lemma}[Euclid]
	\label{lem:euclid}
	\textbf{\textsf{Assume:}}
	\begin{equation}
	N\in\mathbb{N}_{\geq 1}
	\end{equation}
	\textbf{\textsf{Conclude:}}
	\begin{asser}
		\label{asser:lem-euclid} \exists p\in\mathbb{N},\ (p>N)\wedge\mathrm{Prime}(p)
	\end{asser}
\end{lemma}

\begin{proof}[\textbf{Proof of \cref{lem:euclid}}]
	\phantomsection\label{proof:lem:euclid}
	{\scriptsize
		\begin{tabular}{|m{0.7cm}|m{12cm}|m{1.4cm}|m{2.5cm}|}
			\hline
			1 & $N\in\mathbb{N}_{\geq 1}$ & --- & $\mathsf{Hyp}$ \\ \hline
			2 & $M:=N!+1$ & 1 & $\mathsf{def}$ \\
		\end{tabular}%
		\noindent \newline
		\noindent
		\begin{tabular}{|m{0.7cm}|m{12cm}|m{1.4cm}|m{2.5cm}|}
			\hline
			3 & $M\ge 2$ & 2 & $\mathsf{Eq}$ \\ \hline
			4 & $\exists p$ prime, $p\mid M$ & 3 & $\mathsf{Support}$ (FTA) \\
		\end{tabular}%
		\noindent \newline
		\noindent
		\begin{tabular}{|m{0.7cm}|m{12cm}|m{1.4cm}|m{2.5cm}|}
			\hline
			5 & Suppose $p\le N$ & --- & $\mathsf{RAA}$ \\ \hline
			6 & $p\mid N!$ & 5 & $\mathsf{Support}$ (basic NT) \\
		\end{tabular}%
		\noindent \newline
		\noindent
		\begin{tabular}{|m{0.7cm}|m{12cm}|m{1.4cm}|m{2.5cm}|}
			\hline
			7 & $p\mid M$ & 4 & $\wedge\mathsf{E}$ \\ \hline
			8 & $p\mid M-N!=1$ & 6,7 & $\mathsf{Eq}$ \\
		\end{tabular}%
		\noindent \newline
		\noindent
		\begin{tabular}{|m{0.7cm}|m{12cm}|m{1.4cm}|m{2.5cm}|}
			\hline
			9 & $p\ge 2$ and $p\mid 1$: contradiction & 4,8 & $\mathsf{Eq}$ \\ \hline
			10 & Therefore $p>N$ & 5--9 & $\mathsf{RAA}$ \\
		\end{tabular}%
		\noindent \newline
		\noindent
		\begin{tabular}{|m{0.7cm}|m{12cm}|m{1.4cm}|m{2.5cm}|}
			\hline
			11 & $\exists p>N$ prime & 4,10 & $\exists\mathsf{I}$ \\ \hline
			12 & Lemma & 1--11 & $\rightarrow\mathsf{I}$ \\ \hline
		\end{tabular}%
	}
\end{proof}

\begin{lemma}[Power-mean slope]
	\label{lem:slope}
	\textbf{\textsf{Assume:}}
	\begin{equation}
	\beta>1 \quad\text{and}\quad \varepsilon>0.
	\end{equation}
	\textbf{\textsf{Conclude:}}
	\begin{asser}
		\label{asser:lem-slope} \exists N_0:\mathbb{N},\ \forall N:\mathbb{N},\ N\ge N_0\Rightarrow N^{\beta}+N^{\beta-1}<(N+1)^{\beta}-\varepsilon
	\end{asser}
\end{lemma}

\begin{proof}[\textbf{Proof of \cref{lem:slope}}]
	\phantomsection\label{proof:lem:slope}
	{\scriptsize
		\begin{tabular}{|m{0.7cm}|m{12cm}|m{1.4cm}|m{2.5cm}|}
			\hline
			1 & $\beta>1$, $\varepsilon>0$ & --- & $\mathsf{Hyp}$ \\ \hline
			2 & $g(N):=(N+1)^\beta-N^\beta-N^{\beta-1}$ & 1 & $\mathsf{def}$ \\
		\end{tabular}%
		\noindent \newline
		\noindent
		\begin{tabular}{|m{0.7cm}|m{12cm}|m{1.4cm}|m{2.5cm}|}
			\hline
			3 & MVT: $\exists\xi_N\in(N,N+1)$, $(N+1)^\beta-N^\beta=\beta\xi_N^{\beta-1}$ & 1 & $\mathsf{Support}$ (MVT) \\ \hline
			4 & $g(N)=\beta\xi_N^{\beta-1}-N^{\beta-1}$ & 2,3 & $\mathsf{Eq}$ \\
		\end{tabular}%
		\noindent \newline
		\noindent
		\begin{tabular}{|m{0.7cm}|m{12cm}|m{1.4cm}|m{2.5cm}|}
			\hline
			5 & $\beta-1>0$, $\xi_N>N$ $\Rightarrow$ $\xi_N^{\beta-1}>N^{\beta-1}$ & 1,3 & $\mathsf{Eq}$ \\ \hline
			6 & $g(N)>(\beta-1)N^{\beta-1}$ & 4,5 & $\mathsf{Eq}$ \\
		\end{tabular}%
		\noindent \newline
		\noindent
		\begin{tabular}{|m{0.7cm}|m{12cm}|m{1.4cm}|m{2.5cm}|}
			\hline
			7 & $\forall M:\mathbb{R},\ \exists N:\mathbb{N},\ (\beta-1)N^{\beta-1}>M$ & 1 & $\mathsf{Support}$ \\ \hline
			8 & $\exists N_0$, $\forall N\ge N_0$: $g(N)>\varepsilon$ & 6,7 & $\mathsf{def}$ \\
		\end{tabular}%
		\noindent \newline
		\noindent
		\begin{tabular}{|m{0.7cm}|m{12cm}|m{1.4cm}|m{2.5cm}|}
			\hline
			9 & $\forall N\ge N_0$: $N^\beta+N^{\beta-1}<(N+1)^\beta-\varepsilon$ & 2,8 & $\mathsf{Eq}$ \\ \hline
			10 & Lemma & 1--9 & $\rightarrow\mathsf{I}$ \\ \hline
		\end{tabular}%
	}
\end{proof}

\begin{remark}
	\label{rem:slope-beta}
	The restriction $\beta>1$ is essential (at $\beta=1$ the inequality fails
	for all $\varepsilon>0$).  All our applications have $\beta>1$:
	$\beta=3$ in Route 1a, $\beta=c\ge 3$ in Route 2, $\beta=c>8/3$ in
	Route 3, $\beta=c>1/(1-\alpha)>1$ in Route 6.
\end{remark}

\begin{lemma}[Power bound]
	\label{lem:power-bound}
	\textbf{\textsf{Assume:}}
	\begin{equation}
	x\ge 0 \wedge 0<a<1.
	\end{equation}
	\textbf{\textsf{Conclude:}}
	\begin{asser}
		\label{asser:lem-power-bound}
		(1+x)^a\le 1+ax
		\wedge
		\exists C_a\ge 1,\ \forall y\ge 0,\ (1+y)^a\le 1+C_a y
	\end{asser}
\end{lemma}

\begin{proof}[\textbf{Proof of \cref{lem:power-bound}}]
	\phantomsection\label{proof:lem:power-bound}
	{\scriptsize
		\begin{tabular}{|m{0.7cm}|m{12cm}|m{1.4cm}|m{2.5cm}|}
			\hline
			1 & \(a\in(0,1)\), \(x\ge 0\) & --- & \(\mathsf{Hyp}\) \\ \hline
			2 & \(f(t):=1+at-(1+t)^a\) & 1 & \(\mathsf{def}\) \\
		\end{tabular}%
		\noindent \newline
		\noindent
		\begin{tabular}{|m{0.7cm}|m{12cm}|m{1.4cm}|m{2.5cm}|}
			\hline
			3 & \(f(0)=0\) & 2 & \(\mathsf{Eq}\) \\ \hline
			4 & \(f'(t)=a[1-(1+t)^{a-1}]\) & 2 & \(\mathsf{Eq}\) (differentiation) \\
		\end{tabular}%
		\noindent \newline
		\noindent
		\begin{tabular}{|m{0.7cm}|m{12cm}|m{1.4cm}|m{2.5cm}|}
			\hline
			5 & \((1+t)^{a-1}\le 1\) for \(t\ge 0\) & 1 & \(\mathsf{Eq}\) \\ \hline
			6 & \(f'(t)\ge 0\) for \(t\ge 0\) & 4,5 & \(\mathsf{Eq}\) \\
		\end{tabular}%
		\noindent \newline
		\noindent
		\begin{tabular}{|m{0.7cm}|m{12cm}|m{1.4cm}|m{2.5cm}|}
			\hline
			7 & \(f\) non-decreasing on \(\{t\in\mathbb{R}:t\ge 0\}\) & 6 & \(\mathsf{Support}\) (FTC) \\ \hline
			8 & \(f(x)\ge f(0)=0\) & 3,7 & \(\mathsf{Eq}\) \\
		\end{tabular}%
		\noindent \newline
		\noindent
		\begin{tabular}{|m{0.7cm}|m{12cm}|m{1.4cm}|m{2.5cm}|}
			\hline
			9 & \((1+x)^a\le 1+ax\) & 2,8 & \(\mathsf{Eq}\) \\ \hline
			10 & For every \(y\ge 0\), \((1+y)^a\le 1+ay\) & 9 & \(\mathsf{Eq}\) (rename \(x\) as \(y\)) \\
		\end{tabular}%
		\noindent \newline
		\noindent
		\begin{tabular}{|m{0.7cm}|m{12cm}|m{1.4cm}|m{2.5cm}|}
			\hline
			11 & \(ay\le y\) for \(y\ge 0\) & 1 & \(\mathsf{Eq}\) \\ \hline
			12 & \(1+ay\le 1+y\) for \(y\ge 0\) & 11 & \(\mathsf{Eq}\) \\
		\end{tabular}%
		\noindent \newline
		\noindent
		\begin{tabular}{|m{0.7cm}|m{12cm}|m{1.4cm}|m{2.5cm}|}
			\hline
			13 & \(C_a:=1\) works & 10,12 & \(\mathsf{Eq}\) \\ \hline
			14 & Conclude & 1--13 & \(\rightarrow\mathsf{I}\) \\ \hline
		\end{tabular}%
	}
\end{proof}

\begin{lemma}[Mills squeeze]
	\label{lem:squeeze}
	\textbf{\textsf{Assume:}}
	\begin{equation}
	\begin{aligned}
	& c\in\mathbb{R}_{>1},\qquad q:\mathbb{N}_{\geq 1}\to\mathbb{R}_{>0},\\
	& \forall n\in\mathbb{N}_{\geq 1},\qquad
	q_n^{c}<q_{n+1}<(q_n+1)^{c}-1,\\
	& u_n:=q_n^{c^{-n}},\qquad
	v_n:=(q_n+1)^{c^{-n}}
	\qquad(n\in\mathbb{N}_{\geq 1}).
	\end{aligned}
	\end{equation}
	\textbf{\textsf{Conclude:}}
	\begin{asser}
		\label{asser:lem-squeeze}
		\begin{aligned}
			&(a)\ \forall n,\ u_n<u_{n+1}\wedge v_{n+1}<v_n\wedge u_n<v_n \\
			&\wedge (b)\ \exists L>0,\ \forall n,\ u_n<L<v_n \wedge u_n\to L \wedge v_n\to L
		\end{aligned}
	\end{asser}
\end{lemma}

\begin{proof}[\textbf{Proof of \cref{lem:squeeze}}]
	\phantomsection\label{proof:lem:squeeze}
	{\scriptsize
		\begin{tabular}{|m{0.7cm}|m{12cm}|m{1.4cm}|m{2.5cm}|}
			\hline
			1 & $c>1$, $(q_n)>0$, $q_n^{c}<q_{n+1}<(q_n+1)^{c}-1$ & --- & $\mathsf{Hyp}$ \\ \hline
			2 & $u_n:=q_n^{c^{-n}}$, $v_n:=(q_n+1)^{c^{-n}}$ & 1 & $\mathsf{def}$ \\
		\end{tabular}%
		\noindent \newline
		\noindent
		\begin{tabular}{|m{0.7cm}|m{12cm}|m{1.4cm}|m{2.5cm}|}
			\hline
			3 & $u_{n+1}>(q_n^{c})^{c^{-(n+1)}}=u_n$ & 1,2 & $\mathsf{Eq}$ (monotone) \\ \hline
			4 & $v_{n+1}<((q_n+1)^{c})^{c^{-(n+1)}}=v_n$ & 1,2 & $\mathsf{Eq}$ (monotone) \\
		\end{tabular}%
		\noindent \newline
		\noindent
		\begin{tabular}{|m{0.7cm}|m{12cm}|m{1.4cm}|m{2.5cm}|}
			\hline
			5 & $u_n<v_n$ & 2 & $q_n<q_n+1$ \\ \hline
			6 & $(u_n)$ increasing and bounded by $v_1$; $(v_n)$ decreasing and bounded by $u_1$ & 3,4,5 & $\mathsf{Eq}$ \\
		\end{tabular}%
		\noindent \newline
		\noindent
		\begin{tabular}{|m{0.7cm}|m{12cm}|m{1.4cm}|m{2.5cm}|}
			\hline
			7 & $\exists L_u,L_v$: $u_n\to L_u$, $v_n\to L_v$ & 6 & \cref{lem:monotone} \\ \hline
			8 & $q_{n+1}+1\le(q_n+1)^c$ & 1 & $\mathsf{Eq}$ \\
		\end{tabular}%
		\noindent \newline
		\noindent
		\begin{tabular}{|m{0.7cm}|m{12cm}|m{1.4cm}|m{2.5cm}|}
			\hline
			9 & Induction: $q_n\le(q_1+1)^{c^{\,n-1}}$ & 8 & $\mathsf{Ind}$ \\ \hline
			10 & $q_n^{c^{-n}}\le(q_1+1)^{c^{-1}}=:C$ & 9 & $\mathsf{Eq}$ \\
		\end{tabular}%
		\noindent \newline
		\noindent
		\begin{tabular}{|m{0.7cm}|m{12cm}|m{1.4cm}|m{2.5cm}|}
			\hline
			11 & $(1+1/q_n)^{c^{-n}}-1\le c^{-n}/q_n$ & 10,1 & \cref{lem:power-bound} \\ \hline
			12 & $v_n-u_n\le C\cdot c^{-n}/q_n\le C\cdot c^{-n}$ & 2,10,11 & $\mathsf{Eq}$ \\
		\end{tabular}%
		\noindent \newline
		\noindent
		\begin{tabular}{|m{0.7cm}|m{12cm}|m{1.4cm}|m{2.5cm}|}
			\hline
			13 & $v_n-u_n\to 0$ & 12 & $\mathsf{PowRate}$ \\ \hline
			14 & $L_u=L_v=:L$ & 7,13 & $\mathsf{Eq}$ \\
		\end{tabular}%
		\noindent \newline
		\noindent
		\begin{tabular}{|m{0.7cm}|m{12cm}|m{1.4cm}|m{2.5cm}|}
			\hline
			15 & $u_n<L<v_n$ (strict monotonicity) & 3,4,7,14 & $\mathsf{Eq}$ \\ \hline
			16 & (a) and (b) hold & 3,4,14,15 & $\wedge\mathsf{I}$ \\
		\end{tabular}%
		\noindent \newline
		\noindent
		\begin{tabular}{|m{0.7cm}|m{12cm}|m{1.4cm}|m{2.5cm}|}
			\hline
			17 & Lemma & 1--16 & $\rightarrow\mathsf{I}$ \\ \hline
		\end{tabular}%
	}
\end{proof}

\begin{lemma}[Monotone bounded convergence]
	\label{lem:monotone}
	\begin{equation}
	\mathsf{Assume}\!\left(
	\begin{aligned}
	& a\in\mathbb{R}^{\mathbb{N}_{\geq 1}}\ \wedge\\
	& \Bigl[
	\bigl(\forall n\in\mathbb{N}_{\geq 1},\ a_n\le a_{n+1}\bigr)
	\ \vee\
	\bigl(\forall n\in\mathbb{N}_{\geq 1},\ a_{n+1}\le a_n\bigr)
	\Bigr]\ \wedge\\
	& \exists B\in\mathbb{R}_{\ge 0},\ 
	\forall n\in\mathbb{N}_{\geq 1},\ 
	|a_n|\le B
	\end{aligned}
	\right).
	\end{equation}
	\textbf{\textsf{Conclude:}}
	\begin{asser}
		\label{asser:lem-monotone} \exists L:\mathbb{R},\ a_n\to L
	\end{asser}
\end{lemma}

\begin{proof}[\textbf{Proof of \cref{lem:monotone}}]
	\phantomsection\label{proof:lem:monotone}
	{\scriptsize
		\begin{tabular}{|m{0.7cm}|m{12cm}|m{1.4cm}|m{2.5cm}|}
			\hline
			1 & $(a_n)$ monotone and bounded & --- & $\mathsf{Hyp}$ \\ \hline
			2 & Case (increasing) & --- & $\vee\mathsf{I}$ \\
		\end{tabular}%
		\noindent \newline
		\noindent
		\begin{tabular}{|m{0.7cm}|m{12cm}|m{1.4cm}|m{2.5cm}|}
			\hline
			3 & $S:=\{a_n\}$ non-empty, bounded above & 1,2 & $\mathsf{def}$ \\ \hline
			4 & $L:=\sup S$ exists & 3 & $\mathsf{Ax}$ (Dedekind) \\
		\end{tabular}%
		\noindent \newline
		\noindent
		\begin{tabular}{|m{0.7cm}|m{12cm}|m{1.4cm}|m{2.5cm}|}
			\hline
			5 & $\forall n$, $a_n\le L$ & 4 & $\mathsf{def}(\sup)$ \\ \hline
			6 & $\forall\varepsilon>0$, $\exists N$, $a_N>L-\varepsilon$ & 4 & $\mathsf{def}(\sup)$ \\
		\end{tabular}%
		\noindent \newline
		\noindent
		\begin{tabular}{|m{0.7cm}|m{12cm}|m{1.4cm}|m{2.5cm}|}
			\hline
			7 & $\forall n\ge N$: $|a_n-L|<\varepsilon$ & 2,5,6 & $\mathsf{Eq}$ \\ \hline
			8 & $a_n\to L$ & 7 & $\mathsf{def}$ (limit) \\
		\end{tabular}%
		\noindent \newline
		\noindent
		\begin{tabular}{|m{0.7cm}|m{12cm}|m{1.4cm}|m{2.5cm}|}
			\hline
			9 & Case (decreasing): symmetric, with $L:=\inf S$ & 1 & $\mathsf{Ref}$ \\ \hline
			10 & $(a_n)$ converges & 2--8,9 & $\vee\mathsf{E}$ \\
		\end{tabular}%
		\noindent \newline
		\noindent
		\begin{tabular}{|m{0.7cm}|m{12cm}|m{1.4cm}|m{2.5cm}|}
			\hline
			11 & Lemma & 1--10 & $\rightarrow\mathsf{I}$ \\ \hline
		\end{tabular}%
	}
\end{proof}

\begin{lemma}[Wright's nested-interval facts]
	\label{lem:wright-facts}
	{ \scriptsize
		\begin{equation}
		\mathsf{Assume}\!\left(
		\begin{aligned}
		& \Sigma=\bigcap_{n\in\mathbb{N}_{\geq 1}}F_n
		\ \wedge\
		\mathsf{Nested}(F)
		\ \wedge\
		\bigl(\forall n\in\mathbb{N}_{\geq 1},\ F_n\neq\varnothing\bigr)
		\ \wedge\
		\bigl(\forall n\in\mathbb{N}_{\geq 1},\ \mathsf{FiniteUnionClosed}(F_n)\bigr)\\
		& \wedge\ \text{(iii)}\ \bigl(\forall\varepsilon\in\mathbb{R}_{>0},\ \exists N\in\mathbb{N}_{\geq 1},\ \forall n\in\mathbb{N}_{\geq N},\ \mu(F_n)<\varepsilon\bigr)
		\ \wedge\ \text{(iv$'$)}\ \mathsf{EventualBinaryBranching}(F)
		\end{aligned}
		\right),
		\end{equation}
		\begin{equation}
		\begin{aligned}
		\mathsf{Nested}(F)
		&:\iff
		\forall n\in\mathbb{N}_{\geq 1},\ F_{n+1}\subseteq F_n,\\[2pt]
		\mathsf{ClosedInterval}(J)
		&:\iff
		\exists a,b\in\mathbb{R},\ a\le b,\ J=[a,b],\\[2pt]
		\mathsf{FiniteUnionClosed}(E)
		&:\iff
		\exists m\in\mathbb{N}_{\geq 1},\ \exists J_1,\dots,J_m,\ 
		\Bigl(\forall i\in\{1,\dots,m\},\ \mathsf{ClosedInterval}(J_i)\Bigr)
		\wedge E=\bigcup_{i=1}^{m}J_i,\\[2pt]
		\mathsf{Comp}(J,E)
		&:\iff
		\mathsf{ClosedInterval}(J)\wedge J\subseteq E
		\wedge\forall J'\subseteq E,\ 
		\bigl(\mathsf{ClosedInterval}(J')\wedge J\subseteq J'\bigr)
		\Rightarrow J'=J,\\[2pt]
		\mathsf{EventualBinaryBranching}(F)
		&:\iff
		\exists N_0\in\mathbb{N}_{\geq 1},\ \forall n\in\mathbb{N}_{\geq N_0},\ 
		\forall J\subseteq\mathbb{R},\ 
		\mathsf{Comp}(J,F_n)\Rightarrow{}\\
		&\qquad\exists J_0,J_1\subseteq\mathbb{R},\ 
		\mathsf{Comp}(J_0,F_{n+1})\wedge\mathsf{Comp}(J_1,F_{n+1})
		\wedge J_0\cap J_1=\varnothing
		\wedge J_0\cup J_1\subseteq J.
		\end{aligned}
		\end{equation}
	}
	\textbf{\textsf{Conclude:}}
	\begin{asser}
		\label{asser:lem-wright-facts}
		\begin{aligned}
			&(a)\ \Sigma\ne\emptyset \wedge (b)\ \mu(\Sigma)=0 \\
			&\wedge (c)\ \mathrm{NowhereDense}(\Sigma) \wedge (d)\ |\Sigma|=|\mathbb{R}|
		\end{aligned}
	\end{asser}
\end{lemma}

\begin{proof}[\textbf{Proof of \cref{lem:wright-facts}}]
	\phantomsection\label{proof:lem:wright-facts}
	{\scriptsize
		\begin{tabular}{|m{0.7cm}|m{12cm}|m{1.4cm}|m{2.5cm}|}
			\hline
			1 & \(\Sigma=\bigcap_{n\ge1}F_n\), nested, \(F_n\ne\emptyset\), \(F_n\) finite union of closed intervals, \(\mu(F_n)\to0\), and (iv$'$) with \(N_0\) & --- & \(\mathsf{Hyp}\) \\ \hline
			2 & There is \(N_1\ge1\) such that \(\mu(F_n)<1\) for all \(n\ge N_1\) & 1 & \(\mathsf{Eq}\) (\(\mu(F_n)\to0\)) \\
		\end{tabular}%
		\noindent \newline
		\noindent
		\begin{tabular}{|m{0.7cm}|m{12cm}|m{1.4cm}|m{2.5cm}|}
			\hline
			3 & For \(n\ge N_1\), \(F_n\) is bounded; hence \(F_n\) is compact & 1,2 & \(\mathsf{Support}\) (closed finite union of bounded intervals) \\ \hline
			4 & The tail family \((F_n)_{n\ge N_1}\) has the finite intersection property & 1,3 & \(\mathsf{Eq}\) (nested; finite intersection equals \(F_{\max n}\)) \\
		\end{tabular}%
		\noindent \newline
		\noindent
		\begin{tabular}{|m{0.7cm}|m{12cm}|m{1.4cm}|m{2.5cm}|}
			\hline
			5 & \(\Sigma=\bigcap_{n\ge N_1}F_n\ne\emptyset\) & 3,4 & \(\mathsf{Support}\) (compact FIP) \\ \hline
			6 & \(\Sigma\subseteq F_n\) for every \(n\) & 1 & \(\mathsf{def}\) \\
		\end{tabular}%
		\noindent \newline
		\noindent
		\begin{tabular}{|m{0.7cm}|m{12cm}|m{1.4cm}|m{2.5cm}|}
			\hline
			7 & \(\mu(\Sigma)\le\mu(F_n)\) for every \(n\) & 6 & \(\mathsf{Support}\) (monotonicity of \(\mu\)) \\ \hline
			8 & For every \(\varepsilon>0\) there is \(n\) with \(\mu(F_n)<\varepsilon\) & 1 & \(\mathsf{Eq}\) (\(\mu(F_n)\to0\)) \\
		\end{tabular}%
		\noindent \newline
		\noindent
		\begin{tabular}{|m{0.7cm}|m{12cm}|m{1.4cm}|m{2.5cm}|}
			\hline
			9 & For every \(\varepsilon>0\), \(\mu(\Sigma)<\varepsilon\) & 7,8 & \(\mathsf{Eq}\) \\ \hline
			10 & \(\mu(\Sigma)=0\) & 9 & \(\mathsf{Eq}\) \\
		\end{tabular}%
		\noindent \newline
		\noindent
		\begin{tabular}{|m{0.7cm}|m{12cm}|m{1.4cm}|m{2.5cm}|}
			\hline
			11 & \(\Sigma\) is closed & 1 & \(\mathsf{Support}\) (intersection of closed sets) \\ \hline
			12 & \(\mathrm{int}(\Sigma)=\emptyset\) & 10,11 & \(\mathsf{Support}\) (a nonempty open interval has positive measure) \\
		\end{tabular}%
		\noindent \newline
		\noindent
		\begin{tabular}{|m{0.7cm}|m{12cm}|m{1.4cm}|m{2.5cm}|}
			\hline
			13 & \(\mathrm{NowhereDense}(\Sigma)\) & 11,12 & \(\mathsf{Eq}\) (closed with empty interior) \\ \hline
			14 & For every \(n\ge N_0\), every component interval \(J\) of \(F_n\) contains at least two component intervals of \(F_{n+1}\) & 1 & \(\mathsf{Support}\) (subintervals of \(F_{n+1}\) lie in components) \\
		\end{tabular}%
		\noindent \newline
		\noindent
		\begin{tabular}{|m{0.7cm}|m{12cm}|m{1.4cm}|m{2.5cm}|}
			\hline
			15 & Let \(T\) be the tree of component intervals of \(F_n\) for \(n\ge N_0\), ordered by inclusion & 14 & \(\mathsf{def}\) \\ \hline
			16 & Every node of \(T\) has at least two children & 14,15 & \(\mathsf{Eq}\) \\
		\end{tabular}%
		\noindent \newline
		\noindent
		\begin{tabular}{|m{0.7cm}|m{12cm}|m{1.4cm}|m{2.5cm}|}
			\hline
			17 & For every finite binary word \(s\), choose nested component intervals \(I_s\) of \(F_{N_0+|s|}\), with \(I_{s0},I_{s1}\) disjoint children of \(I_s\) & 16 & \(\mathsf{Support}\) (recursion) \\ \hline
			18 & Every component of \(F_n\) has length at most \(\mu(F_n)\), hence lengths tend to \(0\) & 1 & \(\mathsf{Eq}\) \\
		\end{tabular}%
		\noindent \newline
		\noindent
		\begin{tabular}{|m{0.7cm}|m{12cm}|m{1.4cm}|m{2.5cm}|}
			\hline
			19 & For every \(\sigma\in\{0,1\}^{\mathbb{N}}\), \(\bigcap_{k\ge0}I_{\sigma|k}\) is a singleton contained in \(\Sigma\) & 17,18 & \(\mathsf{Support}\) (nested closed intervals with lengths tending to \(0\)) \\ \hline
			20 & The map \(\sigma\mapsto \bigcap_{k\ge0}I_{\sigma|k}\) is injective & 17,19 & \(\mathsf{Eq}\) (disjoint choices at first difference) \\ 
		\end{tabular}%
		\noindent \newline
		\noindent
		\begin{tabular}{|m{0.7cm}|m{12cm}|m{1.4cm}|m{2.5cm}|}
			\hline
			21 & \(|\Sigma|\ge|\{0,1\}^{\mathbb{N}}|=|\mathbb{R}|\) & 20 & \(\mathsf{Eq}\) \\ \hline
			22 & \(|\Sigma|\le|\mathbb{R}|\) & 1 & \(\mathsf{Eq}\) (\(\Sigma\subseteq\mathbb{R}\)) \\
		\end{tabular}%
		\noindent \newline
		\noindent
		\begin{tabular}{|m{0.7cm}|m{12cm}|m{1.4cm}|m{2.5cm}|}
			\hline
			23 & \(|\Sigma|=|\mathbb{R}|\) & 21,22 & \(\mathsf{Eq}\) \\ \hline
			24 & Conclude (a)--(d) & 5,10,13,23 & \(\rightarrow\mathsf{I}\) \\ \hline
		\end{tabular}%
	}
\end{proof}

\begin{lemma}[Wright convergence]
	\label{lem:wright-conv}
	\textbf{\textsf{Assume:}}
	\begin{equation}
	q_1\ge 4,\qquad
	\forall n\in\mathbb{N}_{\geq 1},\ 2^{q_n}<q_{n+1}<2^{q_n+1},
	\end{equation}
	and $u_n:=\log_2^{[n]}(q_n)$, $v_n:=\log_2^{[n]}(q_n+1)$.
	\textbf{\textsf{Conclude:}}
	\begin{asser}
		\label{asser:lem-wright-conv}
		\begin{aligned}
			&(a)\ \forall n,\ u_n<u_{n+1}<v_{n+1}<v_n \\
			&\wedge (b)\ \exists\theta\in(u_1,v_1),\ u_n\to\theta,\ v_n\to\theta,\ u_n<\theta<v_n.
		\end{aligned}
	\end{asser}
\end{lemma}

\begin{proof}[\textbf{Proof of \cref{lem:wright-conv}}]
	\phantomsection\label{proof:lem:wright-conv}
	{\scriptsize
		\begin{tabular}{|m{0.7cm}|m{12cm}|m{1.4cm}|m{2.5cm}|}
			\hline
			1 & $q_1\ge 4$, $2^{q_n}<q_{n+1}<2^{q_n+1}$ for all $n$ & --- & $\mathsf{Hyp}$ \\ \hline
			2 & $q_n\ge 4$ for all $n$ (induction from $q_1\ge 4$) & 1 & $\mathsf{Ind}$ \\
		\end{tabular}%
		\noindent \newline
		\noindent
		\begin{tabular}{|m{0.7cm}|m{12cm}|m{1.4cm}|m{2.5cm}|}
			\hline
			3 & $q_n<q_{n+1}$ (since $2^{q_n}>q_n$ for $q_n\ge 4$) & 1,2 & $\mathsf{Eq}$ \\ \hline
			4 & $q_{n+1}+1\le 2^{q_n+1}$ (since $q_{n+1}<2^{q_n+1}$ in $\mathbb{N}$) & 1 & $\mathsf{Eq}$ \\
		\end{tabular}%
		\noindent \newline
		\noindent
		\begin{tabular}{|m{0.7cm}|m{12cm}|m{1.4cm}|m{2.5cm}|}
			\hline
			5 & $\log_2:\mathbb{R}_{>0}\to\mathbb{R}$ strictly increasing & --- & $\mathsf{Support}$ (log monotone) \\ \hline
			6 & $\log_2^{[n]}$ strictly increasing on $\mathbb{R}_{>0}$ for every $n$ & 5 & $\mathsf{Ind}$ \\
		\end{tabular}%
		\noindent \newline
		\noindent
		\begin{tabular}{|m{0.7cm}|m{12cm}|m{1.4cm}|m{2.5cm}|}
			\hline
			7 & $\log_2(q_{n+1})>q_n$ and $\log_2(q_{n+1}+1)\le q_n+1$ & 1,3,4 & $\mathsf{Eq}$ \\ \hline
			8 & $q_n,q_{n+1}+1\in\mathbb{R}_{>0}$ and $q_n<q_{n+1}+1$ & 1,3 & $\mathsf{Eq}$ \\
		\end{tabular}%
		\noindent \newline
		\noindent
		\begin{tabular}{|m{0.7cm}|m{12cm}|m{1.4cm}|m{2.5cm}|}
			\hline
			9 & $u_n:=\log_2^{[n]}(q_n)$, $v_n:=\log_2^{[n]}(q_n+1)$ & 6,8 & $\mathsf{Def}$ \\ \hline
			10 & $\forall n$: $u_n<u_{n+1}$ & 6,9 & $\mathsf{LogInv}$ (monotone) \\
		\end{tabular}%
		\noindent \newline
		\noindent
		\begin{tabular}{|m{0.7cm}|m{12cm}|m{1.4cm}|m{2.5cm}|}
			\hline
			11 & $\forall n$: $v_{n+1}<v_n$ & 8,9 & $\mathsf{LogInv}$ (monotone) \\ \hline
			12 & $\forall n$: $u_n<v_n$ & 9 & $q_n<q_n+1$ \\
		\end{tabular}%
		\noindent \newline
		\noindent
		\begin{tabular}{|m{0.7cm}|m{12cm}|m{1.4cm}|m{2.5cm}|}
			\hline
			13 & $\forall n$: $u_n<u_{n+1}<v_{n+1}<v_n$ & 10,11,12 & $\wedge\mathsf{I}$ \\ \hline
			14 & $\exists\theta$, $\theta=\lim u_n=\lim v_n$, $u_n<\theta<v_n$ & 13 & $\mathsf{Ref}$ (squeeze) \\
		\end{tabular}%
		\noindent \newline
		\noindent
		\begin{tabular}{|m{0.7cm}|m{12cm}|m{1.4cm}|m{2.5cm}|}
			\hline
			15 & (a) and (b) hold, with $\theta\in(u_1,v_1)$ & 12,13,14 & $\wedge\mathsf{I}$ \\ \hline
			16 & Lemma & 1--15 & $\rightarrow\mathsf{I}$ \\ \hline
		\end{tabular}%
	}
\end{proof}

\begin{lemma}[Polynomial-gap construction]
	\label{lem:polygap-construct}
	\textbf{\textsf{Assume:}} $\alpha\in(0,1)$, $k>0$, $\forall n,\ p_{n+1}-p_n<k\,p_n^{\alpha}$ (\cref{asser:polygap}); $c>1/(1-\alpha)$.
	\textbf{\textsf{Conclude:}}
	\begin{asser}
		\label{asser:lem-polygap-construct}
		$\exists N_0,\ \forall N\ge N_0,\ \exists p$ prime with $N^c<p<(N+1)^c-1$.
	\end{asser}
\end{lemma}

\begin{proof}[\textbf{Proof of \cref{lem:polygap-construct}}]
	\phantomsection\label{proof:lem:polygap-construct}
	{\scriptsize
		\begin{tabular}{|m{0.7cm}|m{12cm}|m{1.4cm}|m{2.5cm}|}
			\hline
			1 & $\alpha\in(0,1)$, $k>0$, $p_{n+1}-p_n<kp_n^\alpha$, $c>1/(1-\alpha)$ & --- & $\mathsf{Hyp}$ \\ \hline
			2 & $\delta:=c(1-\alpha)-1>0$ & 1 & $\mathsf{Eq}$ \\
		\end{tabular}%
		\noindent \newline
		\noindent
		\begin{tabular}{|m{0.7cm}|m{12cm}|m{1.4cm}|m{2.5cm}|}
			\hline
			3 & $N_0:=k^{1/\delta}+1$ & 2 & $\mathsf{def}$ \\ \hline
			4 & $N\ge N_0$, $p_n$ largest prime $<N^c$ & 3 & $\exists\mathsf{E}$ \\
		\end{tabular}%
		\noindent \newline
		\noindent
		\begin{tabular}{|m{0.7cm}|m{12cm}|m{1.4cm}|m{2.5cm}|}
			\hline
			5 & $p_n<N^c<p_{n+1}<p_n+k\,p_n^\alpha$ & 1,4 & $\forall\mathsf{E}$ \\ \hline
			6 & $k\,p_n^\alpha<k\,N^{c\alpha}<N^{\delta}N^{c\alpha}=N^{c-1}$ & 3,4,5 & $\mathsf{Eq}$ \\
		\end{tabular}%
		\noindent \newline
		\noindent
		\begin{tabular}{|m{0.7cm}|m{12cm}|m{1.4cm}|m{2.5cm}|}
			\hline
			7 & $p_{n+1}<N^c+N^{c-1}$ & 5,6 & $\mathsf{Eq}$ \\ \hline
			8 & $N^c+N^{c-1}<(N+1)^c-1$ & 4 & \cref{lem:slope} ($\beta=c>1$, $\varepsilon=1$) \\
		\end{tabular}%
		\noindent \newline
		\noindent
		\begin{tabular}{|m{0.7cm}|m{12cm}|m{1.4cm}|m{2.5cm}|}
			\hline
			9 & $\exists p$ prime in $(N^c,(N+1)^c-1)$ & 4,7,8 & $\exists\mathsf{I}$ \\ \hline
			10 & Conclude & 1--9 & $\rightarrow\mathsf{I}$ \\ \hline
		\end{tabular}%
	}
\end{proof}

% =========================================================================
\subsection{Route 1a: Proof of \cref{thm:mills}}
% =========================================================================

\begin{proof}[\textbf{Proof of \cref{thm:mills}}]
	\phantomsection\label{proof:mills}
	{\scriptsize
		\begin{tabular}{|m{0.7cm}|m{12cm}|m{1.4cm}|m{2.5cm}|}
			\hline
			1 & $\text{Ingham}$ & --- & $\mathsf{Hyp}$ \\ \hline
			2 & $\exists k>0,\ \forall n,\ p_{n+1}-p_n<k\,p_n^{5/8}$ & 1 & $\mathsf{def}$ \\
		\end{tabular}%
		\noindent \newline
		\noindent
		\begin{tabular}{|m{0.7cm}|m{12cm}|m{1.4cm}|m{2.5cm}|}
			\hline
			3 & $k$ as in 2 & 2 & $\exists\mathsf{E}$ \\ \hline
			4 & $N>k^{8}$, $p_n$ largest prime $<N^{3}$ (\cref{lem:euclid}) & 3 & $\mathsf{def}$ \\
		\end{tabular}%
		\noindent \newline
		\noindent
		\begin{tabular}{|m{0.7cm}|m{12cm}|m{1.4cm}|m{2.5cm}|}
			\hline
			5 & $p_n<N^{3}<p_{n+1}$ & 4 & $\mathsf{Def}$ \\ \hline
			6 & $p_{n+1}<p_n+k\,p_n^{5/8}$ & 3 & $\forall\mathsf{E}$ \\
		\end{tabular}%
		\noindent \newline
		\noindent
		\begin{tabular}{|m{0.7cm}|m{12cm}|m{1.4cm}|m{2.5cm}|}
			\hline
			7 & $k\,p_n^{5/8}<N^{1/8}N^{15/8}=N^{2}$ & 5,4 & $\mathsf{Eq}$ \\ \hline
			8 & $N^{3}<p_{n+1}<N^{3}+N^{2}$ & 5,6,7 & $\mathsf{Eq}$ \\
		\end{tabular}%
		\noindent \newline
		\noindent
		\begin{tabular}{|m{0.7cm}|m{12cm}|m{1.4cm}|m{2.5cm}|}
			\hline
			9 & $N^{3}+N^{2}<(N+1)^{3}-1$ & 4 & \cref{lem:slope} ($\beta=3$, $\varepsilon=1$) \\ \hline
			10 & $\exists$ prime in $(N^{3},(N+1)^{3}-1)$ & 5,8,9 & $\exists\mathsf{I}$ \\
		\end{tabular}%
		\noindent \newline
		\noindent
		\begin{tabular}{|m{0.7cm}|m{12cm}|m{1.4cm}|m{2.5cm}|}
			\hline
			11 & $\forall N>k^{8}$: same & 4--10 & $\forall\mathsf{I}$ \\ \hline
			12 & $q_1$ smallest prime $>k^{8}$; $q_{n+1}:=\min\{p:\ q_n^{3}<p<(q_n+1)^{3}-1\}$ & 11 & $\exists\mathsf{E}$ \\
		\end{tabular}%
		\noindent \newline
		\noindent
		\begin{tabular}{|m{0.7cm}|m{12cm}|m{1.4cm}|m{2.5cm}|}
			\hline
			13 & $\mathrm{MillsRec}(q)$ & 12 & $\mathsf{Def}$ \\ \hline
			14 & $u_n:=q_n^{3^{-n}}$, $v_n:=(q_n+1)^{3^{-n}}$ & 12 & $\exists\mathsf{I}$ \\
		\end{tabular}%
		\noindent \newline
		\noindent
		\begin{tabular}{|m{0.7cm}|m{12cm}|m{1.4cm}|m{2.5cm}|}
			\hline
			15 & $u_n<u_{n+1}<v_{n+1}<v_n$ & 13,14 & \cref{thm:millsseq} \\ \hline
			16 & $\exists\theta$, $\theta=\lim u_n=\lim v_n$, $u_n<\theta<v_n$ & 15 & \cref{lem:squeeze} \\
		\end{tabular}%
		\noindent \newline
		\noindent
		\begin{tabular}{|m{0.7cm}|m{12cm}|m{1.4cm}|m{2.5cm}|}
			\hline
			17 & $\forall n$: $q_n<\theta^{3^{n}}<q_n+1$ & 14,16 & $\mathsf{Eq}$ \\ \hline
			18 & $\forall n$: $\mathrm{floor}(\theta^{3^{n}})=q_n$ & 17 & $\mathsf{Def}$ \\
		\end{tabular}%
		\noindent \newline
		\noindent
		\begin{tabular}{|m{0.7cm}|m{12cm}|m{1.4cm}|m{2.5cm}|}
			\hline
			19 & $\forall n$: $\mathrm{Prime}(\mathrm{floor}(\theta^{3^{n}}))$ & 12,18 & $\mathsf{Def}$ \\ \hline
			20 & $\theta>1$ & 12,16 & $\mathsf{Eq}$ \\
		\end{tabular}%
		\noindent \newline
		\noindent
		\begin{tabular}{|m{0.7cm}|m{12cm}|m{1.4cm}|m{2.5cm}|}
			\hline
			21 & $\exists\theta>1,\ \mathrm{IsMills}(\theta)$ & 19,20 & $\exists\mathsf{I}$ \\ \hline
			22 & $\text{Ingham}\Rightarrow\exists\theta>1,\ \mathrm{IsMills}(\theta)$ & 1--21 & $\rightarrow\mathsf{I}$ \\ \hline
		\end{tabular}%
	}
\end{proof}

% =========================================================================
\subsection{Route 1b: Proof of \cref{thm:millsseq}}
% =========================================================================

\begin{proof}[\textbf{Proof of \cref{thm:millsseq}}]
	\phantomsection\label{proof:millsseq}
	{\scriptsize
		\begin{tabular}{|m{0.7cm}|m{12cm}|m{1.4cm}|m{2.5cm}|}
			\hline
			1 & $\mathrm{UV}(q,u,v)\wedge\mathrm{MillsRec}(q)$ & --- & $\mathsf{Hyp}$ \\ \hline
			2 & $\mathrm{UV}(q,u,v)$ & 1 & $\wedge\mathsf{E}$ \\
		\end{tabular}%
		\noindent \newline
		\noindent
		\begin{tabular}{|m{0.7cm}|m{12cm}|m{1.4cm}|m{2.5cm}|}
			\hline
			3 & $\mathrm{MillsRec}(q)$ & 1 & $\wedge\mathsf{E}$ \\ \hline
			4 & $\forall n$: $u_n=q_n^{3^{-n}}$, $v_n=(q_n+1)^{3^{-n}}$ & 2 & $\mathsf{Def}$ \\
		\end{tabular}%
		\noindent \newline
		\noindent
		\begin{tabular}{|m{0.7cm}|m{12cm}|m{1.4cm}|m{2.5cm}|}
			\hline
			5 & $\forall n$: $q_n^{3}<q_{n+1}<(q_n+1)^{3}-1$ & 3 & $\mathsf{Def}$ \\ \hline
			6 & $\forall n$: $q_{n+1}>q_n^{3}\ge q_n$ & 5 & $\mathsf{Eq}$ \\
		\end{tabular}%
		\noindent \newline
		\noindent
		\begin{tabular}{|m{0.7cm}|m{12cm}|m{1.4cm}|m{2.5cm}|}
			\hline
			7 & $\forall n$: $u_{n+1}>(q_n^{3})^{3^{-(n+1)}}=u_n$ & 4,6 & $\mathsf{Eq}$ \\ \hline
			8 & $\forall n$: $v_{n+1}<((q_n+1)^{3})^{3^{-(n+1)}}=v_n$ & 4,5 & $\mathsf{Eq}$ \\
		\end{tabular}%
		\noindent \newline
		\noindent
		\begin{tabular}{|m{0.7cm}|m{12cm}|m{1.4cm}|m{2.5cm}|}
			\hline
			9 & $\forall n$: $u_n<v_n$ & 4 & $q_n<q_n+1$ \\ \hline
			10 & $\forall n$: $u_n<u_{n+1}<v_{n+1}<v_n$ & 7,8,9 & $\wedge\mathsf{I}$ \\
		\end{tabular}%
		\noindent \newline
		\noindent
		\begin{tabular}{|m{0.7cm}|m{12cm}|m{1.4cm}|m{2.5cm}|}
			\hline
			11 & $(\mathrm{UV}\wedge\mathrm{MillsRec})\Rightarrow\forall n,\ u_n<u_{n+1}<v_{n+1}<v_n$ & 1--10 & $\rightarrow\mathsf{I}$ \\ \hline
		\end{tabular}%
	}
\end{proof}

\begin{remark}
	{\scriptsize The label \cref{thm:millsseq} is the theorem; \cref{ineq:millsseq} is its conclusion.  The Dudley transcription writes $5^{-n}$ instead of $3^{-n}$, which violates the strict monotonicity proved in rows 7--8. }
\end{remark}

% =========================================================================
\subsection{Route 2: Proof of \cref{thm:kuipers}}
% =========================================================================

\begin{proof}[\textbf{Proof of \cref{thm:kuipers}}]
	\phantomsection\label{proof:kuipers}
	{\scriptsize
		\begin{tabular}{|m{0.7cm}|m{12cm}|m{1.4cm}|m{2.5cm}|}
			\hline
			1 & $\text{Ingham}$ & --- & $\mathsf{Hyp}$ \\ \hline
			2 & $c\in\mathbb{N}_{\geq 3}$ arbitrary & --- & $\mathsf{def}$ \\
		\end{tabular}%
		\noindent \newline
		\noindent
		\begin{tabular}{|m{0.7cm}|m{12cm}|m{1.4cm}|m{2.5cm}|}
			\hline
			3 & $a:=3c-4$, $b:=3c-1$; $8a-5b=9c-27\ge 0\Rightarrow a/b\ge 5/8$ & 2 & $\mathsf{Eq}$ \\ \hline
			4 & $ca+1=(c-1)b$ & 3 & $\mathsf{Eq}$ \\
		\end{tabular}%
		\noindent \newline
		\noindent
		\begin{tabular}{|m{0.7cm}|m{12cm}|m{1.4cm}|m{2.5cm}|}
			\hline
			5 & $\exists k_1>0$, $\forall n$, $p_{n+1}-p_n<k_1 p_n^{a/b}$ & 1,3 & $\mathsf{def}$ \\ \hline
			6 & $k_1$ as in 5 & 5 & $\exists\mathsf{E}$ \\
		\end{tabular}%
		\noindent \newline
		\noindent
		\begin{tabular}{|m{0.7cm}|m{12cm}|m{1.4cm}|m{2.5cm}|}
			\hline
			7 & $N>k_1^{b}$, $p_n$ largest prime $<N^{c}$ & 6 & $\mathsf{def}$ \\ \hline
			8 & $p_n<N^{c}<p_{n+1}$ & 7 & $\mathsf{Def}$ \\
		\end{tabular}%
		\noindent \newline
		\noindent
		\begin{tabular}{|m{0.7cm}|m{12cm}|m{1.4cm}|m{2.5cm}|}
			\hline
			9 & $p_{n+1}<N^{c}+k_1 N^{ca/b}$ & 6,8 & $\mathsf{Eq}$ \\ \hline
					10 & $k_1<N^{1/b}$ & 7 & $\mathsf{Eq}$ \\
	\end{tabular}%
\noindent \newline
\noindent
\begin{tabular}{|m{0.7cm}|m{12cm}|m{1.4cm}|m{2.5cm}|}
\hline
11 & $k_1 N^{ca/b}<N^{(ca+1)/b}=N^{c-1}$ & 10,4 & $\mathsf{Eq}$ \\ \hline
12 & $p_{n+1}<N^{c}+N^{c-1}$ & 9,11 & $\mathsf{Eq}$ \\
\end{tabular}%
\noindent \newline
\noindent
\begin{tabular}{|m{0.7cm}|m{12cm}|m{1.4cm}|m{2.5cm}|}
\hline
13 & $N^{c}+N^{c-1}<(N+1)^{c}-1$ & 7 & \cref{lem:slope} ($\beta=c\ge 3$, $\varepsilon=1$) \\ \hline
14 & $\exists$ prime in $(N^{c},(N+1)^{c}-1)$ & 8,12,13 & $\exists\mathsf{I}$ \\
\end{tabular}%
\noindent \newline
\noindent
\begin{tabular}{|m{0.7cm}|m{12cm}|m{1.4cm}|m{2.5cm}|}
\hline
15 & $\forall N>k_1^{b}$: same & 7--14 & $\forall\mathsf{I}$ \\ \hline
16 & $q_1>k_1^{b}$; $q_{n+1}:=\min\{p:\ q_n^{c}<p<(q_n+1)^{c}-1\}$ & 15 & $\exists\mathsf{E}$ \\
\end{tabular}%
\noindent \newline
\noindent
\begin{tabular}{|m{0.7cm}|m{12cm}|m{1.4cm}|m{2.5cm}|}
\hline
17 & $u_n:=q_n^{c^{-n}}$, $v_n:=(q_n+1)^{c^{-n}}$ & 16 & $\mathsf{Def}$ \\ \hline
18 & $u_n<u_{n+1}<v_{n+1}<v_n$ & 16,17 & \cref{lem:squeeze}a \\
\end{tabular}%
\noindent \newline
\noindent
\begin{tabular}{|m{0.7cm}|m{12cm}|m{1.4cm}|m{2.5cm}|}
\hline
19 & $\exists\theta$, $\theta=\lim u_n=\lim v_n$, $u_n<\theta<v_n$ & 18 & \cref{lem:squeeze}b \\ \hline
20 & $\forall n$: $q_n<\theta^{c^{n}}<q_n+1$ & 17,19 & $\mathsf{Eq}$ \\
\end{tabular}%
\noindent \newline
\noindent
\begin{tabular}{|m{0.7cm}|m{12cm}|m{1.4cm}|m{2.5cm}|}
\hline
21 & $\forall n$: $\mathrm{floor}(\theta^{c^{n}})=q_n$ prime & 16,20 & $\mathsf{Def}$ \\ \hline
22 & $\theta>1$ & 19,20 & $\mathsf{Eq}$ \\
\end{tabular}%
\noindent \newline
\noindent
\begin{tabular}{|m{0.7cm}|m{12cm}|m{1.4cm}|m{2.5cm}|}
\hline
23 & $\mathrm{IsKuipers}(c,\theta)$ & 21,22 & $\mathsf{Def}$ \\ \hline
24 & $\exists\theta>1,\ \mathrm{IsKuipers}(c,\theta)$ & 23 & $\exists\mathsf{I}$ \\
\end{tabular}%
\noindent \newline
\noindent
\begin{tabular}{|m{0.7cm}|m{12cm}|m{1.4cm}|m{2.5cm}|}
\hline
25 & $\forall c\ge 3$, $\exists\theta>1$, $\mathrm{IsKuipers}(c,\theta)$ & 2--24 & $\forall\mathsf{I}$ \\ \hline
26 & $\text{Ingham}\Rightarrow\forall c\ge 3$, $\exists\theta>1$, $\mathrm{IsKuipers}(c,\theta)$ & 1--25 & $\rightarrow\mathsf{I}$ \\
\end{tabular}%
}
\end{proof}

% =========================================================================
\subsection{Route 3: Proof of \cref{thm:ansari}}
% =========================================================================

\begin{proof}[\textbf{Proof of \cref{thm:ansari}}]
\phantomsection\label{proof:ansari}
{\scriptsize
\begin{tabular}{|m{0.7cm}|m{12cm}|m{1.4cm}|m{2.5cm}|}
\hline
1 & $\text{Ingham}$ & --- & $\mathsf{Hyp}$ \\ \hline
2 & $c\in\mathbb{R}_{>8/3}$ arbitrary & --- & $\mathsf{def}$ \\
\end{tabular}%
\noindent \newline
\noindent
\begin{tabular}{|m{0.7cm}|m{12cm}|m{1.4cm}|m{2.5cm}|}
\hline
3 & $\delta:=3c/8-1>0$ & 2 & $\mathsf{Eq}$ \\ \hline
4 & $\exists k>0$, $\forall n$, $p_{n+1}-p_n<k\,p_n^{5/8}$ & 1 & $\mathsf{def}$ \\
\end{tabular}%
\noindent \newline
\noindent
\begin{tabular}{|m{0.7cm}|m{12cm}|m{1.4cm}|m{2.5cm}|}
\hline
5 & $k$ as in 4 & 4 & $\exists\mathsf{E}$ \\ \hline
6 & $N>k^{1/\delta}$, $p_n$ largest prime $<N^{c}$ & 5 & $\mathsf{def}$ \\
\end{tabular}%
\noindent \newline
\noindent
\begin{tabular}{|m{0.7cm}|m{12cm}|m{1.4cm}|m{2.5cm}|}
\hline
7 & $p_n<N^{c}<p_{n+1}$ & 6 & $\mathsf{Def}$ \\ \hline
8 & $p_{n+1}<N^{c}+k N^{5c/8}$ & 5,7 & $\mathsf{Eq}$ \\
\end{tabular}%
\noindent \newline
\noindent
\begin{tabular}{|m{0.7cm}|m{12cm}|m{1.4cm}|m{2.5cm}|}
\hline
9 & $k<N^{\delta}$ & 6 & $\mathsf{Eq}$ \\ \hline
10 & $k N^{5c/8}<N^{\delta+5c/8}=N^{c-1}$ & 3,9 & $\mathsf{Eq}$ \\
\end{tabular}%
\noindent \newline
\noindent
\begin{tabular}{|m{0.7cm}|m{12cm}|m{1.4cm}|m{2.5cm}|}
\hline
11 & $p_{n+1}<N^{c}+N^{c-1}$ & 8,10 & $\mathsf{Eq}$ \\ \hline
12 & $N^{c}+N^{c-1}<(N+1)^{c}-1$ & 6 & \cref{lem:slope} ($\beta=c>8/3$, $\varepsilon=1$) \\
\end{tabular}%
\noindent \newline
\noindent
\begin{tabular}{|m{0.7cm}|m{12cm}|m{1.4cm}|m{2.5cm}|}
\hline
13 & $\exists$ prime in $(N^{c},(N+1)^{c}-1)$ & 7,11,12 & $\exists\mathsf{I}$ \\ \hline
14 & $\forall N>k^{1/\delta}$: same & 6--13 & $\forall\mathsf{I}$ \\
\end{tabular}%
\noindent \newline
\noindent
\begin{tabular}{|m{0.7cm}|m{12cm}|m{1.4cm}|m{2.5cm}|}
\hline
15 & $q_1>k^{1/\delta}$; $q_{n+1}:=\min\{p:\ q_n^{c}<p<(q_n+1)^{c}-1\}$ & 14 & $\exists\mathsf{E}$ \\ \hline
16 & $u_n:=q_n^{c^{-n}}$, $v_n:=(q_n+1)^{c^{-n}}$ & 15 & $\mathsf{Def}$ \\
\end{tabular}%
\noindent \newline
\noindent
\begin{tabular}{|m{0.7cm}|m{12cm}|m{1.4cm}|m{2.5cm}|}
\hline
17 & $u_n<u_{n+1}<v_{n+1}<v_n$ & 15,16 & \cref{lem:squeeze}a \\ \hline
18 & $\exists\theta$, $\theta=\lim u_n=\lim v_n$, $u_n<\theta<v_n$ & 17 & \cref{lem:squeeze}b \\
\end{tabular}%
\noindent \newline
\noindent
\begin{tabular}{|m{0.7cm}|m{12cm}|m{1.4cm}|m{2.5cm}|}
\hline
19 & $\forall n$: $q_n<\theta^{c^{n}}<q_n+1$ & 16,18 & $\mathsf{Eq}$ \\ \hline
20 & $\forall n$: $\mathrm{floor}(\theta^{c^{n}})=q_n$ prime & 15,19 & $\mathsf{Def}$ \\
\end{tabular}%
\noindent \newline
\noindent
\begin{tabular}{|m{0.7cm}|m{12cm}|m{1.4cm}|m{2.5cm}|}
\hline
21 & $\theta>1$ & 18,19 & $\mathsf{Eq}$ \\ \hline
22 & $\mathrm{IsAnsari}(c,\theta)$ & 20,21 & $\mathsf{Def}$ \\
\end{tabular}%
\noindent \newline
\noindent
\begin{tabular}{|m{0.7cm}|m{12cm}|m{1.4cm}|m{2.5cm}|}
\hline
23 & $\exists\theta>1,\ \mathrm{IsAnsari}(c,\theta)$ & 22 & $\exists\mathsf{I}$ \\ \hline
24 & $\forall c>8/3$, $\exists\theta>1$, $\mathrm{IsAnsari}(c,\theta)$ & 2--23 & $\forall\mathsf{I}$ \\
\end{tabular}%
\noindent \newline
\noindent
\begin{tabular}{|m{0.7cm}|m{12cm}|m{1.4cm}|m{2.5cm}|}
\hline
25 & $\text{Ingham}\Rightarrow\forall c>8/3$, $\exists\theta>1$, $\mathrm{IsAnsari}(c,\theta)$ & 1--24 & $\rightarrow\mathsf{I}$ \\ \hline
\end{tabular}%
}
\end{proof}

% =========================================================================
\subsection{Route 4: Proof of \cref{thm:wright}}
% =========================================================================

\begin{proof}[\textbf{Proof of \cref{thm:wright}}]
\phantomsection\label{proof:wright}
{\scriptsize
\begin{tabular}{|m{0.7cm}|m{12cm}|m{1.4cm}|m{2.5cm}|}
\hline
1 & $\text{Bertrand}$ & --- & $\mathsf{Hyp}$ \\ \hline
2 & $\forall N\ge2$, $\exists p$ prime, $N<p<2N$ & 1 & $\mathsf{def}$ \\
\end{tabular}%
\noindent \newline
\noindent
\begin{tabular}{|m{0.7cm}|m{12cm}|m{1.4cm}|m{2.5cm}|}
\hline
3 & $q_1:=5$ & 2 & $\mathsf{def}$ \\ \hline
4 & $q_{n+1}:=\min\{p \text{ prime}:\ 2^{q_n}<p<2^{q_n+1}\}$ (exists by 2 at $N=2^{q_n}$) & 2,3 & $\exists\mathsf{E}$ \\
\end{tabular}%
\noindent \newline
\noindent
\begin{tabular}{|m{0.7cm}|m{12cm}|m{1.4cm}|m{2.5cm}|}
\hline
5 & $u_n:=\log_2^{[n]}(q_n)$, $v_n:=\log_2^{[n]}(q_n+1)$ & 4 & $\mathsf{def}$ \\ \hline
6 & $u_n<u_{n+1}<v_{n+1}<v_n$ & 3,4,5 & \cref{lem:wright-conv}a \\
\end{tabular}%
\noindent \newline
\noindent
\begin{tabular}{|m{0.7cm}|m{12cm}|m{1.4cm}|m{2.5cm}|}
\hline
7 & $\exists\theta\in(u_1,v_1)$, $u_n\to\theta$, $v_n\to\theta$, $u_n<\theta<v_n$ & 5,6 & \cref{lem:wright-conv}b \\ \hline
8 & $\forall n$: $q_n<\exp_2^{[n]}(\theta)<q_n+1$ & 5,7 & $\mathsf{Eq}$ \\
\end{tabular}%
\noindent \newline
\noindent
\begin{tabular}{|m{0.7cm}|m{12cm}|m{1.4cm}|m{2.5cm}|}
\hline
9 & $\forall n$: $\mathrm{floor}(\exp_2^{[n]}(\theta))=q_n$ prime & 4,8 & $\mathsf{Def}$ \\ \hline
10 & $\theta>1$ & 3,7 & $\mathsf{Eq}$ \\
\end{tabular}%
\noindent \newline
\noindent
\begin{tabular}{|m{0.7cm}|m{12cm}|m{1.4cm}|m{2.5cm}|}
\hline
11 & $\exists\theta>1,\ \mathrm{IsWright}(\theta)$ & 9,10 & $\exists\mathsf{I}$ \\ \hline
12 & $\text{Bertrand}\Rightarrow\exists\theta>1,\ \mathrm{IsWright}(\theta)$ & 1--11 & $\rightarrow\mathsf{I}$ \\
\end{tabular}%
}
\end{proof}

% =========================================================================
\subsection{Route 5: Proof of \cref{thm:ore}}
% =========================================================================

\begin{proof}[\textbf{Proof of \cref{thm:ore}}]
\phantomsection\label{proof:ore}
{\scriptsize
\begin{tabular}{|m{0.7cm}|m{12cm}|m{1.4cm}|m{2.5cm}|}
\hline
1 & $f,q,\delta,\epsilon$ with $\mathrm{IsOreSel}(f)\wedge\mathrm{OreBound}(f,q,\delta,\epsilon)$ & --- & $\mathsf{Hyp}$ \\ \hline
2 & $\mathrm{Continuous}(f)\wedge\mathrm{EventuallyIncreasing}(f)$ & 1 & $\mathsf{Def}$ \\
\end{tabular}%
\noindent \newline
\noindent
\begin{tabular}{|m{0.7cm}|m{12cm}|m{1.4cm}|m{2.5cm}|}
\hline
3 & $\forall n$: $f(q_n-\delta_n)<q_{n+1}-\delta_{n+1}\wedge q_{n+1}+\epsilon_{n+1}<f(q_n+\epsilon_n)$ & 1 & $\mathsf{Def}$ \\ \hline
4 & $I_n:=[q_n-\delta_n,q_n+\epsilon_n]$ & 3 & $\mathsf{def}$ \\
\end{tabular}%
\noindent \newline
\noindent
\begin{tabular}{|m{0.7cm}|m{12cm}|m{1.4cm}|m{2.5cm}|}
\hline
5 & $I_n\subseteq\mathbb{R}$ closed, nonempty & 2,3 & $\mathsf{Eq}$ \\ \hline
6 & $f(I_n)\supseteq I_{n+1}$ (IVT on 3 using 2) & 2,3,4 & $\mathsf{Supp}$ (IVT) \\
\end{tabular}%
\noindent \newline
\noindent
\begin{tabular}{|m{0.7cm}|m{12cm}|m{1.4cm}|m{2.5cm}|}
\hline
7 & $J_n:=f^{-[n]}(I_n)$ & 6 & $\mathsf{def}$ \\ \hline
8 & $J_{n+1}\subseteq J_n$ (nested) & 6,7 & $\mathsf{Eq}$ \\
\end{tabular}%
\noindent \newline
\noindent
\begin{tabular}{|m{0.7cm}|m{12cm}|m{1.4cm}|m{2.5cm}|}
\hline
9 & Each $J_n$ closed (preimage of closed under continuous $f^{[n]}$) & 2,7 & $\mathsf{Supp}$ (continuity) \\ \hline
10 & W.l.o.g. $J_1$ bounded (else intersect with $[-M,M]$ for large $M$; conclusion unaffected) & 7 & $\mathsf{Supp}$ \\
\end{tabular}%
\noindent \newline
\noindent
\begin{tabular}{|m{0.7cm}|m{12cm}|m{1.4cm}|m{2.5cm}|}
\hline
11 & Finite intersection property: $J_{n_1}\cap\dots\cap J_{n_k}=J_{\max n_i}\ne\emptyset$ & 8,9 & $\mathsf{def}$ \\ \hline
12 & $\bigcap_n J_n\ne\emptyset$ (compactness of $J_1$) & 9,10,11 & $\mathsf{Supp}$ (compact FIP) \\
\end{tabular}%
\noindent \newline
\noindent
\begin{tabular}{|m{0.7cm}|m{12cm}|m{1.4cm}|m{2.5cm}|}
\hline
13 & $\theta\in\bigcap_n J_n$ & 12 & $\exists\mathsf{E}$ \\ \hline
14 & $\forall n$: $f^{[n]}(\theta)\in I_n$, i.e.\ $q_n-\delta_n\le f^{[n]}(\theta)\le q_n+\epsilon_n$ & 7,13 & $\mathsf{Eq}$ \\
\end{tabular}%
\noindent \newline
\noindent
\begin{tabular}{|m{0.7cm}|m{12cm}|m{1.4cm}|m{2.5cm}|}
\hline
15 & $\exists\theta>0,\ \forall n,\ q_n-\delta_n\le f^{[n]}(\theta)\le q_n+\epsilon_n$ & 13,14 & $\exists\mathsf{I}$ \\ \hline
16 & Conclude & 1--15 & $\rightarrow\mathsf{I}$ \\
\end{tabular}%
}
\end{proof}

% =========================================================================
\subsection{Route 6: Proof of \cref{thm:direct}}
% =========================================================================

\begin{proof}[\textbf{Proof of \cref{thm:direct}}]
\phantomsection\label{proof:direct}
{\scriptsize
\begin{tabular}{|m{0.7cm}|m{12cm}|m{1.4cm}|m{2.5cm}|}
\hline
1 & $\alpha\in(0,1)$, $k>0$, $\mathrm{PolyGap}(\alpha,k)$, $c>1/(1-\alpha)$ & --- & $\mathsf{Hyp}$ \\ \hline
2 & $\exists N_0,\ \forall N\ge N_0,\ \exists p$ prime, $N^c<p<(N+1)^c-1$ & 1 & \cref{lem:polygap-construct} \\
\end{tabular}%
\noindent \newline
\noindent
\begin{tabular}{|m{0.7cm}|m{12cm}|m{1.4cm}|m{2.5cm}|}
\hline
3 & $q_1>N_0$ (arbitrary); $q_{n+1}:=\min\{p \text{ prime}:\ q_n^c<p<(q_n+1)^c-1\}$ & 2 & $\exists\mathsf{E}$ \\ \hline
4 & $\mathrm{MillsRec}_c(q)$: $\forall n,\ q_n^c<q_{n+1}<(q_n+1)^c-1$ & 3 & $\mathsf{Def}$ \\
\end{tabular}%
\noindent \newline
\noindent
\begin{tabular}{|m{0.7cm}|m{12cm}|m{1.4cm}|m{2.5cm}|}
\hline
5 & $u_n:=q_n^{c^{-n}}$, $v_n:=(q_n+1)^{c^{-n}}$ & 3 & $\mathsf{def}$ \\ \hline
6 & $u_n<u_{n+1}<v_{n+1}<v_n$ & 4,5 & \cref{lem:squeeze}a \\
\end{tabular}%
\noindent \newline
\noindent
\begin{tabular}{|m{0.7cm}|m{12cm}|m{1.4cm}|m{2.5cm}|}
\hline
7 & $\exists\theta>0$, $u_n\to\theta$, $v_n\to\theta$, $u_n<\theta<v_n$ & 5,6 & \cref{lem:squeeze}b \\ \hline
8 & $\forall n$: $q_n<\theta^{c^n}<q_n+1$ & 5,7 & $\mathsf{Eq}$ \\
\end{tabular}%
\noindent \newline
\noindent
\begin{tabular}{|m{0.7cm}|m{12cm}|m{1.4cm}|m{2.5cm}|}
\hline
9 & $\forall n$: $\mathrm{Prime}(\mathrm{floor}(\theta^{c^n}))$ & 3,8 & $\mathsf{Def}$ \\ \hline
10 & $\exists\theta>0,\ \forall n,\ \mathrm{Prime}(\mathrm{floor}(\theta^{c^n}))$ & 7,9 & $\exists\mathsf{I}$ \\
\end{tabular}%
\noindent \newline
\noindent
\begin{tabular}{|m{0.7cm}|m{12cm}|m{1.4cm}|m{2.5cm}|}
\hline
11 & Conclude & 1--10 & $\rightarrow\mathsf{I}$ \\ \hline
\end{tabular}%
}
\end{proof}

% =========================================================================
\subsection{Route 7: Proof of \cref{thm:niven}}
% =========================================================================

\begin{proof}[\textbf{Proof of \cref{thm:niven}}]
\phantomsection\label{proof:niven}
{\scriptsize
\begin{tabular}{|m{0.7cm}|m{12cm}|m{1.4cm}|m{2.5cm}|}
\hline
1 & $\text{Ingham}$ & --- & $\mathsf{Hyp}$ \\ \hline
2 & $c\in\mathbb{R}_{>1}$ arbitrary & --- & $\mathsf{def}$ \\
\end{tabular}%
\noindent \newline
\noindent
\begin{tabular}{|m{0.7cm}|m{12cm}|m{1.4cm}|m{2.5cm}|}
\hline
3 & $\exists k>0$, $\forall n$, $p_{n+1}-p_n<k\,p_n^{5/8}$ & 1 & $\mathsf{def}$ \\ \hline
4 & $k$ as in 3; $N_0$ such that $\forall x\ge N_0$, $k\,x^{5/8}<c^{x}(c-1)/2$ & 2,3 & $\mathsf{def}$ \\
\end{tabular}%
\noindent \newline
\noindent
\begin{tabular}{|m{0.7cm}|m{12cm}|m{1.4cm}|m{2.5cm}|}
\hline
5 & $q_1>N_0$; $q_{n+1}:=\min\{p \text{ prime}:\ c^{q_n}<p<c^{q_n+1}\}$ (exists by 3,4 for $q_n\ge N_0$) & 3,4 & $\exists\mathsf{E}$ \\ \hline
6 & $u_n:=\log_c^{[n]}(q_n)$, $v_n:=\log_c^{[n]}(q_n+1)$ & 5 & $\mathsf{def}$ \\
\end{tabular}%
\noindent \newline
\noindent
\begin{tabular}{|m{0.7cm}|m{12cm}|m{1.4cm}|m{2.5cm}|}
\hline
7 & $u_n<u_{n+1}<v_{n+1}<v_n$ (Wright argument, base $c$) & 5,6 & \cref{lem:wright-conv}a ($c\ge2$) \\ \hline
8 & $\exists\theta\in(u_1,v_1)$, $u_n\to\theta$, $v_n\to\theta$, $u_n<\theta<v_n$ & 6,7 & \cref{lem:wright-conv}b \\
\end{tabular}%
\noindent \newline
\noindent
\begin{tabular}{|m{0.7cm}|m{12cm}|m{1.4cm}|m{2.5cm}|}
\hline
9 & $\forall n$: $q_n<c^{\theta^n}<q_n+1$ & 5,6,8 & $\mathsf{Eq}$ \\ \hline
10 & $\forall n$: $\mathrm{floor}(c^{\theta^n})=q_n$ prime & 5,9 & $\mathsf{Def}$ \\
\end{tabular}%
\noindent \newline
\noindent
\begin{tabular}{|m{0.7cm}|m{12cm}|m{1.4cm}|m{2.5cm}|}
\hline
11 & $\theta>1$ & 5,8 & $\mathsf{Eq}$ \\ \hline
12 & $\mathrm{IsNiven}(c,\theta)$ & 10,11 & $\mathsf{Def}$ \\
\end{tabular}%
\noindent \newline
\noindent
\begin{tabular}{|m{0.7cm}|m{12cm}|m{1.4cm}|m{2.5cm}|}
\hline
13 & $\exists\theta>0,\ \mathrm{IsNiven}(c,\theta)$ & 12 & $\exists\mathsf{I}$ \\ \hline
14 & $\forall c>1,\ \exists\theta>0,\ \mathrm{IsNiven}(c,\theta)$ & 2--13 & $\forall\mathsf{I}$ \\
\end{tabular}%
\noindent \newline
\noindent
\begin{tabular}{|m{0.7cm}|m{12cm}|m{1.4cm}|m{2.5cm}|}
\hline
15 & Conclude & 1--14 & $\rightarrow\mathsf{I}$ \\ \hline
\end{tabular}%
}
\end{proof}

% =========================================================================
\subsection{Route 8: Proof of \cref{thm:wright-metrical}}
% =========================================================================

\begin{proof}[\textbf{Proof of \cref{thm:wright-metrical}}]
\phantomsection\label{proof:wright-metrical}
{\scriptsize
\begin{tabular}{|m{0.7cm}|m{12cm}|m{1.4cm}|m{2.5cm}|}
\hline
1 & Each route yields, at every stage $n$, a nonempty admissible interval for $q_{n+1}$ & --- & $\mathsf{Hyp}$ \\ \hline
2 & For each of Ansari, Wright, Niven: $\Theta:=\{\theta:\ \text{construction succeeds}\}$ satisfies (i)--(iv$'$) of \cref{lem:wright-facts} & 1 & $\mathsf{Supp}$ (per-route check) \\
\end{tabular}%
\noindent \newline
\noindent
\begin{tabular}{|m{0.7cm}|m{12cm}|m{1.4cm}|m{2.5cm}|}
\hline
3 & $\mathrm{IsNull}(\mathsf{ThetaAnsari})\wedge\mathrm{NowhereDense}(\mathsf{ThetaAnsari})\wedge\mathrm{ContinuumCard}(\mathsf{ThetaAnsari})$ & 2 & \cref{lem:wright-facts}(b,c,d) \\ \hline
4 & $\mathrm{IsNull}(\mathsf{ThetaWright})\wedge\mathrm{NowhereDense}(\mathsf{ThetaWright})\wedge\mathrm{ContinuumCard}(\mathsf{ThetaWright})$ & 2 & \cref{lem:wright-facts}(b,c,d) \\
\end{tabular}%
\noindent \newline
\noindent
\begin{tabular}{|m{0.7cm}|m{12cm}|m{1.4cm}|m{2.5cm}|}
\hline
5 & $\mathrm{IsNull}(\mathsf{ThetaNiven})\wedge\mathrm{NowhereDense}(\mathsf{ThetaNiven})\wedge\mathrm{ContinuumCard}(\mathsf{ThetaNiven})$ & 2 & \cref{lem:wright-facts}(b,c,d) \\ \hline
6 & Conjunction of 3,4,5 & 3,4,5 & $\wedge\mathsf{I}$ \\
\end{tabular}%
\noindent \newline
\noindent
\begin{tabular}{|m{0.7cm}|m{12cm}|m{1.4cm}|m{2.5cm}|}
\hline
7 & Conclude & 1--6 & $\rightarrow\mathsf{I}$ \\ \hline
\end{tabular}%
}
\end{proof}

% =========================================================================
\subsection{Route 9: Proofs of the Lean-form theorems}
\label{sec:route-lean}
% =========================================================================

\begin{proof}[\textbf{Proof of \cref{thm:lean-exists}}]
\phantomsection\label{proof:lean-exists}
{\scriptsize
\begin{tabular}{|m{0.7cm}|m{12cm}|m{1.4cm}|m{2.5cm}|}
\hline
1 & $\text{Ingham}$ & --- & $\mathsf{Hyp}$ \\ \hline
2 & $\exists\theta>1,\ \mathrm{IsMills}(\theta)$ & 1 & \cref{thm:mills} \\
\end{tabular}%
\noindent \newline
\noindent
\begin{tabular}{|m{0.7cm}|m{12cm}|m{1.4cm}|m{2.5cm}|}
\hline
3 & $\exists A:\mathbb{R}_{>1},\ \mathrm{IsMills}(A)$ & 2 & $\exists\mathsf{I}$ \\ \hline
4 & Conclude & 1--3 & $\rightarrow\mathsf{I}$ \\
\end{tabular}%
}
\end{proof}

\begin{proof}[\textbf{Proof of \cref{thm:lean-least}}]
\phantomsection\label{proof:lean-least}
{\scriptsize
\begin{tabular}{|m{0.7cm}|m{12cm}|m{1.4cm}|m{2.5cm}|}
\hline
1 & $\exists A:\mathbb{R}_{>1},\ \mathrm{IsMills}(A)$ & --- & $\mathsf{Hyp}$ \\ \hline
2 & $S:=\{A\in\mathbb{R}:\ \mathrm{IsMills}(A)\}\ne\emptyset$, $S\subseteq(1,\infty)$ & 1 & $\mathsf{Def}$ \\
\end{tabular}%
\noindent \newline
\noindent
\begin{tabular}{|m{0.7cm}|m{12cm}|m{1.4cm}|m{2.5cm}|}
\hline
3 & $S$ closed in $[1+\epsilon,\infty)$ for every $\epsilon>0$ (finite intersection of closed level sets of $\lfloor A^{3^n}\rfloor$) & 2 & $\mathsf{Supp}$ (Caldwell--Cheng 2005) \\ \hline
4 & $S$ bounded below by 1 & 2 & $\mathsf{Def}$ \\
\end{tabular}%
\noindent \newline
\noindent
\begin{tabular}{|m{0.7cm}|m{12cm}|m{1.4cm}|m{2.5cm}|}
\hline
5 & $L:=\inf S$ exists & 3,4 & $\mathsf{Ax}$ (Dedekind) \\ \hline
6 & $L\in S$ (compactness of $S\cap[1,L+1]$) & 3,5 & $\mathsf{Supp}$ (closedness + boundedness) \\
\end{tabular}%
\noindent \newline
\noindent
\begin{tabular}{|m{0.7cm}|m{12cm}|m{1.4cm}|m{2.5cm}|}
\hline
7 & $\forall B\in S$, $L\le B$ & 5 & $\mathsf{def}(\inf)$ \\ \hline
8 & $\mathrm{IsMinMills}(L)$ & 6,7 & $\mathsf{Def}$ \\
\end{tabular}%
\noindent \newline
\noindent
\begin{tabular}{|m{0.7cm}|m{12cm}|m{1.4cm}|m{2.5cm}|}
\hline
9 & $\exists A:\mathbb{R},\ \mathrm{IsMinMills}(A)$ & 8 & $\exists\mathsf{I}$ \\ \hline
10 & Conclude & 1--9 & $\rightarrow\mathsf{I}$ \\
\end{tabular}%
}
\end{proof}

\begin{proof}[\textbf{Proof of \cref{thm:lean-irrational}}]
\phantomsection\label{proof:lean-irrational}
{\scriptsize
\begin{tabular}{|m{0.7cm}|m{12cm}|m{1.4cm}|m{2.5cm}|}
\hline
1 & $\exists A:\mathbb{R},\ \mathrm{IsMinMills}(A)$ & --- & $\mathsf{Hyp}$ \\ \hline
2 & $A$ least Mills constant & 1 & $\mathsf{Def}$ \\
\end{tabular}%
\noindent \newline
\noindent
\begin{tabular}{|m{0.7cm}|m{12cm}|m{1.4cm}|m{2.5cm}|}
\hline
3 & Suppose $A=p/q\in\mathbb{Q}$, $p,q\in\mathbb{N}$, $q\ge1$ & 2 & $\mathsf{RAA}$ \\ \hline
4 & $qA=p\in\mathbb{Z}$, so $q^{3^n}A^{3^n}=p^{3^n}$ & 3 & $\mathsf{Eq}$ \\
\end{tabular}%
\noindent \newline
\noindent
\begin{tabular}{|m{0.7cm}|m{12cm}|m{1.4cm}|m{2.5cm}|}
\hline
5 & $\mathrm{floor}(A^{3^n})\in(p^{3^n}/q^{3^n}-1,\ p^{3^n}/q^{3^n})$ for all $n$ & 4 & $\mathsf{Def}$ \\ \hline
6 & For $n$ large, denominator $q^{3^n}$ dominates: no prime in the interval (Saito 2025, via gap estimates) & 3,5 & $\mathsf{Supp}$ (Saito) \\
\end{tabular}%
\noindent \newline
\noindent
\begin{tabular}{|m{0.7cm}|m{12cm}|m{1.4cm}|m{2.5cm}|}
\hline
7 & Contradiction with $\mathrm{IsMills}(A)$ & 2,6 & $\mathsf{Eq}$ \\ \hline
8 & Therefore $A\notin\mathbb{Q}$, i.e.\ $\mathrm{Irrational}(A)$ & 3--7 & $\mathsf{RAA}$ \\
\end{tabular}%
\noindent \newline
\noindent
\begin{tabular}{|m{0.7cm}|m{12cm}|m{1.4cm}|m{2.5cm}|}
\hline
9 & $\forall A,\ \mathrm{IsMinMills}(A)\Rightarrow\mathrm{Irrational}(A)$ & 1--8 & $\forall\mathsf{I}$ \\ \hline
10 & Conclude & 1--9 & $\rightarrow\mathsf{I}$ \\
\end{tabular}%
}
\end{proof}

\begin{proof}[\textbf{Proof of \cref{thm:lean-rh}}]
\phantomsection\label{proof:lean-rh}
{\scriptsize
\begin{tabular}{|m{0.7cm}|m{12cm}|m{1.4cm}|m{2.5cm}|}
\hline
1 & $\mathrm{RH}$, $\exists A:\mathbb{R},\ \mathrm{IsMinMills}(A)$ & --- & $\mathsf{Hyp}$ \\ \hline
2 & $A$ least Mills constant; under RH, $\pi(x)=\mathrm{li}(x)+O(\sqrt{x}\ln x)$ & 1 & $\mathsf{Supp}$ (von Koch 1901) \\
\end{tabular}%
\noindent \newline
\noindent
\begin{tabular}{|m{0.7cm}|m{12cm}|m{1.4cm}|m{2.5cm}|}
\hline
3 & Explicit gap bound under RH: $p_{n+1}-p_n<2\sqrt{p_n}\ln p_n$ for $p_n\ge p_0$ & 2 & $\mathsf{Supp}$ (RH consequence) \\ \hline
4 & Re-run \cref{proof:mills} with gap constant $k=2$: some $q_1>k^8$ giving a Mills constant $A_0$ & 3 & $\mathsf{Supp}$ (proof:mills with explicit $k$) \\
\end{tabular}%
\noindent \newline
\noindent
\begin{tabular}{|m{0.7cm}|m{12cm}|m{1.4cm}|m{2.5cm}|}
\hline
5 & Numerical evaluation of $A_0$: $1.3063778838<A_0<1.3063778839$ & 4 & $\mathsf{Supp}$ (Caldwell--Cheng 2005, \S 3) \\ \hline
6 & Any Mills constant $\ge A_0\ge 1.3063778838$ & 4,5 & $\mathsf{def}$ (minimality) \\
\end{tabular}%
\noindent \newline
\noindent
\begin{tabular}{|m{0.7cm}|m{12cm}|m{1.4cm}|m{2.5cm}|}
\hline
7 & $A\le A_0<1.3063778839$ and $A\ge A_0>1.3063778838$ & 5,6 & $\mathsf{Eq}$ \\ \hline
8 & $\forall A,\ \mathrm{IsMinMills}(A)\Rightarrow 1.3063778838<A<1.3063778839$ & 1--7 & $\forall\mathsf{I}$ \\
\end{tabular}%
\noindent \newline
\noindent
\begin{tabular}{|m{0.7cm}|m{12cm}|m{1.4cm}|m{2.5cm}|}
\hline
9 & Conclude & 1--8 & $\rightarrow\mathsf{I}$ \\ \hline
\end{tabular}%
}
\end{proof}

\begin{proof}[\textbf{Proof of \cref{thm:lean-detailed}}]
\phantomsection\label{proof:lean-detailed}
{\scriptsize
\begin{tabular}{|m{0.7cm}|m{12cm}|m{1.4cm}|m{2.5cm}|}
\hline
1 & $\mathrm{MillsRec}(q)\wedge q_1>k^8$ ($k$ from Ingham) & --- & $\mathsf{Hyp}$ \\ \hline
2 & $\forall n$: $q_n^3<q_{n+1}<(q_n+1)^3-1$ & 1 & $\mathsf{Def}$ \\
\end{tabular}%
\noindent \newline
\noindent
\begin{tabular}{|m{0.7cm}|m{12cm}|m{1.4cm}|m{2.5cm}|}
\hline
3 & $u_n:=q_n^{3^{-n}}$, $v_n:=(q_n+1)^{3^{-n}}$ & 2 & $\mathsf{def}$ \\ \hline
4 & $u_n<u_{n+1}<v_{n+1}<v_n$ & 2,3 & \cref{thm:millsseq} \\
\end{tabular}%
\noindent \newline
\noindent
\begin{tabular}{|m{0.7cm}|m{12cm}|m{1.4cm}|m{2.5cm}|}
\hline
5 & $\exists\theta>0$, $u_n\to\theta$, $v_n\to\theta$, $u_n<\theta<v_n$ & 3,4 & \cref{lem:squeeze}b \\ \hline
6 & $\forall n$: $q_n<\theta^{3^n}<q_n+1$ & 3,5 & $\mathsf{Eq}$ \\
\end{tabular}%
\noindent \newline
\noindent
\begin{tabular}{|m{0.7cm}|m{12cm}|m{1.4cm}|m{2.5cm}|}
\hline
7 & $\forall n$: $q_n=\mathrm{floor}(\theta^{3^n})$ & 6 & $\mathsf{Def}$ \\ \hline
8 & $\forall n$: $\mathrm{Prime}(\mathrm{floor}(\theta^{3^n}))$ & 2,7 & $\mathsf{Def}$ \\
\end{tabular}%
\noindent \newline
\noindent
\begin{tabular}{|m{0.7cm}|m{12cm}|m{1.4cm}|m{2.5cm}|}
\hline
9 & $\theta>1$, so $\mathrm{IsMills}(\theta)$ & 2,5,8 & $\mathsf{Eq}$ \\ \hline
10 & $\exists\theta>0$, $(\forall n,\ q_n=\mathrm{floor}(\theta^{3^n}))\wedge\mathrm{IsMills}(\theta)$ & 7,9 & $\exists\mathsf{I}$ \\
\end{tabular}%
\noindent \newline
\noindent
\begin{tabular}{|m{0.7cm}|m{12cm}|m{1.4cm}|m{2.5cm}|}
\hline
11 & Conclude & 1--10 & $\rightarrow\mathsf{I}$ \\ \hline
\end{tabular}%
}
\end{proof}

% =========================================================================
\section{Scripts}
\label{sec:scripts}

\subsection{Lean scripts}
\label{sec:scripts-lean}

The following Lean~4 / Mathlib-style formalisation is the intended target.
It is not a compilable file in this repository (no Mathlib import is
provided); it records the type signatures and theorem statements
corresponding to \cref{sec:thm-lean}.

\begin{lstlisting}[style=lean]
-- Lean 4 / Mathlib-style sketch of Mills' theorem.
import Mathlib

namespace Mills

/-- Primality (decidable predicate on Nat). -/
def IsPrime (p : Nat) : Prop :=
2 <= p ∧ ∀ d : Nat, 2 <= d -> d ∣ p -> d = p

/-- The Mills predicate on reals: floor(θ^(3^n)) is prime for all n ≥ 1. -/
def IsMills (θ : ℝ) : Prop :=
1 < θ ∧ ∀ n : Nat, 0 < n ->
IsPrime (Int.toNat (Int.floor (θ ^ (3 ^ n))))

/-- The least Mills constant predicate. -/
def IsMinMills (θ : ℝ) : Prop :=
IsMills θ ∧ ∀ θ' : ℝ, IsMills θ' -> θ <= θ'

/-- Ingham's estimate: ∃ k > 0, ∀ n, p_{n+1} - p_n < k * p_n^{5/8}. -/
axiom Ingham :
∃ k : ℝ, 0 < k ∧
∀ n : Nat, 0 < n ->
(Nat.nth Nat.Prime (n+1) - Nat.nth Nat.Prime n : ℝ)
< k * (Nat.nth Nat.Prime n : ℝ) ^ (5/8 : ℝ)

/-- Mills' theorem in Lean-form: from Ingham's estimate, there exists a
real θ > 1 with the Mills property. -/
theorem mills (h : Ingham) : ∃ θ : ℝ, 1 < θ ∧ IsMills θ := by
sorry

/-- Existence of the least Mills constant. -/
theorem least_mills (h : ∃ θ : ℝ, 1 < θ ∧ IsMills θ) :
∃ θ : ℝ, IsMinMills θ := by
sorry

/-- The least Mills constant is irrational (Saito 2025). -/
theorem min_mills_irrational
(h : ∃ θ : ℝ, IsMinMills θ) :
∀ θ : ℝ, IsMinMills θ -> Irrational θ := by
sorry

/-- If RH holds, the least Mills constant lies in a narrow interval
(Caldwell–Cheng 2005). -/
axiom RH : Prop
theorem min_mills_under_rh
(hRH : RH) (h : ∃ θ : ℝ, IsMinMills θ) :
∀ θ : ℝ, IsMinMills θ ->
(1.3063778838 : ℝ) < θ ∧ θ < (1.3063778839 : ℝ) := by
sorry

/-- Lean-adapted Mills limit: the limit of the minimal Mills sequence
satisfies the floor identity. -/
theorem mills_limit
(q : Nat -> Nat) (hrec : ∀ n : Nat, 0 < n ->
(q n : ℝ) ^ 3 < q (n+1) ∧
(q (n+1) : ℝ) < ((q n : ℝ) + 1) ^ 3 - 1) :
∃ θ : ℝ, 0 < θ ∧
(∀ n : Nat, 0 < n -> q n = Int.toNat (Int.floor (θ ^ (3 ^ n))))
∧ IsMills θ := by
sorry

end Mills
\end{lstlisting}

\subsection{Python scripts}
\label{sec:scripts-python}

\begin{lstlisting}[style=python]
#!/usr/bin/env python3
"""
verify_mills.py -- numerical verification of the ten routes to
Mills' theorem (see Sections 5-6 of the paper).

Usage:
python3 verify_mills.py [N]

N is the threshold (max integer N and max prime index).  Default N = 50.
Only the Python standard library is required.
"""

import sys
import math

# ----------------------------------------------------------------
# Primality and prime utilities (trial division, fine for small N)
# ----------------------------------------------------------------
def is_prime(m: int) -> bool:
if m < 2:
return False
if m % 2 == 0:
return m == 2
r = math.isqrt(m)
for d in range(3, r + 1, 2):
if m % d == 0:
return False
return True

def first_primes(count: int):
out, m = [], 2
while len(out) < count:
if is_prime(m):
out.append(m)
m += 1
return out

def smallest_prime_in(lo: int, hi: int):
p = lo + 1
while p < hi:
if is_prime(p):
return p
p += 1
return None

def prime_in(lo: int, hi: int) -> bool:
return smallest_prime_in(lo, hi) is not None

# ----------------------------------------------------------------
# Route 1: Ingham
# ----------------------------------------------------------------
def check_ingham(N: int):
ps = first_primes(N + 1)
k_min = 0.0
for n in range(N):
p, q = ps[n], ps[n + 1]
k_min = max(k_min, (q - p) / (p ** (5 / 8)))
ok = all(math.isfinite((ps[n+1]-ps[n]) / (ps[n]**(5/8)))
for n in range(N))
return ok, f"max gap/p^(5/8) = {k_min:.6f}; any k >= {k_min:.4f} works"

# ----------------------------------------------------------------
# Route 1 (Mills interval + recurrence)
# ----------------------------------------------------------------
def check_mills_interval(N: int):
bad = [n for n in range(2, N + 1)
if not prime_in(n ** 3, (n + 1) ** 3 - 1)]
return (not bad), f"{N-1} intervals; failures: {bad or 'none'}"

def check_mills_sequence(N: int, seed: int = 2):
qs = [seed]
for _ in range(N - 1):
nxt = smallest_prime_in(qs[-1] ** 3, (qs[-1] + 1) ** 3 - 1)
if nxt is None:
return False, f"no prime after q={qs[-1]}"
qs.append(nxt)
def mono(exp):
us = [q ** (exp ** -n) for n, q in enumerate(qs, 1)]
vs = [(q + 1) ** (exp ** -n) for n, q in enumerate(qs, 1)]
return (all(us[i] < us[i+1] for i in range(len(us)-1))
and all(vs[i] > vs[i+1] for i in range(len(vs)-1))
and all(us[i] < vs[i] for i in range(len(us))))
ok3, ok5 = mono(3), mono(5)
msg = (f"q_1..q_{len(qs)}={qs[:5]}{'...' if len(qs)>5 else ''}\n"
f"  exponent 3^-n (paper): {'PASS' if ok3 else 'FAIL'}\n"
f"  exponent 5^-n (Dudley typo): {'PASS' if ok5 else 'FAIL'}")
return ok3, msg

# ----------------------------------------------------------------
# Route 2: Kuipers
# ----------------------------------------------------------------
def check_kuipers(N: int, c: int):
bad = [n for n in range(2, N + 1)
if not prime_in(n ** c, (n + 1) ** c - 1)]
return (not bad), f"c={c}: {N-1} intervals; failures: {bad or 'none'}"

# ----------------------------------------------------------------
# Route 3: Ansari
# ----------------------------------------------------------------
def check_ansari(N: int, c: float = 3.0):
delta = 3 * c / 8 - 1
if delta <= 0:
return False, f"c={c} does not satisfy c>8/3"
bad = [n for n in range(2, N + 1)
if not (n**c + n**(c-1) < (n+1)**c - delta)]
return (not bad), f"c={c}, delta={delta:.4f}; failures: {bad or 'none'}"

# ----------------------------------------------------------------
# Route 4: Bertrand
# ----------------------------------------------------------------
def check_bertrand(N: int):
bad = [n for n in range(2, N + 1) if not prime_in(n, 2 * n)]
return (not bad), f"{N-1} values; failures: {bad or 'none'}"

# ----------------------------------------------------------------
# Route 4 (Wright tower)
# ----------------------------------------------------------------
def check_wright(N: int, cap: int = 10**7):
qs = [2]
for _ in range(N - 1):
lo, hi = 2 ** qs[-1], 2 ** (qs[-1] + 1)
if hi > cap:
break
nxt = smallest_prime_in(lo, hi)
if nxt is None:
return False, f"no prime in (2^{qs[-1]}, 2^{qs[-1]+1})"
qs.append(nxt)
def log2(x): return math.log2(x)
def it_log(n, x):
for _ in range(n):
x = log2(x)
return x
us = [it_log(n, q) for n, q in enumerate(qs, 1)]
vs = [it_log(n, q + 1) for n, q in enumerate(qs, 1)]
inc_u = all(us[i] < us[i+1] for i in range(len(us)-1))
dec_v = all(vs[i] > vs[i+1] for i in range(len(vs)-1))
return (inc_u and dec_v), f"{len(qs)} terms; u inc={inc_u}, v dec={dec_v}"

# ----------------------------------------------------------------
# Route 6: polynomial gap
# ----------------------------------------------------------------
def check_polygap(N: int, alpha: float, k: float, c: float):
if not (c > 1 / (1 - alpha)):
return False, f"c={c} does not satisfy c>1/(1-alpha)={1/(1-alpha):.4f}"
# Sample check: existence of prime in (N^c, (N+1)^c - 1) for small N.
bad = []
for n in range(2, N + 1):
lo, hi = n ** c, (n + 1) ** c - 1
# For non-integer c, we just check (N^c, (N+1)^c - 1) has a
# prime; this is a sanity check, not a proof.
if hi > 10 ** 7:
break
if not prime_in(int(lo), int(hi) + 1):
bad.append(n)
return (not bad), f"alpha={alpha}, c={c}: sampling failures={bad or 'none'}"

# ----------------------------------------------------------------
# Driver
# ----------------------------------------------------------------
def main():
N = int(sys.argv[1]) if len(sys.argv) > 1 else 50
print(f"=== Mills' Theorem: numerical verification, N = {N} ===\n")
checks = [
("1. Ingham", lambda: check_ingham(N)),
("1. Mills interval", lambda: check_mills_interval(N)),
("1. Mills sequence (3^-n vs 5^-n)",
lambda: check_mills_sequence(min(N, 15))),
("2. Kuipers c=3", lambda: check_kuipers(N, 3)),
("2. Kuipers c=4", lambda: check_kuipers(N, 4)),
("3. Ansari c=3", lambda: check_ansari(N, 3.0)),
("4. Bertrand", lambda: check_bertrand(N)),
("4. Wright tower", lambda: check_wright(N)),
("6. Polynomial gap (alpha=5/8, c=3)",
lambda: check_polygap(min(N, 15), 5/8, 1.0, 3.0)),
]
for label, fn in checks:
try:
ok, msg = fn()
print(f"[{'PASS' if ok else 'FAIL'}] {label}\n    {msg}\n")
except Exception as e:
print(f"[ERROR] {label}: {e}\n")

if __name__ == "__main__":
main()
\end{lstlisting}

\section{Appendix}
\label{sec:appendix}
% =========================================================================

\subsection{Tables}

	% ---------------------------------------------------------------
% Definition table
% ---------------------------------------------------------------
{\scriptsize \setstretch{1}  
\begin{longtable}{|p{4.2cm}|p{11.2cm}|}
\hline	\textbf{Predicate} & \textbf{Formula}\\
	\hline
	\endfirsthead
	\hline
	\textbf{Predicate} & \textbf{Formula}\\
	\hline
	\endhead
	\hline
	\endfoot
	$\mathsf{Prime}(p)$ &
	$p\ge2\ \wedge\ \forall d\in\mathbb N,\,
	((d\ge2\wedge d\mid p)\Rightarrow d=p)$\\
	\hline
	$\mathsf{Floor}(x,m)$ &
	$m\le x<m+1$\\
	\hline
	$\mathsf{PSeq}(p)$ &
	$p_1=2\ \wedge\
	(\forall n\ge1,\mathsf{Prime}(p_n))\ \wedge\
	(\forall n\ge1,p_n<p_{n+1})$\\
	\hline
	$\mathsf{IsMills}(A)$ &
	$A>1\ \wedge\
	\forall n\ge1,\mathsf{Prime}(\floor{A^{3^n}})$\\
	\hline
	$\mathsf{MillsRec}_c(q)$ &
	$c>1\ \wedge\ \forall n\ge1,\ q_n^{c}<q_{n+1}<(q_n+1)^{c}-1$\\
	\hline
	$\mathsf{MillsRec}(q)$ &
	$\mathsf{MillsRec}_3(q)$\\
	\hline
	$\mathsf{UV}_c(q,u,v)$ &
	$\forall n\ge1,\
	u_n=q_n^{c^{-n}}\ \wedge\
	v_n=(q_n+1)^{c^{-n}}$\\
	\hline
	$\mathsf{UV}(q,u,v)$ &
	$\mathsf{UV}_3(q,u,v)$\\
	\hline
	$\mathsf{Ingham}$ &
	$\exists k>0,\ \forall n\ge1,\
	p_{n+1}-p_n<kp_n^{5/8}$\\
	\hline
	$\mathsf{Ring}(R,+,\cdot)$ &
	$\begin{aligned}\setstretch{1} &\bigl(\forall a,b,c\in R,\ (a+b)+c=a+(b+c)\bigr)\ \wedge\
	\bigl(\exists 0\in R,\ \forall a\in R,\ a+0=a\bigr)\ \\& \wedge\
	\bigl(\forall a\in R,\ \exists(-a)\in R,\ a+(-a)=0\bigr)\ \wedge\
	\bigl(\forall a,b\in R,\ a+b=b+a\bigr)\ \\& \wedge\
	\bigl(\forall a,b,c\in R,\ (ab)c=a(bc)\bigr)\ \wedge\
	\bigl(\forall a,b,c\in R,\ a(b+c)=ab+ac\\& \wedge\ (a+b)c=ac+bc\bigr)\end{aligned}$\\
	\hline
	$\mathsf{LinearOrder}(R,<)$ &
	$\begin{aligned}\setstretch{1} &\bigl(\forall a\in R,\ \neg(a<a)\bigr)\ \wedge\
	\bigl(\forall a,b,c\in R,\ (a<b\wedge b<c)\Rightarrow a<c\bigr)\ \\& \wedge\
	\bigl(\forall a,b\in R,\ a<b\ \vee\ a=b\ \vee\ b<a\bigr)\end{aligned}$\\
	\hline
	$\mathsf{RingOrder}(R,+,\cdot,<)$ &
	$\mathsf{Ring}(R,+,\cdot)\ \wedge\
	\mathsf{LinearOrder}(R,<)\ \wedge\
	\forall a,b,c\in R,\
	(a<b\wedge 0<c)\Rightarrow ac<bc$\\
	\hline
	$\mathsf{OrderPreservation}(f,I)$ &
	$\forall x,y\in I,\ (x<y)\Rightarrow f(x)<f(y)$\\
	\hline
	$\mathsf{PowMonotone}(r)$ &
	$\forall x,y>0,\ (x<y\wedge r>0)\Rightarrow x^r<y^r$\\
	\hline
	$\mathsf{PowerLaws}$ &
	$\forall x>0,\ \forall a,b\in\mathbb R,\
	(x^a)^b=x^{ab}$\\
	\hline
	$\mathsf{Archimedean}(\mathbb R)$ &
	$\forall x\in\mathbb R,\ \exists N\in\mathbb N,\ x<N$\\
	\hline
	$\mathsf{PowerGrowth}$ &
	$\forall \gamma>0,\ \forall C>0,\ \exists N_0,\
	\forall N\ge N_0,\ C<N^\gamma$\\
	\hline
	$\mathsf{Induction}(P)$ &
	$\bigl(P(1)\wedge
	\forall n\ge1\,(P(n)\Rightarrow P(n+1))\bigr)
	\Rightarrow\forall n\ge1,\ P(n)$\\
	\hline
	$\mathsf{WellOrder}(\mathbb N)$ &
	$\forall S\subseteq\mathbb N,\
	(S\ne\varnothing)\Rightarrow\exists m\in S,\
	\forall n\in S,\ m\le n$\\
	\hline
	$\mathsf{Conv}(u,\theta)$ &
	$\forall \varepsilon>0,\ \exists N_1\in\mathbb N,\
	\forall n\ge N_1,\ |u_n-\theta|<\varepsilon$\\
	\hline
	$u_n\to\theta$ &
	$\mathsf{Conv}(u,\theta)$\\
	\hline
	$\mathsf{Squeeze}$ &
	$\bigl(\forall n\ge1,\ u_n\le w_n\le v_n\bigr)\ \wedge\
	\mathsf{Conv}(u,\theta)\ \wedge\ \mathsf{Conv}(v,\theta)
	\Rightarrow \mathsf{Conv}(w,\theta)$\\
	\hline
	$\mathsf{MCT}$ &
	$\bigl(\forall n,\ u_n\le u_{n+1}\bigr)\ \wedge\
	\bigl(\forall n,\ v_{n+1}\le v_n\bigr)\ \wedge\
	\bigl(\forall n,\ u_n\le v_n\bigr)
	\Rightarrow \exists\theta,\ \mathsf{Conv}(u,\theta)\wedge\mathsf{Conv}(v,\theta)$\\
	\hline
	$\mathsf{MillsSqueeze}(q,u,v,\theta)$ &
	$\forall n\ge1,\
	u_n<u_{n+1}<v_{n+1}<v_n
	\wedge\ u_n<\theta<v_n
	\wedge\ (u_n\to\theta)
	\wedge\ (v_n\to\theta)$\\
	\hline
	$\mathsf{Bertrand}$ &
	$\forall N\ge2,\ \exists p,\ \mathsf{Prime}(p)\wedge N<p<2N$\\
	\hline
	$\mathsf{slope}(\beta)$ &
	$\beta>1\ \wedge\ \exists N_0,\ \forall N\ge N_0,\ N^{\beta}+N^{\beta-1}<(N+1)^{\beta}-1$\\
	\hline
	$\mathsf{PolyGap}(\alpha,k)$ &
	$\alpha\in(0,1)\ \wedge\ k>0\ \wedge\ \forall n\ge1,\
	p_{n+1}-p_n<k\,p_n^{\alpha}$\\
	\hline
	$\mathsf{IsKuipers}(c,\theta)$ &
	$c\ge3\ \wedge\ \theta>1\ \wedge\ \forall n\ge1,\
	\mathsf{Prime}(\floor{\theta^{c^n}})$\\
	\hline
	$\mathsf{IsAnsari}(c,\theta)$ &
	$c>8/3\ \wedge\ \theta>1\ \wedge\ \forall n\ge1,\
	\mathsf{Prime}(\floor{\theta^{c^n}})$\\
	\hline
	$\mathsf{IsWright}(\theta)$ &
	$\theta>1\ \wedge\ \forall n\ge1,\
	\mathsf{Prime}(\floor{\exp_2^{[n]}(\theta)})$\\
	\hline
	$\log_2^{[n]}(x)$ &
	$\log_2^{[0]}(x)=x,\ \ \log_2^{[n+1]}(x)=\log_2\bigl(\log_2^{[n]}(x)\bigr)$\\
	\hline
	$\exp_2^{[n]}(x)$ &
	$\exp_2^{[0]}(x)=x,\ \ \exp_2^{[n+1]}(x)=2^{\exp_2^{[n]}(x)}$\\
	\hline
	$\mathsf{Lem}(\cdot)$ &
	$\mathsf{Lem}(\cdot)\equiv\mathsf{Lemma}(\cdot)$\\
	\hline
	$\mathsf{Def}(\cdot)$ &
	$\mathsf{Def}(\cdot)\equiv\mathsf{Unfold}(\cdot)$\\
	\hline
	$\mathsf{Eq}(\cdot)$ &
	$\mathsf{Eq}(\cdot)\equiv\mathsf{Rewrite}(\cdot)$\\
	\hline
	$\mathsf{Hyp}$ &
	$\mathsf{Hyp}\equiv\mathsf{Assume}$\\
	\hline
	$\mathsf{def}(\cdot)$ &
	$\mathsf{def}(\cdot)\equiv\mathsf{Unfold}(\cdot)$\\
	\hline
	$\exists\mathsf{E}$, $\exists\mathsf{I}$, $\forall\mathsf{I}$, $\rightarrow\mathsf{I}$, $\wedge\mathsf{E}$ &
	$\mathsf{Rule}(\cdot)$\\
	\hline
\end{longtable}

% ---------------------------------------------------------------
% Unified proof table: all theorems and lemmas
% ---------------------------------------------------------------

\begin{longtable}{|p{3cm}|c|p{10cm}|p{4cm}|}
\hline	\label{tab:all-proofs}
\textbf{Block} & \textbf{Step} & \textbf{Formula} & \textbf{Dependency}\\
	\hline
	\endfirsthead
	\hline
	\textbf{Block} & \textbf{Step} & \textbf{Formula} & \textbf{Dependency}\\
	\hline
	\endhead
	\hline
	\endfoot
	% ----- Block 1: MillsMono -----
	\multirow{14}{*}{$\mathsf{MillsMono}$} &
	1 & $\mathsf{UV}(q,u,v)\wedge\mathsf{MillsRec}(q)$ & $\mathsf{Hyp}$\\
	\cline{2-4}
	& 2 & $\mathsf{UV}(q,u,v)$ & $1,\ \wedge\mathsf{E}$\\
	\cline{2-4}
	& 3 & $\mathsf{MillsRec}(q)$ & $1,\ \wedge\mathsf{E}$\\
	\cline{2-4}
	& 4 & $\forall n\ge1,\ u_n=q_n^{3^{-n}}\wedge v_n=(q_n+1)^{3^{-n}}$ & $2,\ \mathsf{def}(\mathsf{UV})$\\
	\cline{2-4}
	& 5 & $\forall n\ge1,\ q_n^3<q_{n+1}<(q_n+1)^3-1$ & $3,\ \mathsf{def}(\mathsf{MillsRec})$\\
	\cline{2-4}
	& 6 & $\forall n\ge1,\ q_n^3<q_{n+1}\wedge q_{n+1}+1<(q_n+1)^3$ & $5,\ \mathsf{RingOrder}$\\
	\cline{2-4}
	& 7 & $\forall n\ge1,\ (q_n^3)^{3^{-(n+1)}}<q_{n+1}^{3^{-(n+1)}}$ & $6,\ \mathsf{PowMonotone}$\\
	\cline{2-4}
	& 8 & $\forall n\ge1,\ q_n^{3^{-n}}<q_{n+1}^{3^{-(n+1)}}$ & $7,\ \mathsf{PowerLaws}$\\
	\cline{2-4}
	& 9 & $\forall n\ge1,\ (q_{n+1}+1)^{3^{-(n+1)}}<((q_n+1)^3)^{3^{-(n+1)}}$ & $6,\ \mathsf{PowMonotone}$\\
	\cline{2-4}
	& 10 & $\forall n\ge1,\ (q_{n+1}+1)^{3^{-(n+1)}}<(q_n+1)^{3^{-n}}$ & $9,\ \mathsf{PowerLaws}$\\
	\cline{2-4}
	& 11 & $\forall n\ge1,\ q_n^{3^{-n}}<(q_n+1)^{3^{-n}}$ & $\mathsf{RingOrder},\ \mathsf{PowMonotone}$\\
	\cline{2-4}
	& 12 & $\forall n\ge1,\ u_n<u_{n+1}\wedge v_{n+1}<v_n\wedge u_n<v_n$ & $4,8,10,11$\\
	\cline{2-4}
	& 13 & $\forall n\ge1,\ u_n<u_{n+1}<v_{n+1}<v_n$ & $4,12,\ \mathsf{PowerLaws}$\\
	\cline{2-4}
	& 14 & $(\mathsf{UV}(q,u,v)\wedge\mathsf{MillsRec}(q))\Rightarrow\forall n\ge1,\ u_n<u_{n+1}<v_{n+1}<v_n$ & $1\text{--}13,\ \rightarrow\mathsf{I}$\\
	\hline
	
	% ----- Block 2: KuipersRoute -----
	\multirow{11}{*}{$\mathsf{KuipersRoute}$} &
	1 & $\Ingham$ & $\mathsf{Hyp}$\\
	\cline{2-4}
	& 2 & $c\in\mathbb{N}_{\ge 3}$ & $\mathsf{def}$\\
	\cline{2-4}
	& 3 & $a:=3c-4,\ b:=3c-1,\ 8a-5b\ge 0$ & $2,\ \mathsf{Eq}$\\
	\cline{2-4}
	& 4 & $ca+1=(c-1)b$ & $3,\ \mathsf{Eq}$\\
	\cline{2-4}
	& 5 & $\exists k_1>0,\ \forall n\ge1,\ p_{n+1}-p_n<k_1p_n^{a/b}$ & $1,3,\ \mathsf{def}$\\
	\cline{2-4}
	& 6 & $N>k_1^{b}\wedge p_n=\max\{p\in\mathbb{N}:\mathsf{Prime}(p)\wedge p<N^c\}$ & $5,\ \exists\mathsf{E}$\\
	\cline{2-4}
	& 7 & $p_{n+1}<N^c+N^{c-1}$ & $5,6,4,\ \mathsf{Eq}$\\
	\cline{2-4}
	& 8 & $N^c+N^{c-1}<(N+1)^c-1$ & $6,\ \mathsf{slope}(c)$\\
	\cline{2-4}
	& 9 & $\exists p,\ \mathsf{Prime}(p)\wedge N^c<p<(N+1)^c-1$ & $6,7,8,\ \exists\mathsf{I}$\\
	\cline{2-4}
	& 10 & $\forall c\ge3,\ \exists\theta>1,\ \mathsf{IsKuipers}(c,\theta)$ & $2\text{--}9,\ \forall\mathsf{I}$\\
	\cline{2-4}
	& 11 & $\Ingham\Rightarrow\forall c\ge3,\ \exists\theta>1,\ \mathsf{IsKuipers}(c,\theta)$ & $1\text{--}10,\ \rightarrow\mathsf{I}$\\
	\hline
	
	% ----- Block 3: AnsariRoute -----
	\multirow{10}{*}{$\mathsf{AnsariRoute}$} &
	1 & $\Ingham$ & $\mathsf{Hyp}$\\
	\cline{2-4}
	& 2 & $c\in\mathbb{R}_{>8/3}$ & $\mathsf{def}$\\
	\cline{2-4}
	& 3 & $\delta:=3c/8-1>0$ & $2,\ \mathsf{Eq}$\\
	\cline{2-4}
	& 4 & $\exists k>0,\ \forall n\ge1,\ p_{n+1}-p_n<kp_n^{5/8}$ & $1,\ \mathsf{def}$\\
	\cline{2-4}
	& 5 & $N>k^{1/\delta}\wedge p_n=\max\{p\in\mathbb{N}:\mathsf{Prime}(p)\wedge p<N^c\}$ & $4,\ \exists\mathsf{E}$\\
	\cline{2-4}
	& 6 & $p_{n+1}<N^c+N^{c-1}$ & $4,5,3,\ \mathsf{Eq}$\\
	\cline{2-4}
	& 7 & $N^c+N^{c-1}<(N+1)^c-1$ & $5,\ \mathsf{slope}(c)$\\
	\cline{2-4}
	& 8 & $\exists p,\ \mathsf{Prime}(p)\wedge N^c<p<(N+1)^c-1$ & $5,6,7,\ \exists\mathsf{I}$\\
	\cline{2-4}
	& 9 & $\forall c>8/3,\ \exists\theta>1,\ \mathsf{IsAnsari}(c,\theta)$ & $2\text{--}8,\ \forall\mathsf{I}$\\
	\cline{2-4}
	& 10 & $\Ingham\Rightarrow\forall c>8/3,\ \exists\theta>1,\ \mathsf{IsAnsari}(c,\theta)$ & $1\text{--}9,\ \rightarrow\mathsf{I}$\\
	\hline
	
	% ----- Block 4: WrightRoute -----
	\multirow{12}{*}{$\mathsf{WrightRoute}$} &
	1 & $\mathsf{Bertrand}$ & $\mathsf{Hyp}$\\
	\cline{2-4}
	& 2 & $\forall N\ge2,\ \exists p,\ \mathsf{Prime}(p)\wedge N<p<2N$ & $1,\ \mathsf{def}$\\
	\cline{2-4}
	& 3 & $q_1=5$ & $2,\ \mathsf{def}$\\
	\cline{2-4}
	& 4 & $\mathsf{Prime}(q_{n+1})\wedge 2^{q_n}<q_{n+1}<2^{q_n+1}$ & $2,3,\ \exists\mathsf{E}$\\
	\cline{2-4}
	& 5 & $u_n=\log_2^{[n]}(q_n)\wedge v_n=\log_2^{[n]}(q_n+1)$ & $4,\ \mathsf{def}$\\
	\cline{2-4}
	& 6 & $\forall n\ge1,\ u_n\le u_{n+1}\wedge v_{n+1}\le v_n\wedge u_n\le v_n$ & $3,4,5,\ \mathsf{Lem}(\log_2^{[n]})$\\
	\cline{2-4}
	& 7 & $\exists\theta,\ \mathsf{Conv}(u,\theta)\wedge\mathsf{Conv}(v,\theta)$ & $6,\ \mathsf{MCT}$\\
	\cline{2-4}
	& 8 & $q_n<\exp_2^{[n]}(\theta)<q_n+1$ & $5,7,\ \mathsf{Lem}(\exp_2^{[n]}\circ\log_2^{[n]})$\\
	\cline{2-4}
	& 9 & $\mathsf{Prime}(\floor{\exp_2^{[n]}(\theta)})\wedge\floor{\exp_2^{[n]}(\theta)}=q_n$ & $4,8,\ \mathsf{Def}$\\
	\cline{2-4}
	& 10 & $\theta>1$ & $3,7,\ \mathsf{Eq}$\\
	\cline{2-4}
	& 11 & $\exists\theta>1,\ \mathsf{IsWright}(\theta)$ & $9,10,\ \exists\mathsf{I}$\\
	\cline{2-4}
	& 12 & $\mathsf{Bertrand}\Rightarrow\exists\theta>1,\ \mathsf{IsWright}(\theta)$ & $1\text{--}11,\ \rightarrow\mathsf{I}$\\
	\hline
	
	% ----- Block 5: PolyGapRoute -----
	\multirow{11}{*}{$\mathsf{PolyGapRoute}$} &
	1 & $\mathsf{PolyGap}(\alpha,k)\wedge c>\tfrac{1}{1-\alpha}$ & $\mathsf{Hyp}$\\
	\cline{2-4}
	& 2 & $\exists N_0,\ \forall N\ge N_0,\ \exists p,\ \mathsf{Prime}(p)\wedge N^c<p<(N+1)^c-1$ & $1,\ \mathsf{Lem}(\mathsf{PolyGap})$\\
	\cline{2-4}
	& 3 & $q_1>N_0\wedge\forall n\ge1,\ \mathsf{Prime}(q_{n+1})\wedge q_n^c<q_{n+1}<(q_n+1)^c-1$ & $2,\ \exists\mathsf{E}$\\
	\cline{2-4}
	& 4 & $\mathsf{MillsRec}_c(q)$ & $3,\ \mathsf{Def}$\\
	\cline{2-4}
	& 5 & $\mathsf{UV}_c(q,u,v)$ & $3,\ \mathsf{def}$\\
	\cline{2-4}
	& 6 & $\forall n\ge1,\ u_n<u_{n+1}<v_{n+1}<v_n$ & $4,5,\ \mathsf{MillsMono}_c$\\
	\cline{2-4}
	& 7 & $\exists\theta>0,\ \mathsf{Conv}(u,\theta)\wedge\mathsf{Conv}(v,\theta)$ & $6,\ \mathsf{MCT}$\\
	\cline{2-4}
	& 8 & $q_n<\theta^{c^n}<q_n+1$ & $5,7,\ \mathsf{Eq}$\\
	\cline{2-4}
	& 9 & $\forall n\ge1,\ \mathsf{Prime}(\floor{\theta^{c^n}})$ & $3,8,\ \mathsf{Def}$\\
	\cline{2-4}
	& 10 & $\exists\theta>0,\ \forall n\ge1,\ \mathsf{Prime}(\floor{\theta^{c^n}})$ & $7,9,\ \exists\mathsf{I}$\\
	\cline{2-4}
	& 11 & $\mathsf{PolyGap}(\alpha,k)\wedge c>\tfrac{1}{1-\alpha}\Rightarrow\exists\theta>0,\ \forall n\ge1,\ \mathsf{Prime}(\floor{\theta^{c^n}})$ & $1\text{--}10,\ \rightarrow\mathsf{I}$\\
	\hline
	
\end{longtable}
\label{sec:appendix-tables}

\begin{table}[H]
	\centering
	\caption{Numerical check of the exponent \(3^{-n}\) (paper) versus \(5^{-n}\)
		(Dudley typo) in the Mills recurrence, using the minimal sequence
		\(q_1=2\), \(q_{n+1}=\min\{p\text{ prime}: q_n^3<p<(q_n+1)^3-1\}\).}
	\label{tab:exponent-typo}
	{\small
		\begin{tabular}{|c|r|r|r|r|}
			\hline
			\(n\) & \(q_n\) & \(u_n=q_n^{3^{-n}}\) & \(u_n^{(5)}=q_n^{5^{-n}}\)
			& Monotone? \\
			\hline
			1 & 2 & \(2^{1/3}\approx 1.2599\) & \(2^{1/5}\approx 1.1487\) & --- \\
			2 & 11 & \(11^{1/9}\approx 1.3048\) & \(11^{1/25}\approx 1.1006\) &
			\(3^{-n}\): inc, \(5^{-n}\): dec \\
			3 & 1361 & \(1361^{1/27}\approx 1.3036\) & \(1361^{1/125}\approx 1.0586\) &
			\(3^{-n}\): inc, \(5^{-n}\): dec \\
			4 & 2521097 & \(2521097^{1/81}\approx 1.3035\) & --- &
			\(3^{-n}\): inc \\
			\hline
		\end{tabular}
	}
\end{table}

\text{\ } \newline
\text{\ } \vspace{0.3in}\text{\ }

\begin{table}[H]
	\centering
	\caption{Summary of the ten routes to \(\exists\theta,\ \mathrm{IsMills}(\theta)\).}
	\label{tab:routes}
	{\small
		\begin{tabular}{|c|c|c|m{2.7cm}|}
			\hline
			\textbf{Route} & \textbf{Author(s)} & \textbf{Year} & \textbf{Asser} \\
			\hline
			1 & Mills & 1947 & \cref{asser:mills} \\ \hline
			2 & Kuipers & 1950 & \cref{asser:kuipers} \\ \hline
			3 & Ansari & 1951 & \cref{asser:ansari} \\ \hline
			4 & Wright & 1951 & \cref{asser:wright} \\ \hline
			5 & Ore & 1952 & \cref{asser:ore} \\ \hline
			6 & Polynomial-gap & --- & \cref{asser:direct} \\ \hline
			7 & Niven & 1951 & \cref{asser:niven} \\ \hline
			8 & Wright (metrical) & 1954 & \cref{asser:wright-metrical} \\ \hline
			9 & Lean (pointer) & 2025 & \cref{asser:lean-exists} \\ \hline
			10 & Lean (detailed) & 2025 & \cref{asser:lean-detailed} \\
			\hline
		\end{tabular}
	}
\end{table}


\text{\ } \newline
\text{\ } \vspace{0.3in}\text{\ }
\begin{table}[H]
	\centering
	\caption{ Ingham ratios}
	\begin{tabular}{rrrr}
		\toprule$n$ & $p_{n}$ & gap & $R_{n}$ \\ 
		\midrule1 & 2 & 1 & 0.648420 \\ 
		2 & 3 & 2 & 1.006560%
	\end{tabular}%
	\noindent \hspace{1em}\\[0pt]
	\noindent 
	\begin{tabular}{rrrr}
		3 & 5 & 2 & 0.731573 \\ 
		4 & 7 & 4 & 1.182766 \\ 
		
	\end{tabular}%
	\noindent \hspace{1em}\\[0pt]
	\noindent 
	\begin{tabular}{rrrr}
		5 & 11 & 2 & 0.446515 \\ 
		6 & 13 & 4 & 0.806455 \\ 
		
	\end{tabular}%
	\noindent \hspace{1em}\\[0pt]
	\noindent 
	\begin{tabular}{rrrr}
		7 & 17 & 2 & 0.340872 \\ 
		8 & 19 & 4 & 0.634271 \\ 
		
	\end{tabular}%
	\noindent \hspace{1em}\\[0pt]
	\noindent 
	\begin{tabular}{rrrr}
		9 & 23 & 6 & 0.842902 \\ 
		10 & 29 & 2 & 0.255064 \\ 
		\bottomrule 
	\end{tabular}
\end{table}
\text{\ } \newline
\text{\ } \vspace{0.3in}\text{\ }



\begin{table}[H]
	\centering
	\caption{Mills intervals}
	\begin{tabular}{rrr}
		\toprule $N$ & $N^{3}$ & \multicolumn{1}{r}{$(N+1)^{3}-1$} \\ 
		\midrule 2 & 8 & \multicolumn{1}{r}{26} \\ 
		3 & 27 & \multicolumn{1}{r}{63} \\ 
		
	\end{tabular}%
	\noindent \hspace{1em}\\[0pt]
	\noindent 
	\begin{tabular}{rrr}
		4 & 64 & \multicolumn{1}{r}{124} \\ 
		5 & 125 & \multicolumn{1}{r}{215} \\ 
		
	\end{tabular}%
	\noindent \hspace{1em}\\[0pt]
	\noindent 
	\begin{tabular}{rrr}
		6 & 216 & \multicolumn{1}{r}{342} \\ 
		7 & 343 & \multicolumn{1}{r}{511} \\ 
		
	\end{tabular}%
	\noindent \hspace{1em}\\[0pt]
	\noindent 
	\begin{tabular}{rrr}
		8 & 512 & \multicolumn{1}{r}{728} \\ 
		9 & 729 & \multicolumn{1}{r}{999} \\ 
		
	\end{tabular}%
	\noindent \hspace{1em}\\[0pt]
	\noindent 
	\begin{tabular}{rrr}
		10 & 1000 & \multicolumn{1}{r}{1330} \\ 
		11 & 1331 & \multicolumn{1}{r}{1727} \\ 
		
	\end{tabular}
	\noindent \hspace{1em}\\[0pt]
	\noindent 
	\begin{tabular}{rrr}
		12 & 1728 & 2196 \\ 
		13 & 2197 & 2743 \\ 
		
	\end{tabular}%
	\noindent \hspace{1em}\\[0pt]
	\noindent 
	\begin{tabular}{rrr}
		14 & 2744 & 3374 \\ 
		15 & 3375 & 4095 \\ 
		
	\end{tabular}%
	\noindent \hspace{1em}\\[0pt]
	\noindent 
	\begin{tabular}{rrr}
		16 & 4096 & 4912 \\ 
		17 & 4913 & 5831 \\ 
		
	\end{tabular}%
	\noindent \hspace{1em}\\[0pt]
	\noindent 
	\begin{tabular}{rrr}
		18 & 5832 & 6858 \\ 
		19 & 6859 & 7999 \\ 
		
	\end{tabular}%
	\noindent \hspace{1em}\\[0pt]
	\noindent 
	\begin{tabular}{rrr}
		20 & 8000 & 9260 \\ 
		
	\end{tabular}
\end{table}

\text{\ } \newline
\text{\ } \vspace{0.3in}\text{\ }




\begin{table}[H]
	\centering
	\caption{Mills monotonicity}
	\begin{tabular}{rrrrrr}
		\toprule $n$ & $q_n$ & $u_n^{(5)}$ & $v_n^{(5)}$ & $u_n^{(3)}$ & $v_n^{(3)}$
		\\ 
		\midrule 1 & 2 & 1.148698 & 1.245731 & 1.259921 & 1.442250 \\ 
		2 & 11 & 1.098745 & 1.102914 & 1.300706 & 1.318699 \\ 
		
	\end{tabular}%
	\noindent \hspace{1em}\\[0pt]
	\noindent 
	\begin{tabular}{rrrrrr}
		3 & 1361 & 1.077244 & 1.077299 & 1.304759 & 1.304914 \\ 
		\bottomrule   
	\end{tabular}
	
\end{table}



\text{\ } \newline
\text{\ } \vspace{0.3in}\text{\ }

\begin{table}[H]
	\caption{Kuipers intervals}
	\centering
	\begin{tabular}{rrr}
		\toprule $N$ & $c=3$ interval & \multicolumn{1}{r}{witness} \\ 
		\midrule 2 & [8,26] & \multicolumn{1}{r}{11} \\ 
		5 & [125,215] & \multicolumn{1}{r}{127} \\ 
		
	\end{tabular}%
	\noindent \hspace{1em}\\[0pt]
	\noindent 
	\begin{tabular}{rrr}
		10 & [1000,1330] & \multicolumn{1}{r}{1009} \\ 
		15 & [3375,4095] & \multicolumn{1}{r}{3391} \\ 
		
	\end{tabular}%
	\noindent \hspace{1em}\\[0pt]
	\noindent 
	\begin{tabular}{rrr}
		20 & [8000,9260] & \multicolumn{1}{r}{8009} \\ 
		
	\end{tabular}
	\noindent \hspace{1em}\\[0pt]
	\noindent 
	\begin{tabular}{rrr}
		$N$ & $c=4$ interval & witness \\ 
		2 & [16,80] & 17 \\ 
		5 & [625,1295] & 631 \\ 
		
	\end{tabular}%
	\noindent \hspace{1em}\\[0pt]
	\noindent 
	\begin{tabular}{rrr}
		10 & [10000,14640] & 10007 \\ 
		15 & [50625,65535] & 50627 \\ 
		
	\end{tabular}%
	\noindent \hspace{1em}\\[0pt]
	\noindent 
	\begin{tabular}{rrr}
		20 & [160000,194480] & 160001 \\ 
		
	\end{tabular}
	
	
	
	\text{\ } \newline
	\text{\ } \vspace{0.3in}\text{\ }
\end{table}

\begin{table}[H]
	\centering
	\caption{Ansari inequality}
	\begin{tabular}{rrrr}
		\toprule $N$ & $N^3+N^2$ & $(N+1)^3-\delta$ & holds? \\ 
		\midrule 2 & 12 & 26.875 & yes \\ 
		5 & 150 & 215.875 & yes \\ 
		
	\end{tabular}%
	\noindent \hspace{1em}\\[0pt]
	\noindent 
	\begin{tabular}{rrrr}
		10 & 1100 & 1330.875 & yes \\ 
		15 & 3600 & 4095.875 & yes \\ 
		
	\end{tabular}%
	\noindent \hspace{1em}\\[0pt]
	\noindent 
	\begin{tabular}{rrrr}
		20 & 8400 & 9260.875 & yes \\ 
		\bottomrule 
	\end{tabular}
\end{table}


\text{\ } \newline
\text{\ } \vspace{0.3in}\text{\ }
\begin{table}[H]
	\centering
	\caption{Bertrand's postulate}
	
	
	\text{\ } \newline
	\text{\ } \vspace{0.3in}\text{\ }
	
	\begin{tabular}{rr|rr}
		\toprule $N$ & prime in $(N,2N)$ & $N$ & prime in $(N,2N)$ \\ 
		\midrule 2 & 3 & 11 & 13 \\ 
		3 & 5 & 13 & 17 \\ 
		
	\end{tabular}%
	\noindent \hspace{1em}\\[0pt]
	\noindent 
	\begin{tabular}{rr|rr}
		5 & 7 & 17 & 19 \\ 
		7 & 11 & 19 & 23 \\ 
		
	\end{tabular}%
	\noindent \hspace{1em}\\[0pt]
	\noindent 
	\begin{tabular}{rr|rr}
		9 & 11 & 20 & 23 \\ 
		\noindent 
	\end{tabular}
	
\end{table}



\text{\ } \newline
\text{\ } \vspace{0.3in}\text{\ }
\begin{table}[H]
	\centering
	\caption{Wright sequence and iterated logarithms}
	\begin{tabular}{rrrr}
		\toprule $n$ & $q_n$ & $u_n = \log^{(n)} q_n$ & $v_n = \log^{(n)}(q_n+1)$ \\ 
		\midrule 1 & 2 & 1.000000 & 1.584963 \\ 
		2 & 5 & 1.215323 & 1.369919 \\ 
		
	\end{tabular}%
	\noindent \hspace{1em}\\[0pt]
	\noindent 
	\begin{tabular}{rrrr}
		3 & 37 & 1.251834 & 1.258213 \\ 
		\bottomrule 
	\end{tabular}
	
\end{table}
}
\text{\ } \newline
\text{\ } \vspace{0.3in}\text{\ }

{	\scriptsize
	
	\caption{ \large{Table of used Formal predicates}}
	
	\medskip
	
	% ===================== RCA_0 =====================
	\noindent \newline
	\noindent
	\begin{tabular}{|m{4cm}|m{5cm}|m{9cm}|}
		\hline
		\textbf{Object} & \textbf{Type} & \textbf{Detailed formulation} \\ \hline
		$\mathrm{RCA}_0$ & $\mathsf{System}$ & $\mathrm{RCA}_0 \iff \mathsf{Q} \wedge \mathrm{I}\Sigma^0_1 \wedge \Delta^0_1\text{-}\mathsf{CA}$ \\
	\end{tabular}%
	\noindent \newline
	\noindent
	\begin{tabular}{|m{4cm}|m{5cm}|m{9cm}|}
		\hline
		$\mathsf{Q}$ & $\mathsf{Prop}$ & $\mathsf{Q} \iff \bigl(\forall n,m,\ S(n)=S(m)\Rightarrow n=m\bigr) \wedge \bigl(\forall n,\ S(n)\neq 0\bigr) \wedge \bigl(\forall n,\ n+0=n\bigr) \wedge \bigl(\forall n,m,\ n+S(m)=S(n+m)\bigr) \wedge \bigl(\forall n,\ n\cdot 0=0\bigr) \wedge \bigl(\forall n,m,\ n\cdot S(m)=n\cdot m+n\bigr) \wedge \bigl(\forall n,\ \neg(n<0)\bigr) \wedge \bigl(\forall n,m,\ n<S(m)\Leftrightarrow(n<m\vee n=m)\bigr)$ \\
	\end{tabular}%
	\noindent \newline
	\noindent
	\begin{tabular}{|m{4cm}|m{5cm}|m{9cm}|}
		\hline
		$\mathrm{I}\Sigma^0_1$ & $\mathsf{Prop}$ & $\mathrm{I}\Sigma^0_1 \iff \forall k,\ \forall\varphi\in\Sigma^0_1,\ \Bigl[\bigl(\varphi(0,k)\wedge\forall n,(\varphi(n,k)\Rightarrow\varphi(n+1,k))\bigr)\Rightarrow\forall n,\ \varphi(n,k)\Bigr]$ \\ \hline
		$\Delta^0_1\text{-}\mathsf{CA}$ & $\mathsf{Prop}$ & $\Delta^0_1\text{-}\mathsf{CA} \iff \forall k,\ \forall\varphi,\psi\in\Sigma^0_1,\ \Bigl[\bigl(\forall n,\ \varphi(n,k)\Leftrightarrow\neg\psi(n,k)\bigr)\Rightarrow\exists X,\ \forall n,\ \bigl(n\in X\Leftrightarrow\varphi(n,k)\bigr)\Bigr]$ \\
	\end{tabular}%
	\noindent \newline
	\noindent
	\begin{tabular}{|m{4cm}|m{5cm}|m{9cm}|}
		\hline
		$\Sigma^0_0$ & $2^{\mathbb{N}}\to\mathsf{Prop}$ & $\Sigma^0_0 := \mathsf{Cl}_{\{\wedge,\vee,\neg,\Rightarrow,\exists n<t,\ \forall n<t\}}\bigl(\mathsf{At}_{\mathcal{L}_2}\bigr)$ \\ \hline
		$\Sigma^0_1$ & $2^{\mathbb{N}}\to\mathsf{Prop}$ & $\Sigma^0_1 := \bigl\{\,\exists m\,\psi(m,\cdot) \;:\; \psi\in\Sigma^0_0\,\bigr\}$ \\
	\end{tabular}%
	\noindent \newline
	\noindent
	\begin{tabular}{|m{4cm}|m{5cm}|m{9cm}|}
		\hline
		$\Pi^0_1$ & $2^{\mathbb{N}}\to\mathsf{Prop}$ & $\Pi^0_1 := \bigl\{\,\forall m\,\psi(m,\cdot) \;:\; \psi\in\Sigma^0_0\,\bigr\}$ \\ \hline
		$\Sigma^0_n$ & $\mathbb{N}_{\geq 1}\to\bigl(2^{\mathbb{N}}\to\mathsf{Prop}\bigr)$ & $\Sigma^0_n := \mathsf{Cl}_{\{\wedge,\vee,\neg,\Rightarrow,\exists k,\ \forall k\}}\bigl(\Pi^0_{n-1}\bigr)$ \\
	\end{tabular}%
	\noindent \newline
	\noindent
	\begin{tabular}{|m{4cm}|m{5cm}|m{9cm}|}
		\hline
		$\Pi^0_n$ & $\mathbb{N}_{\geq 1}\to\bigl(2^{\mathbb{N}}\to\mathsf{Prop}\bigr)$ & $\Pi^0_n := \bigl\{\,\neg\psi \;:\; \psi\in\Sigma^0_n\,\bigr\}$ \\ \hline
		$\mathsf{Arith}^0$ & $2^{\mathbb{N}}\to\mathsf{Prop}$ & $\mathsf{Arith}^0 := \bigcup_{n\in\mathbb{N}}\Sigma^0_n$ \\
	\end{tabular}%
	\noindent \newline
	\noindent
	\begin{tabular}{|m{4cm}|m{5cm}|m{9cm}|}
		\hline
		$\mathrm{ACA}_0$ & $\mathsf{System}$ & $\mathrm{ACA}_0 \iff \mathrm{RCA}_0 \wedge \mathsf{ACA}$ \\
	\end{tabular}%
	\noindent \newline
	\noindent
	\begin{tabular}{|m{4cm}|m{5cm}|m{9cm}|}
		\hline
		$\mathsf{ACA}$ & $\mathsf{Prop}$ & $\mathsf{ACA} \iff \forall k,\ \forall\varphi\in\mathsf{Arith}^0,\ \exists X,\ \forall n,\ \bigl(n\in X\Leftrightarrow\varphi(n,k)\bigr)$ \\ \hline
		$\mathsf{ACA}^{\mathsf{alt}}$ & $\mathsf{Prop}$ & $\mathsf{ACA}^{\mathsf{alt}} \iff \forall X,\ \exists Y,\ \forall n,\ \bigl(n\in Y\Leftrightarrow\Phi_n^X(n)\downarrow\bigr)$ \\
	\end{tabular}%
	\noindent \newline
	\noindent
	\begin{tabular}{|m{4cm}|m{5cm}|m{9cm}|}
		\hline
		$\mathsf{ACA}^{\mathsf{eq}}$ & $\mathsf{Prop}$ & $\mathsf{RCA}_0 \vdash \mathsf{ACA} \Leftrightarrow \mathsf{ACA}^{\mathsf{alt}} \Leftrightarrow \bigl(\forall X,\ \exists Y,\ Y=\mathsf{TJ}(X)\bigr)$ \\
	\end{tabular}%
	\noindent \newline
	\noindent
	\begin{tabular}{|m{4cm}|m{5cm}|m{9cm}|}
		\hline
		$\mathrm{WKL}_0$ & $\mathsf{System}$ & $\mathrm{WKL}_0 \iff \mathrm{RCA}_0 \wedge \mathsf{WKL}$ \\
	\end{tabular}%
	\noindent \newline
	\noindent
	\begin{tabular}{|m{4cm}|m{5cm}|m{9cm}|}
		\hline
		$\mathsf{WKL}$ & $\mathsf{Prop}$ & $\mathsf{WKL} \iff \forall T\subseteq\{0,1\}^{<\mathbb{N}},\ \bigl(\mathsf{InfBinTree}(T)\Rightarrow\exists f:\mathbb{N}\to\{0,1\},\ \forall n\in\mathbb{N},\ f\!\upharpoonright n\in T\bigr)$ \\ \hline
		$\mathsf{InfBinTree}$ & $2^{\{0,1\}^{<\mathbb{N}}}\to\mathsf{Prop}$ & $\forall T,\ \mathsf{InfBinTree}(T) \iff \bigl(\forall s\in T,\ \forall m<|s|,\ s\!\upharpoonright m\in T\bigr) \wedge \bigl(\forall n\in\mathbb{N},\ \exists s\in T,\ |s|=n\bigr)$ \\
	\end{tabular}
	\noindent \newline
	\noindent
	\begin{tabular}{|m{4cm}|m{5cm}|m{9cm}|}
		\hline
		$\mathsf{Type}$ & $\mathsf{Universe}$ &
		$\bigl(\forall n:\mathbb{N},\ \mathsf{Type}_n:\mathsf{Type}_{n+1}\bigr)
		\wedge\bigl(\forall n:\mathbb{N},\ \forall A:\mathsf{Type}_n,\ A:\mathsf{Type}_{n+1}\bigr)
		\wedge\mathsf{Type}=\bigcup_{n:\mathbb{N}}\mathsf{Type}_n$ \\ \hline
		$\mathsf{Prop}$ & $\mathsf{Type}$ &
		$\mathsf{Prop}:\mathsf{Type}\ \wedge\ \mathsf{LEM}\in\mathsf{Ax}\ \wedge\ \mathsf{LEM}\notin\mathsf{Def}(\mathsf{Prop})$ \\
	\end{tabular}%
	\noindent \newline
	\noindent
	\begin{tabular}{|m{4cm}|m{5cm}|m{9cm}|}
		\hline
		$\mathsf{LEM}$ & $\mathsf{Prop}$ &
		$\mathsf{LEM}:\iff\forall P:\mathsf{Prop},\ P\vee\neg P$ \\
	\end{tabular}%
	\noindent \newline
	\noindent
	\begin{tabular}{|m{4cm}|m{5cm}|m{9cm}|}
		\hline
		$\mathsf{Op}$ & $\mathsf{Type}$ &
		$\mathsf{Op}:=\bigsqcup_{A,B:\mathsf{Type}}(A\to B)$ \\ \hline
		$\mathsf{Rel}$ & $\mathsf{Type}$ &
		$\mathsf{Rel}:=\bigsqcup_{A,B:\mathsf{Type}}(A\times B\to\mathsf{Prop})$ \\
	\end{tabular}\noindent \newline
	\noindent
	\begin{tabular}{|m{4cm}|m{5cm}|m{9cm}|}
		\hline
		$\mathsf{dom}$ & $\mathsf{Op}\to\mathsf{Type}$ &
		$\mathsf{dom}(\langle A,B,f\rangle):=A$ \\ \hline
		$\mathsf{cod}$ & $\mathsf{Op}\to\mathsf{Type}$ &
		$\mathsf{cod}(\langle A,B,f\rangle):=B$ \\ 
	\end{tabular}\noindent \newline
	\noindent
	\begin{tabular}{|m{4cm}|m{5cm}|m{9cm}|}
		\hline
		$\mathcal{P}(T)$ & $\mathsf{Type}\to\mathsf{Type}$ &
		$\mathcal{P}(T):=T\to\mathsf{Prop}$ \\ \hline
		$\varnothing$ & $\mathcal{P}(T)$ &
		$\varnothing:=\lambda x:T,\ \bot$ \\ 
	\end{tabular}	\noindent \newline	\noindent
	\begin{tabular}{|m{4cm}|m{5cm}|m{9cm}|}
		\hline
		$=_X$ & $X\times X\to\mathsf{Prop}$ &
		$x=y:\iff\forall P:X\to\mathsf{Prop},\ P(x)\Rightarrow P(y)$ \\ \hline
		$\neq_X$ & $X\times X\to\mathsf{Prop}$ &
		$x\neq y:\iff\neg(x=y)$ \\ 
	\end{tabular}	\noindent \newline
	\noindent
	\begin{tabular}{|m{4cm}|m{5cm}|m{9cm}|}
		\hline
		$\in$ & $X\times\mathcal{P}(X)\to\mathsf{Prop}$ &
		$x\in S:\iff S(x)$ \\ \hline
		$\subseteq$ & $\mathcal{P}(X)\times\mathcal{P}(X)\to\mathsf{Prop}$ &
		$S\subseteq T:\iff\forall x:X,\ x\in S\Rightarrow x\in T$ \\ 
	\end{tabular}\noindent \newline
	\noindent
	\begin{tabular}{|m{4cm}|m{5cm}|m{9cm}|}
		\hline
		$<_F$ & $F\times F\to\mathsf{Prop}$ &
		$\mathsf{LinearOrder}(F,<)\wedge\mathsf{OrderCompatible}(F,+,\cdot,<)$
		(for a field $F$) \\ \hline
		$\leq_F$ & $F\times F\to\mathsf{Prop}$ &
		$x\leq y:\iff(x<y)\vee(x=y)$ \\ 
	\end{tabular}
	\noindent \newline
	\noindent
	\begin{tabular}{|m{4cm}|m{5cm}|m{9cm}|}
		\hline
		$\mathrm{succ}$ & $\mathbb{N}\to\mathbb{N}$ &
		$\forall n:\mathbb{N},\ \mathrm{succ}(n)=n+1$ \\ \hline
		$\mathrm{char}$ & $(\mathsf{Type},\mathsf{Op},\mathsf{Op},T,T)\to\mathbb{N}$ &
		$\bigl(\mathrm{char}(R,+,\cdot,0,1)=0 \iff \forall k\in\mathbb{N}_{\geq 1},\ k\cdot 1_R\neq 0_R\bigr)
		\wedge\bigl(\mathrm{char}(R,+,\cdot,0,1)=n\wedge n\in\mathbb{N}_{\geq 1} \iff n\cdot 1_R=0_R\wedge\forall k\in\mathbb{N}_{\geq 1},\ k<n\Rightarrow k\cdot 1_R\neq 0_R\bigr)$ \\
	\end{tabular}%
	\noindent \newline
	\noindent
	\begin{tabular}{|m{4cm}|m{5cm}|m{9cm}|}
		\hline
		$\mathrm{deg}$ & $R[x]\to\mathbb{N}\cup\{\bot\}$ &
		$\mathrm{deg}(0):=\bot$;\quad
		$\mathrm{deg}\bigl(\textstyle\sum_{k=0}^{m}a_kx^k\bigr)=m:\iff a_m\neq 0$ \\ \hline
		$\mathsf{AddAssoc}(R,+)$ & $(\mathsf{Type},\mathsf{Op})\to\mathsf{Prop}$ &
		$\forall x,y,z\in R,\ (x+y)+z=x+(y+z)$ \\ 
	\end{tabular}%	%==================== Ring axioms ====================
	\noindent \newline
	\noindent
	\begin{tabular}{|m{4cm}|m{5cm}|m{9cm}|}
		\hline
		$\mathsf{AddComm}(R,+)$ & $(\mathsf{Type},\mathsf{Op})\to\mathsf{Prop}$ &
		$\forall x,y\in R,\ x+y=y+x$ \\ \hline
		$\mathsf{AddId}(R,+,0)$ & $(\mathsf{Type},\mathsf{Op},T)\to\mathsf{Prop}$ &
		$\forall x\in R,\ x+0=x$ \\ 
	\end{tabular}
	\noindent \newline
	\noindent
	\begin{tabular}{|m{4cm}|m{5cm}|m{9cm}|}
		\hline
		$\mathsf{AddInv}(R,+,0)$ & $(\mathsf{Type},\mathsf{Op},T)\to\mathsf{Prop}$ &
		$\forall x\in R,\ \exists w\in R,\ x+w=0$ \\ \hline
		$\mathsf{MulAssoc}(R,\cdot)$ & $(\mathsf{Type},\mathsf{Op})\to\mathsf{Prop}$ &
		$\forall x,y,z\in R,\ (x\cdot y)\cdot z=x\cdot(y\cdot z)$ \\ 
	\end{tabular}
	\noindent \newline
	\noindent
	\begin{tabular}{|m{4cm}|m{5cm}|m{9cm}|}
		\hline
		$\mathsf{MulComm}(R,\cdot)$ & $(\mathsf{Type},\mathsf{Op})\to\mathsf{Prop}$ &
		$\forall x,y\in R,\ x\cdot y=y\cdot x$ \\ \hline
		$\mathsf{MulId}(R,\cdot,1)$ & $(\mathsf{Type},\mathsf{Op},T)\to\mathsf{Prop}$ &
		$\forall x\in R,\ x\cdot 1=x$ \\ 
	\end{tabular}
	\noindent \newline
	\noindent
	\begin{tabular}{|m{4cm}|m{5cm}|m{9cm}|}
		\hline
		$\mathsf{MulInv}(R,\cdot,0,1)$ & $(\mathsf{Type},\mathsf{Op},T,T)\to\mathsf{Prop}$ &
		$\forall a\in R,\ a\neq 0\Rightarrow\exists a^{-1}\in R,\ a\cdot a^{-1}=1$ \\ \hline
		$\mathsf{Distrib}(R,+,\cdot)$ & $(\mathsf{Type},\mathsf{Op},\mathsf{Op})\to\mathsf{Prop}$ &
		$\forall x,y,z\in R,\ x\cdot(y+z)=x\cdot y+x\cdot z$ \\ 
	\end{tabular}
	\noindent \newline
	\noindent
	\begin{tabular}{|m{4cm}|m{5cm}|m{9cm}|}
		\hline
		$\mathsf{Nontrivial}(R,0,1)$ & $(\mathsf{Type},T,T)\to\mathsf{Prop}$ &
		$0\neq 1$ \\ \hline
		$\mathsf{NoZeroDiv}(R,\cdot,0)$ & $(\mathsf{Type},\mathsf{Op},T)\to\mathsf{Prop}$ &
		$\forall a,b\in R,\ a\cdot b=0\Rightarrow a=0\vee b=0$ \\
	\end{tabular}%
	%==================== Rings ====================
	\noindent \newline
	\noindent
	\begin{tabular}{|m{4cm}|m{5cm}|m{9cm}|}
		\hline
		$\mathsf{CommutativeRing}$ & $(\mathsf{Type},\mathsf{Op},\mathsf{Op},T,T)\to\mathsf{Prop}$ &
		$\mathsf{AddAssoc}(R,+)\wedge\mathsf{AddComm}(R,+)\wedge\mathsf{AddId}(R,+,0)
		\wedge\mathsf{AddInv}(R,+,0)\wedge\mathsf{MulAssoc}(R,\cdot)
		\wedge\mathsf{MulComm}(R,\cdot)\wedge\mathsf{MulId}(R,\cdot,1)
		\wedge\mathsf{Distrib}(R,+,\cdot)\wedge\mathsf{Nontrivial}(R,0,1)$ \\ \hline
		$\mathsf{IntegralDomain}$ & $(\mathsf{Type},\mathsf{Op},\mathsf{Op},T,T)\to\mathsf{Prop}$ &
		$\mathsf{CommutativeRing}(R,+,\cdot,0,1)\wedge\mathsf{NoZeroDiv}(R,\cdot,0)$ \\ 
	\end{tabular}
	\noindent \newline
	\noindent
	\begin{tabular}{|m{4cm}|m{5cm}|m{9cm}|}
		\hline
		$\mathsf{Field}$ & $(\mathsf{Type},\mathsf{Op},\mathsf{Op},T,T)\to\mathsf{Prop}$ &
		$\mathsf{CommutativeRing}(F,+,\cdot,0,1)\wedge\mathsf{MulInv}(F,\cdot,0,1)$ \\ \hline
		$\mathsf{HomAdd}(\varphi)$ & $\mathsf{Op}\to\mathsf{Prop}$ &
		$\forall x,y\in\mathsf{dom}(\varphi),\ \varphi(x+y)=\varphi(x)+\varphi(y)$ \\ 
	\end{tabular}
	\noindent \newline
	\noindent
	\begin{tabular}{|m{4cm}|m{5cm}|m{9cm}|}
		\hline
		$\mathsf{HomMul}(\varphi)$ & $\mathsf{Op}\to\mathsf{Prop}$ &
		$\forall x,y\in\mathsf{dom}(\varphi),\ \varphi(x\cdot y)=\varphi(x)\cdot\varphi(y)$ \\ \hline
		$\mathsf{HomOne}(\varphi)$ & $\mathsf{Op}\to\mathsf{Prop}$ &
		$\varphi(1_{\mathsf{dom}(\varphi)})=1_{\mathsf{cod}(\varphi)}$ \\ 
	\end{tabular}
	\noindent \newline
	\noindent
	\begin{tabular}{|m{4cm}|m{5cm}|m{9cm}|}
		\hline
		$\mathsf{RingHom}(\varphi)$ & $\mathsf{Op}\to\mathsf{Prop}$ &
		$\mathsf{HomAdd}(\varphi)\wedge\mathsf{HomMul}(\varphi)\wedge\mathsf{HomOne}(\varphi)$ \\ \hline
		$\mathsf{ClosedUnder}(T,S)$ & $(\mathsf{Type},\mathcal{P}(\mathsf{Op}))\to\mathsf{Prop}$ &
		$\forall op\in S,\ \forall x\in\mathsf{dom}(op)\cap T^{\mathrm{arity}(op)},\
		op(x)\in T$ (where $\mathrm{arity}(op)$ is the number of components of
		$\mathsf{dom}(op)$) \\ 
	\end{tabular}%
	%==================== Peano ====================
	\noindent \newline
	\noindent
	\begin{tabular}{|m{4cm}|m{5cm}|m{9cm}|}
		\hline
		$\mathsf{ZeroMem}(N,0)$ & $(\mathsf{Type},T)\to\mathsf{Prop}$ &
		$0\in N$ \\ \hline
		$\mathsf{SuccClosed}(N,s)$ & $(\mathsf{Type},T\to T)\to\mathsf{Prop}$ &
		$\forall n\in N,\ s(n)\in N$ \\ 
	\end{tabular}
	\noindent \newline
	\noindent
	\begin{tabular}{|m{4cm}|m{5cm}|m{9cm}|}
		\hline
		$\mathsf{SuccNeqZero}(N,s,0)$ & $(\mathsf{Type},T\to T,T)\to\mathsf{Prop}$ &
		$\forall n\in N,\ s(n)\neq 0$ \\ \hline
		$\mathsf{SuccInj}(N,s)$ & $(\mathsf{Type},T\to T)\to\mathsf{Prop}$ &
		$\forall m,n\in N,\ s(m)=s(n)\Rightarrow m=n$ \\ 
	\end{tabular}
	\noindent \newline
	\noindent
	\begin{tabular}{|m{4cm}|m{5cm}|m{9cm}|}
		\hline
		$\mathsf{Induction}(N,s,0)$ & $(\mathsf{Type},T\to T,T)\to\mathsf{Prop}$ &
		$\forall S\subseteq N,\ \bigl(0\in S\wedge(\forall n\in S,\ s(n)\in S)\bigr)
		\Rightarrow S=N$ \\ \hline
		$\mathsf{Peano}$ & $(\mathsf{Type},T,T\to T)\to\mathsf{Prop}$ &
		$\mathsf{ZeroMem}(N,0)\wedge\mathsf{SuccClosed}(N,s)
		\wedge\mathsf{SuccNeqZero}(N,s,0)\wedge\mathsf{SuccInj}(N,s)
		\wedge\mathsf{Induction}(N,s,0)$ \\ 
	\end{tabular}%
	%==================== Initial ====================
	\noindent \newline
	\noindent
	\begin{tabular}{|m{4cm}|m{5cm}|m{9cm}|}
		\hline
		$\mathsf{Initial}(F)$ & $\mathsf{Type}\to\mathsf{Prop}$ &
		$\forall R,\ \mathsf{CommutativeRing}(R)\Rightarrow
		\exists!\,\varphi:F\to R,\ \mathsf{RingHom}(\varphi)$ \\ \hline
		$\mathsf{Irrefl}(F,<)$ & $(\mathsf{Type},\mathsf{Rel})\to\mathsf{Prop}$ &
		$\forall x\in F,\ \neg(x<x)$ \\ 
	\end{tabular}%
	%==================== Order ====================
	\noindent \newline
	\noindent
	\begin{tabular}{|m{4cm}|m{5cm}|m{9cm}|}
		\hline
		$\mathsf{Trans}(F,<)$ & $(\mathsf{Type},\mathsf{Rel})\to\mathsf{Prop}$ &
		$\forall x,y,z\in F,\ x<y\wedge y<z\Rightarrow x<z$ \\ \hline
		$\mathsf{Total}(F,<)$ & $(\mathsf{Type},\mathsf{Rel})\to\mathsf{Prop}$ &
		$\forall x,y\in F,\ x<y\vee x=y\vee y<x$ \\ 
	\end{tabular}
	\noindent \newline
	\noindent
	\begin{tabular}{|m{4cm}|m{5cm}|m{9cm}|}
		\hline
		$\mathsf{OrdAdd}(F,+,<)$ & $(\mathsf{Type},\mathsf{Op},\mathsf{Rel})\to\mathsf{Prop}$ &
		$\forall a,b,c\in F,\ a<b\Rightarrow a+c<b+c$ \\ \hline
		$\mathsf{OrdMul}(F,\cdot,<)$ & $(\mathsf{Type},\mathsf{Op},\mathsf{Rel})\to\mathsf{Prop}$ &
		$\forall a,b,c\in F,\ a<b\wedge 0<c\Rightarrow a\cdot c<b\cdot c$ \\ 
	\end{tabular}
	\noindent \newline
	\noindent
	\begin{tabular}{|m{4cm}|m{5cm}|m{9cm}|}
		\hline
		$\mathsf{LinearOrder}(F,<)$ & $(\mathsf{Type},\mathsf{Rel})\to\mathsf{Prop}$ &
		$\mathsf{Irrefl}(F,<)\wedge\mathsf{Trans}(F,<)\wedge\mathsf{Total}(F,<)$ \\ \hline
		$\mathsf{OrderCompatible}(F,+,\cdot,<)$ & $(\mathsf{Type},\mathsf{Op},\mathsf{Op},\mathsf{Rel})\to\mathsf{Prop}$ &
		$\mathsf{OrdAdd}(F,+,<)\wedge\mathsf{OrdMul}(F,\cdot,<)$ \\ 
	\end{tabular}%
	%==================== Dedekind ====================
	\noindent \newline
	\noindent
	\begin{tabular}{|m{4cm}|m{5cm}|m{9cm}|}
		\hline
		$\mathsf{BddAbove}(F,S,\leq)$ & $(\mathsf{Type},\mathcal{P}(F),\mathsf{Rel})\to\mathsf{Prop}$ &
		$\exists u\in F,\ \forall x\in S,\ x\leq u$ \\ \hline
		$\mathsf{HasSup}(F,S,\leq)$ & $(\mathsf{Type},\mathcal{P}(F),\mathsf{Rel})\to\mathsf{Prop}$ &
		$\exists s\in F,\ (\forall x\in S,\ x\leq s)\wedge
		(\forall u\in F,\ (\forall x\in S,\ x\leq u)\Rightarrow s\leq u)$ \\ 
	\end{tabular}
	\noindent \newline
	\noindent
	\begin{tabular}{|m{4cm}|m{5cm}|m{9cm}|}
		\hline
		$\mathsf{DedekindComplete}(F,<)$ & $(\mathsf{Type},\mathsf{Rel})\to\mathsf{Prop}$ &
		$\forall S\subseteq F,\ \bigl(S\neq\varnothing\wedge
		\mathsf{BddAbove}(F,S,\leq)\bigr)\Rightarrow\mathsf{HasSup}(F,S,\leq)$ \\ \hline
		$\mathsf{ContainsBase}(F,C)$ & $(\mathsf{Type},\mathsf{Type})\to\mathsf{Prop}$ &
		$F\subseteq C$ \\ 
	\end{tabular}%
	%%%==================== Algebraic closure ====================
	\noindent \newline
	\noindent
	\begin{tabular}{|m{4cm}|m{5cm}|m{9cm}|}
		\hline
		$\mathsf{AlgClosed}(C)$ & $\mathsf{Type}\to\mathsf{Prop}$ &
		$\forall p\in C[x],\ \mathrm{deg}(p)\geq 1\Rightarrow\exists r\in C,\ p(r)=0$ \\ \hline
		$\mathsf{AlgebraicClosure}(C,F)$ & $(\mathsf{Type},\mathsf{Type})\to\mathsf{Prop}$ &
		$\mathsf{Field}(C)\wedge\mathsf{ContainsBase}(F,C)\wedge\mathsf{AlgClosed}(C)$ \\ 
	\end{tabular}%
	%==================== Measure ====================
	\noindent \newline
	\noindent
	\begin{tabular}{|m{4cm}|m{5cm}|m{9cm}|}
		\hline
		$\mathbb{R}_{\ge 0}^{\mathsf{ext}}$ & $\mathsf{Type}$ &
		$\mathbb{R}_{\ge 0}^{\mathsf{ext}} := \mathbb{R}_{\ge 0}\cup\{\mathsf{Top}\}$ \\ \hline
		$\mathsf{TransInv}(\mu)$ & $(\mathcal{P}(\mathbb{R})\to\mathbb{R}_{\ge 0}^{\mathsf{ext}})\to\mathsf{Prop}$ &
		$\forall a:\mathbb{R},\ \forall S:\mathcal{P}(\mathbb{R}),\ \mu(S+a)=\mu(S)$ \\ \hline
		$\mathsf{CountAdd}(\mu)$ & $(\mathcal{P}(\mathbb{R})\to\mathbb{R}_{\ge 0}^{\mathsf{ext}})\to\mathsf{Prop}$ &
		$\forall(S_n)_{n:\mathbb{N}}\subseteq\mathcal{P}(\mathbb{R}),\
		(\forall m\neq n,\ S_m\cap S_n=\varnothing)\Rightarrow
		\mu\bigl(\bigcup_n S_n\bigr)=\sum_n\mu(S_n)$ \\
	\end{tabular}%
	\noindent \newline
	\noindent
	\begin{tabular}{|m{4cm}|m{5cm}|m{9cm}|}
		\hline
		$\sum_n$ & $(\mathbb{N}\to\mathbb{R}_{\ge 0}^{\mathsf{ext}})\to\mathbb{R}_{\ge 0}^{\mathsf{ext}}$ &
		$\sum_n a_n := \sup\Bigl\{\sum_{n\in F}a_n \;:\; F\in\mathcal{P}_{\mathsf{fin}}(\mathbb{N})\Bigr\}$ \\ \hline
		$\mathcal{P}_{\mathsf{fin}}(\mathbb{N})$ & $\mathcal{P}(\mathcal{P}(\mathbb{N}))$ &
		$F\in\mathcal{P}_{\mathsf{fin}}(\mathbb{N}) \iff \exists m\in\mathbb{N},\ \exists f:\{1,\dots,m\}\to\mathbb{N},\ F=f(\{1,\dots,m\})$ \\
	\end{tabular}%
	\noindent \newline
	\noindent
	\noindent \newline
	\noindent
	\begin{tabular}{|m{4cm}|m{5cm}|m{9cm}|}
		\hline
		\textbf{Object} & \textbf{Type} & \textbf{Detailed formulation} \\ \hline
		$\mathsf{MeasureEmptyZero}(\mu)$ & $(\mathcal{P}(\mathbb{R})\to\mathbb{R}_{\ge 0}^{\mathsf{ext}})\to\mathsf{Prop}$ &
		$\mu(\varnothing)=0$ \\
	\end{tabular}%
	\noindent \newline
	\noindent
	\begin{tabular}{|m{4cm}|m{5cm}|m{9cm}|}
		\hline
		$\mathsf{MeasureUnitInterval}(\mu)$ & $(\mathcal{P}(\mathbb{R})\to\mathbb{R}_{\ge 0}^{\mathsf{ext}})\to\mathsf{Prop}$ &
		$\mu([0,1])=1$ \\
	\end{tabular}%
	\noindent \newline
	\noindent
	\begin{tabular}{|m{4cm}|m{5cm}|m{9cm}|}
		\hline
		$\mathsf{LebesgueMeasure}(\mu)$ & $((\mathbb{R}\to\mathsf{Prop})\to\mathbb{R}_{\ge 0}^{\mathsf{ext}})\to\mathsf{Prop}$ &
		$\mathsf{TransInv}(\mu)\wedge\mathsf{CountAdd}(\mu)
		\wedge\mathsf{MeasureEmptyZero}(\mu)\wedge\mathsf{MeasureUnitInterval}(\mu)$ \\ \hline
		$\mathsf{IsOpen}(U)$ & $(\mathbb{R}\to\mathsf{Prop})\to\mathsf{Prop}$ &
		$\forall x\in U,\ \exists\varepsilon>0,\ \forall y:\mathbb{R},\
		|y-x|<\varepsilon\Rightarrow y\in U$ \\
	\end{tabular}%
	%==================== Topology ====================
	\noindent \newline
	\noindent
	\begin{tabular}{|m{4cm}|m{5cm}|m{9cm}|}
		\hline
		$\mathsf{IsClosed}(F)$ & $(\mathbb{R}\to\mathsf{Prop})\to\mathsf{Prop}$ &
		$\mathsf{IsOpen}(\lambda x,\ \neg F(x))$ \\ \hline
		$\mathbb{N}$ & $\mathsf{Type}$ &
		$\mathsf{Peano}(\mathbb{N},0,\mathrm{succ})$ \\ 
	\end{tabular}%
	%==================== Number systems ====================
	\noindent \newline
	\noindent
	\begin{tabular}{|m{4cm}|m{5cm}|m{9cm}|}
		\hline
		$\mathbb{Z}$ & $\mathsf{Type}$ &
		$\mathsf{IntegralDomain}(\mathbb{Z},+,\cdot,0,1)\wedge\mathsf{Initial}(\mathbb{Z})
		\wedge\mathrm{char}(\mathbb{Z},+,\cdot,0,1)=0$ \\ \hline
		$\mathbb{R}$ & $\mathsf{Type}$ &
		$\mathsf{Field}(\mathbb{R},+,\cdot,0,1)\wedge\mathsf{LinearOrder}(\mathbb{R},<)
		\wedge\mathsf{OrderCompatible}(\mathbb{R},+,\cdot,<)
		\wedge\mathsf{DedekindComplete}(\mathbb{R},<)$ \\
	\end{tabular}
	\noindent \newline
	\noindent
	\begin{tabular}{|m{4cm}|m{5cm}|m{9cm}|}
		\hline
		$\mathbb{C}$ & $\mathsf{Type}$ &
		$\mathsf{AlgebraicClosure}(\mathbb{C},\mathbb{R})\wedge\mathbb{R}\subseteq\mathbb{C}$ \\ \hline
		$\exp$ & $\mathbb{R}\to\mathbb{R}$ &
		$\exp(0)=1\wedge(\forall x,\ \exp'(x)=\exp(x))
		\wedge(\forall x,y,\ \exp(x+y)=\exp(x)\exp(y))$
		(differentiability is implicit in the notation $\exp'$) \\ 
	\end{tabular}%
	%==================== Elementary functions ====================
	\noindent \newline
	\noindent
	\noindent \newline
	\noindent
	\begin{tabular}{|m{4cm}|m{5cm}|m{9cm}|}
		\hline
		$\log$ & $\mathbb{R}_{>0}\to\mathbb{R}$ &
		$(\forall x,\ \log(\exp(x))=x)\wedge(\forall y>0,\ \exp(\log(y))=y)$ \\ \hline
		$\zeta$ & $\mathbb{C}\to\mathbb{C}$ &
		$\forall s:\mathbb{C},\ \Re(s)>1\Rightarrow\zeta(s)=\sum_{n\in\mathbb{N}_{\geq 1}}n^{-s}$ \\
	\end{tabular}%
	\noindent \newline
	\noindent
	\begin{tabular}{|m{4cm}|m{5cm}|m{9cm}|}
		\hline
		$\exp_2$ & $\mathbb{R}\to\mathbb{R}$ &
		$\exp_2(x):=\exp(x\log 2)$ \\ \hline
		$\log_2$ & $\mathbb{R}_{>0}\to\mathbb{R}$ &
		$\log_2(x):=\log(x)/\log 2$ \\ 
	\end{tabular}
	\noindent \newline
	\noindent
	\begin{tabular}{|m{4cm}|m{5cm}|m{9cm}|}
		\hline
		$\log_c$ & $\mathbb{R}_{>0}\to\mathbb{R}$ &
		$\log_c(x):=\log(x)/\log c$ \quad (for $c>0$, $c\neq 1$) \\ \hline
		$\Re$ & $\mathbb{C}\to\mathbb{R}$ &
		$\forall a,b\in\mathbb{R},\ \Re(a+bi)=a$ \\
	\end{tabular}%
	%==================== Calculus ====================
	\noindent \newline
	\noindent
	\begin{tabular}{|m{4cm}|m{5cm}|m{9cm}|}
		\hline
		$\mathsf{Continuous}(f)$ & $(\mathbb{R}\to\mathbb{R})\to\mathsf{Prop}$ &
		$\forall x,\ \forall\varepsilon>0,\ \exists\delta>0,\ \forall y,\
		|y-x|<\delta\Rightarrow|f(y)-f(x)|<\varepsilon$ \\ \hline
		$\mathsf{EventuallyIncreasing}(f)$ & $(\mathbb{R}\to\mathbb{R})\to\mathsf{Prop}$ &
		$\exists N,\ \forall x,y,\ N<x<y\Rightarrow f(x)<f(y)$ \\ 
		\hline
		$\mathsf{Bijective}(f)$ & $\forall\alpha,\beta:\mathsf{Type},\ (\alpha\to\beta)\to\mathsf{Prop}$ &
		$(\forall x,y,\ f(x)=f(y)\Rightarrow x=y)\wedge(\forall y,\ \exists x,\ f(x)=y)$ \\ 
	\end{tabular}%
	\noindent \newline
	\noindent 
	\begin{tabular}{|m{4cm}|m{5cm}|m{9cm}|}
		\hline
		$\mu$ & $(\mathbb{R}\to\mathsf{Prop})\to\mathbb{R}_{\ge 0}^{\mathsf{ext}}$ &
		$\mathsf{LebesgueMeasure}(\mu)$ \\ \hline
		$\mathrm{int}$ & $(\mathbb{R}\to\mathsf{Prop})\to(\mathbb{R}\to\mathsf{Prop})$ &
		$\mathrm{int}(S)\subseteq S\wedge\mathsf{IsOpen}(\mathrm{int}(S))
		\wedge\forall U,\ (\mathsf{IsOpen}(U)\wedge U\subseteq S)\Rightarrow
		U\subseteq\mathrm{int}(S)$ \\
	\end{tabular}
	\noindent \newline
	\noindent
	\begin{tabular}{|m{4cm}|m{5cm}|m{9cm}|}
		\hline
		$\mathrm{cl}$ & $(\mathbb{R}\to\mathsf{Prop})\to(\mathbb{R}\to\mathsf{Prop})$ &
		$S\subseteq\mathrm{cl}(S)\wedge\mathsf{IsClosed}(\mathrm{cl}(S))
		\wedge\forall F,\ (\mathsf{IsClosed}(F)\wedge S\subseteq F)\Rightarrow
		\mathrm{cl}(S)\subseteq F$ \\ \hline
		$\mathrm{floor}$ & $\mathbb{R}\to\mathbb{Z}$ &
		$\forall x:\mathbb{R},\ \forall n:\mathbb{Z},\ \mathrm{floor}(x)=n
		\iff(n\leq x<n+1)$ \\ 
	\end{tabular}%
	%==================== Primes ====================
	\noindent \newline
	\noindent
	\begin{tabular}{|m{4cm}|m{5cm}|m{9cm}|}
		\hline
		$\mathrm{Prime}$ & $\mathbb{N}\to\mathsf{Prop}$ &
		$(p\geq 2)\wedge(\forall d:\mathbb{N},\ (d\geq 2\wedge d\mid p)
		\Rightarrow d=p)$ \\ \hline
		$\mathrm{PSeq}$ & $(\mathbb{N}\to\mathbb{N})\to\mathsf{Prop}$ &
		$(p_1=2)\wedge(\forall n\geq 1,\ \mathrm{Prime}(p_n))
		\wedge(\forall n\geq 1,\ p_n<p_{n+1})$ \\ 
	\end{tabular}%
	%==================== Mills ====================
	\noindent \newline
	\noindent
	\begin{tabular}{|m{4cm}|m{5cm}|m{9cm}|}
		\hline
		$\mathrm{IsMills}$ & $\mathbb{R}\to\mathsf{Prop}$ &
		$(A>1)\wedge(\forall n\geq 1,\ \mathrm{Prime}(\mathrm{floor}(A^{3^n})))$ \\ \hline
		$\mathrm{IsMinMills}$ & $\mathbb{R}\to\mathsf{Prop}$ &
		$\mathrm{IsMills}(A)\wedge(\forall B:\mathbb{R},\ \mathrm{IsMills}(B)
		\Rightarrow A\leq B)$ \\ 
	\end{tabular}
	\noindent \newline
	\noindent
	\begin{tabular}{|m{4cm}|m{5cm}|m{9cm}|}
		\hline
		$\mathrm{MillsRec}$ & $(\mathbb{N}\to\mathbb{N})\to\mathsf{Prop}$ &
		$\forall n\geq 1,\ q_n^3<q_{n+1}<(q_n+1)^3-1$ \\ \hline
		$\mathrm{UV}$ & $(\mathbb{N}\to\mathbb{N})\times(\mathbb{N}\to\mathbb{R})^2\to\mathsf{Prop}$ &
		$\forall n\geq 1,\ (u_n=q_n^{3^{-n}})\wedge(v_n=(q_n+1)^{3^{-n}})$ \\ 
	\end{tabular}%
	%==================== Kuipers / Ansari / Wright / Niven ====================
	\noindent \newline
	\noindent
	\begin{tabular}{|m{4cm}|m{5cm}|m{9cm}|}
		\hline
		$\mathrm{IsKuipers}$ & $\mathbb{N}_{\geq 3}\times\mathbb{R}\to\mathsf{Prop}$ &
		$(A>1)\wedge(\forall n\geq 1,\ \mathrm{Prime}(\mathrm{floor}(A^{c^n})))$ \\ \hline
		$\mathrm{IsAnsari}$ & $\mathbb{R}_{>8/3}\times\mathbb{R}\to\mathsf{Prop}$ &
		$(A>1)\wedge(\forall n\geq 1,\ \mathrm{Prime}(\mathrm{floor}(A^{c^n})))$ \\ 
	\end{tabular}
	\noindent \newline
	\noindent
	\begin{tabular}{|m{4cm}|m{5cm}|m{9cm}|}
		\hline
		$\mathrm{IsWright}$ & $\mathbb{R}\to\mathsf{Prop}$ &
		$(A>1)\wedge(\forall n\geq 1,\
		\mathrm{Prime}(\mathrm{floor}(\mathrm{exp}_2^{[n]}(A))))$ \\ \hline
		$\mathrm{IsNiven}$ & $\mathbb{R}_{>1}\times\mathbb{R}\to\mathsf{Prop}$ &
		$(c>1)\wedge(\theta>0)\wedge(\forall n\geq 1,\
		\mathrm{Prime}(\mathrm{floor}(c^{\theta^n})))$ \\ 
	\end{tabular}%
	%==================== Ore ====================
	\noindent \newline
	\noindent
	\begin{tabular}{|m{4cm}|m{5cm}|m{9cm}|}
		\hline
		$\mathrm{IsOreSel}$ & $(\mathbb{R}\to\mathbb{R})\to\mathsf{Prop}$ &
		$\mathsf{Continuous}(f)\wedge\mathsf{EventuallyIncreasing}(f)$ \\ \hline
		$\mathrm{OreBound}$ & $(\mathbb{R}\to\mathbb{R})\times(\mathbb{N}\to\mathbb{N})
		\times(\mathbb{N}\to\mathbb{R}_{\geq 0})^2\to\mathsf{Prop}$ &
		$\forall n\geq 1,\ f(q_n-\delta_n)<q_{n+1}-\delta_{n+1}\wedge
		q_{n+1}+\epsilon_{n+1}<f(q_n+\epsilon_n)$ \\ 
	\end{tabular}%
	%==================== Prime gaps / null ====================
	\noindent \newline
	\noindent
	\begin{tabular}{|m{4cm}|m{5cm}|m{9cm}|}
		\hline
		$\mathrm{PolyGap}$ & $(0,1)\times\mathbb{R}_{>0}\to\mathsf{Prop}$ &
		$\forall n\geq 1,\ p_{n+1}-p_n<k\,p_n^{\alpha}$ \\ \hline
		$\mathsf{IsNull}$ & $(\mathbb{R}\to\mathsf{Prop})\to\mathsf{Prop}$ &
		$\mu(S)=0$ \\ 
	\end{tabular}%
	%==================== Density / cardinality ====================
	\noindent \newline
	\noindent
	\begin{tabular}{|m{4cm}|m{5cm}|m{9cm}|}
		\hline
		$\mathrm{NowhereDense}$ & $(\mathbb{R}\to\mathsf{Prop})\to\mathsf{Prop}$ &
		$\mathrm{int}(\mathrm{cl}(S))=\varnothing$ \\ \hline
		$\mathrm{ContinuumCard}$ & $(\mathbb{R}\to\mathsf{Prop})\to\mathsf{Prop}$ &
		$\exists f:\mathbb{R}\to(\Sigma\,x:\mathbb{R},\ P(x)),\ \mathsf{Bijective}(f)$ \\ 
	\end{tabular}%
	%==================== Irrationality / RH ====================
	\noindent \newline
	\noindent
	\begin{tabular}{|m{4cm}|m{5cm}|m{9cm}|}
		\hline
		$\mathrm{Irrational}$ & $\mathbb{R}\to\mathsf{Prop}$ &
		$\neg(\exists p:\mathbb{Z},\ \exists q:\mathbb{N}_{\geq 1},\ A=p/q)$ \\ \hline
		$\mathrm{RH}$ & $\mathsf{Prop}$ &
		$\forall s:\mathbb{C},\ (\zeta(s)=0\wedge\Re(s)>1/2)\Rightarrow\Re(s)=1/2$ \\ 
	\end{tabular}%
	%==================== Theta constants ====================
	\noindent \newline
	\noindent
	\begin{tabular}{|m{4cm}|m{5cm}|m{9cm}|}
		\hline
		$\mathrm{ThetaMills}$ & $\mathbb{R}\to\mathsf{Prop}$ &
		$\mathrm{IsMills}(\theta)$ \\ \hline
		$\mathrm{ThetaAnsari}$ & $\mathbb{R}\to\mathsf{Prop}$ &
		$\exists c:\mathbb{R}_{>8/3},\ \mathrm{IsAnsari}(c,\theta)$ \\ 
	\end{tabular}
	\noindent \newline
	\noindent
	\begin{tabular}{|m{4cm}|m{5cm}|m{9cm}|}
		\hline
		$\mathrm{ThetaWright}$ & $\mathbb{R}\to\mathsf{Prop}$ &
		$\mathrm{IsWright}(\theta)$ \\ \hline
		$\mathrm{ThetaNiven}$ & $\mathbb{R}\to\mathsf{Prop}$ &
		$\exists c:\mathbb{R}_{>1},\ \mathrm{IsNiven}(c,\theta)$ \\ 
	\end{tabular}%
	%==================== Iterated exp / log (base 2) ====================
	\noindent \newline
	\noindent
	\begin{tabular}{|m{4cm}|m{5cm}|m{9cm}|}
		\hline
		$\mathrm{exp}_2^{[n]}$ & $\mathbb{R}\to\mathbb{R}$ &
		$(\forall x,\ \mathrm{exp}_2^{[0]}(x)=x)\wedge(\forall n,\
		\mathrm{exp}_2^{[n+1]}=\mathrm{exp}_2\circ\mathrm{exp}_2^{[n]})$ \\ \hline
		$\mathrm{log}_2^{[n]}$ & $\mathbb{R}_{>0}\to\mathbb{R}$ &
		$(\forall x>0,\ \mathrm{log}_2^{[0]}(x)=x)\wedge(\forall n,\
		\mathrm{log}_2^{[n+1]}=\mathrm{log}_2^{[n]}\circ\mathrm{log}_2)$ \\ \hline
	\end{tabular}
}
\newpage
\subsection{Figures}\label{sec:appendix-figures}
\noindent \vspace{-0.3in}
\begin{figure}[H]
	\centering
	\begin{tikzpicture}[scale=0.9]
	\draw[->] (0,0) -- (10,0) node[right] {$n$};
	\draw[->] (0,0) -- (0,4) node[above] {};
	\draw[thick,blue] (1,1.0) -- (2,1.5) -- (3,2.0) -- (4,2.3) -- (5,2.45) -- (6,2.52) -- (7,2.55) -- (8,2.57);
	\draw[thick,red]  (1,3.2) -- (2,2.8) -- (3,2.6) -- (4,2.6) -- (5,2.58) -- (6,2.57) -- (7,2.57) -- (8,2.57);
	\node[blue] at (8.3,2.45) {$u_n$};
	\node[red]  at (8.3,2.7) {$v_n$};
	\draw[dashed] (0,2.57) -- (9,2.57) node[right] {$\theta$};
	\end{tikzpicture}
	\caption{The two monotone sequences $u_n=q_n^{3^{-n}}$ and $v_n=(q_n+1)^{3^{-n}}$ converge to the common limit $\theta$; the Mills recurrence \cref{eq:mills-seq-intro} ensures $u_n<\theta<v_n$ for every $n$.}
	\label{fig:mills-recurrence3}
\end{figure}

\begin{figure}[H]
	\centering
	\begin{tikzpicture}[scale=1.0]
	% x-axis
	\draw[->] (0,0) -- (11,0) node[right] {\(x\)};
	% interval for n=1
	\draw[thick, blue] (1,1) -- (4,1)
	node[midway, above, font=\scriptsize] {\((q_1^3,(q_1+1)^3-1)\)};
	% interval for n=2 (nested)
	\draw[thick, red] (2,2) -- (3.5,2)
	node[midway, above, font=\scriptsize] {\((q_2^3,(q_2+1)^3-1)\)};
	% interval for n=3 (nested)
	\draw[thick, orange] (2.5,3) -- (3.2,3)
	node[midway, above, font=\scriptsize] {\((q_3^3,(q_3+1)^3-1)\)};
	% limit θ
	\draw[dashed, thick, purple] (2.9,-0.5) -- (2.9,3.5);
	\node at (2.9,4) {\(\theta\)};
	% third-root projections
	\foreach \y in {1,2,3} {
		\draw[gray, dotted] (2.9,\y) -- (2.9,0);
	}
	\node[font=\scriptsize, align=center] at (5.5,1.7)
	{Nested intervals \((q_n^3,(q_n+1)^3-1)\)};
	\node[font=\scriptsize, align=center] at (5.5,0.7)
	{accumulate at \(\theta\), giving \(q_n=[\theta^{3^n}]\)};
	\end{tikzpicture}
	\caption{ The intervals
		\((q_n^3,(q_n+1)^3-1)\) are nested and shrink geometrically; their common
		point is \(\theta^{3^n}\) for the Mills constant \(\theta\). Taking
		\(3^n\)-th roots gives the definition of \(\theta\) as
		\(\lim q_n^{3^{-n}}\).}
	\label{fig:mills}
\end{figure}
\begin{figure}[H]
	\centering
	\begin{tikzpicture}[xscale=0.55,yscale=0.9]
	\draw[->] (0,0) -- (11,0) node[right] {$n$};
	\draw[->] (0,0) -- (0,5) node[above] {$\log q_n$};
	\draw[blue,thick] (0.5,0.5) -- (3,2.0) -- (6,3.6) -- (10,4.6);
	\fill (0.5,0.5) circle (1.5pt) node[below right, inner sep=2pt] {$q_1$};
	\fill (3,2.0)   circle (1.5pt) node[below right, inner sep=2pt] {$q_2$};
	\fill (6,3.6)   circle (1.5pt) node[below right, inner sep=2pt] {$q_3$};
	\fill (10,4.6)  circle (1.5pt) node[below left,  inner sep=2pt] {$q_4$};
	\draw[dashed] (0.5,0) -- (0.5,0.5);
	\draw[dashed] (3,0)   -- (3,2.0);
	\draw[dashed] (6,0)   -- (6,3.6);
	\draw[dashed] (10,0)  -- (10,4.6);
	\node[blue,anchor=west] at (11,3.5) {$\log q_n$};
	\end{tikzpicture}
	\caption{The Mills recurrence}
	\label{fig:mills-recurrence}
	
\end{figure}
\normalsize
In \autoref{fig:mills-recurrence} $\log q_{n}$ grows by a factor $3$ at each
step, so 
\begin{equation}
\log \theta =\lim 3^{-n}\log q_{n}.\newline
\end{equation}%
(The printed $5^{-n}$ in the historical transcription is a typo; see %
\Cref{rem:typo}.) \newline

% =========================================================================
\begin{thebibliography}{99}
	
	\bibitem{Mills1947} W.~H.~Mills. A prime-representing function. \emph{Bull.
		Amer. Math. Soc.} \textbf{53} (1947), 604.
	\url{https://doi.org/10.1090/S0002-9904-1947-08849-2}
	
	\bibitem{Kuipers1950} L.~Kuipers. Prime-representing functions. \emph{Indag.
		Math.} \textbf{12} (1950), 57--58.
	\url{https://zbmath.org/?q=an:0036.05102}
	
	\bibitem{Ansari1951} A.~R.~Ansari. On prime representing function. \emph{Ganita}
	\textbf{2} (1951), 81--82.
	\url{https://zbmath.org/?q=an:0043.04302}
	
	\bibitem{Wright1951} E.~M.~Wright. A prime-representing function. \emph{Amer.
		Math. Monthly} \textbf{58} (1951), 616--618.
	\url{https://doi.org/10.1080/00029890.1951.11999747}
	
	\bibitem{Ore1952} O.~Ore. On the selection of subsequences. \emph{Proc. Amer.
		Math. Soc.} \textbf{3} (1952), 706--712.
	\url{https://doi.org/10.1090/S0002-9939-1952-0050546-2}
	
	\bibitem{Niven1951} I.~Niven. Functions which represent prime numbers. \emph{Proc.
		Amer. Math. Soc.} \textbf{2} (1951), 753--755.
	\url{https://doi.org/10.1090/S0002-9939-1951-0043807-4}
	
	\bibitem{Wright1954} E.~M.~Wright. A class of representing functions. \emph{J.
		London Math. Soc.} \textbf{29} (1954), 63--71.
	\url{https://doi.org/10.1112/jlms/s1-29.1.63}
	
	\bibitem{Ingham1937} A.~E.~Ingham. On the difference between consecutive
	primes. \emph{Quart. J. Math. Oxford Ser.} (2) \textbf{8} (1937), 255--266.
	\url{https://doi.org/10.1093/qmath/os-8.1.255}
	
	\bibitem{Dudley1969} U.~Dudley. History of a formula for primes. \emph{Amer.
		Math. Monthly} \textbf{76} (1969), no.~1, 23--28.
	\url{https://doi.org/10.1080/00029890.1969.12000121}
	
	\bibitem{HardyWright1960} G.~H.~Hardy and E.~M.~Wright. \emph{An Introduction to
		the Theory of Numbers}, 4th ed. Oxford University Press, 1960. Theorem~418.
	\url{https://search.worldcat.org/search?q=An+Introduction+to+the+Theory+of+Numbers+Hardy+Wright}
	
	\bibitem{Simpson2009} S.~G.~Simpson. \emph{Subsystems of Second Order
		Arithmetic}, 2nd ed. Cambridge University Press, 2009.
	\url{https://doi.org/10.1017/CBO9780511581007}
	
	\bibitem{FormalConjectures} Google DeepMind. \emph{formal-conjectures}. Lean
	repository, 2026.
	\url{https://github.com/google-deepmind/formal-conjectures}
	
	\bibitem{MathlibGraph} X.~Li and P.~Shafto. The Multinetwork of
	Mathlib. NeurIPS 2026 Datasets and Benchmarks Track (under review).
	Dataset files used for the dependency depth: \texttt{v2/declaration/metrics.csv}
	(column \texttt{dag\_layer} = topological depth of each declaration in the
	declaration DAG) and \texttt{mathlib\_edges.csv} (raw dependency edges).
	{\scriptsize
		\url{https://huggingface.co/datasets/MathNetwork/MathlibGraph}.\newline
		\url{https://huggingface.co/datasets/MathNetwork/MathlibGraph/resolve/main/v2/declaration/metrics.csv}\newline
		\url{https://huggingface.co/datasets/MathNetwork/MathlibGraph/resolve/main/mathlib_edges.csv}.}
	
	\bibitem{Bang1952} T.~Bang. A function representing all prime numbers. \emph{Nordisk Mat.
		Tidskr.} \textbf{34} (1952), 117--118.
	\url{https://zbmath.org/?q=an:0048.03108}
	
	\bibitem{CaldwellCheng2005} C.~K.~Caldwell and Y.~Cheng. Determining Mills'
	constant and a note on Honaker's problem. \emph{J. Integer Seq.} \textbf{8}
	(2005), Article 05.4.1.
	\url{https://cs.uwaterloo.ca/journals/JIS/VOL8/Caldwell/caldwell78.pdf}
	
	\bibitem{Haneke1962} W.~Haneke. Versch\"arfung der Absch\"atzung von
	$\zeta(1/2+it)$. \emph{Acta Arith.} \textbf{8} (1962/1963), 357--430.
	\url{https://doi.org/10.4064/aa-8-4-357-430}
	
	\bibitem{Saito2025} K.~Saito. Mills' constant is irrational. \emph{Mathematika}
	\textbf{71} (2025), no.~3, e70027.
	\url{https://arxiv.org/abs/2404.19461}
	
	\bibitem{Sierpinski1952} W.~Sierpi\'nski. Sur une formule donnant tous les
	nombres premiers. \emph{C. R. Acad. Sci. Paris} \textbf{235} (1952), 1078--1079.
	\url{https://gallica.bnf.fr/ark:/12148/bpt6k31936/f1078}
	
	\bibitem{Srinivasan1961} B.~R.~Srinivasan. Formulae for the $n$th prime.
	\emph{J. Indian Math. Soc.} (N.S.) \textbf{25} (1961), 33--39; \textbf{26}
	(1962), 248.
	\url{https://zbmath.org/?q=an:0111.25002}
	
	\bibitem{Srinivasan1962} B.~R.~Srinivasan. An arithmetical function and an
	associated formula for the $n$th prime, I, II. \emph{Norske Vid. Selsk. Forh.
		Trondheim} \textbf{35} (1962), 68--71, 72--75.
	\url{https://zbmath.org/?q=an:0117.03003}
	
	\bibitem{Willans1964} C.~P.~Willans. On formulae for the $n$th prime number.
	\emph{Math. Gaz.} \textbf{48} (1964), 413--415.
	\url{https://doi.org/10.2307/3611815}
	
\end{thebibliography}

\end{document}